\documentclass[11pt,a4paper]{article}
\usepackage{amsmath,amsxtra,amsthm,amssymb,mathrsfs}
\usepackage{hyperref}
\usepackage[margin=2cm]{geometry}
\usepackage{graphicx}
\usepackage{xcolor}
\usepackage[normalem]{ulem}
\usepackage[shortlabels]{enumitem}
%\usepackage[small,nohug,noPostScript]{diagrams}
%\usepackage{stmaryrd}
\usepackage[all]{xy}
\usepackage{float}
\usepackage{tocloft}
\usepackage{tikz}
\usepackage{mdframed}
\usepackage{comment}
%\usetikzlibrary{positioning}
\usetikzlibrary{shapes,arrows}
\usetikzlibrary{patterns}
\usetikzlibrary{decorations.shapes}
\usetikzlibrary{decorations.pathreplacing}
\usetikzlibrary{decorations.pathmorphing}
\numberwithin{equation}{section}

\setlength{\parindent}{0mm}
\setlength{\parskip}{11pt}

%set default enumerate environment
\setlist[enumerate]{label={\rm(\arabic*)},itemsep=3pt,topsep=-6pt}
\setlist[itemize]{itemsep=3pt,topsep=-6pt}

\newenvironment{claim}[1][]
{\begin{list}{}{\setlength{\itemsep}{0pt}
	\setlength{\listparindent}{0pt}
	\setlength{\labelsep}{0.2cm}
	\setlength{\topsep}{2pt}
	\setlength{\leftmargin}{0.5cm}}

	\item[] \textbf{Claim:$\mbox{}\quad$}#1
	\item[] \textbf{Proof of Claim: }\\
}
{\end{list}}

%%% set up proof environment
\makeatletter
\renewenvironment{proof}[1][\proofname]{\par
    \pushQED{\qed}%
    \normalfont \topsep0\p@\@plus6\p@ \labelsep1em\relax
    \trivlist
    \item[\hskip\labelsep\indent
        \bfseries #1]\ignorespaces
    \mbox{}
}{
    \popQED\endtrivlist\@endpefalse
}
\makeatother

%%% Remove section numbering but retain section counting
%\makeatletter
%\def\@seccntformat#1{%
%  \expandafter\ifx\csname c@#1\endcsname\c@section\else
%  \csname the#1\endcsname\quad
%  \fi}
%\makeatother

%add dots to sections in TOC
\renewcommand{\cftsecleader}{\cftdotfill{\cftdotsep}}
%show section numbers in TOC
\renewcommand{\cftsecpresnum}{}
\renewcommand{\cftsecaftersnum}{}

%abovespace,belowspace,bodyfont,indent,headfont,headpunct(e.g.\newline),headspace,custom-head-spec
\newtheoremstyle{all}{10pt}{0pt}{\itshape}{0pt}
{\bfseries}{.}{5pt plus 1pt minus 1pt}{\thmname{#1}\thmnumber{ #2}\thmnote{ (#3)}}
\theoremstyle{all}

\newtheorem{theorem}{Theorem}[section]
\newtheorem{proposition}[theorem]{Proposition}
\newtheorem{lemma}[theorem]{Lemma}
\newtheorem{corollary}[theorem]{Corollary}
\newtheorem{example}[theorem]{Example}
\newtheorem*{example*}{Example}
\newtheorem{definition}[theorem]{Definition}
\newtheorem{exercise}[theorem]{Exercise}
\newtheorem*{exercise*}{Exercise}
\newtheorem{assumption}[theorem]{Assumption}
\newtheorem*{convention}{Convention}
\newtheorem*{notation}{Notation}

\theoremstyle{remark}
\newtheorem{remark}[theorem]{Remark}
\newtheorem*{remark*}{Remark}


\newcommand{\comm}[1]{{\color[rgb]{0.1,0.9,0.1}#1}}
\newcommand{\old}[1]{{\color{red}#1}}
\newcommand{\new}[1]{{\color{blue}#1}}
\definecolor{darkblue}{rgb}{0,0,0.7}
\newcommand{\defn}[1]{\textsl{\color{darkblue} #1}}

\definecolor{supplementarygreen}{rgb}{0,0.35,0.20}
\definecolor{supplementarybackground}{HTML}{E8F5E9}
\definecolor{advancedbackground}{HTML}{EDE8D0}
\definecolor{advancedred}{HTML}{B22222}
\newmdenv[
  leftline=true,
  rightline=false,
  topline=false,
  bottomline=false,
  linecolor=advancedred,
  linewidth=1.5pt,
  backgroundcolor=advancedbackground,
  suppressfirstparskip=false,
  leftmargin=1em,
  rightmargin=0pt,
  innerleftmargin=0.8em,
  innerrightmargin=0.8em,
  innertopmargin=10pt,
  innerbottommargin=10pt,
  splittopskip=\dimexpr\topskip+10pt\relax,
  splitbottomskip=10pt,
  skipabove=8pt,
  skipbelow=8pt
]{advanced}
\newmdenv[
  leftline=true,
  rightline=false,
  topline=false,
  bottomline=false,
  linecolor=supplementarygreen,
  linewidth=1.5pt,
  backgroundcolor=supplementarybackground,
  suppressfirstparskip=false,
  leftmargin=1em,
  rightmargin=0pt,
  innerleftmargin=0.8em,
  innerrightmargin=0.8em,
  innertopmargin=10pt,
  innerbottommargin=10pt,
  splittopskip=\dimexpr\topskip+10pt\relax,
  splitbottomskip=10pt,
  skipabove=8pt,
  skipbelow=8pt
]{supplementary}

\newif\ifshowhidden
\showhiddenfalse % Change to \showhiddentrue to include hidden material.
% Put \begin{hidden} and \end{hidden} on separate, unindented lines.
\ifshowhidden
  \includecomment{hidden}
\else
  \excludecomment{hidden}
\fi

\DeclareMathOperator{\isom}{\cong}
\DeclareMathOperator{\ifonlyif}{\Leftrightarrow}
\DeclareMathOperator{\imply}{\Rightarrow}
\DeclareMathOperator{\onto}{\twoheadrightarrow}
\DeclareMathOperator{\into}{\hookrightarrow}
\newcommand{\ses}[3]{0\rightarrow #1 \rightarrow #2 \rightarrow #3 \rightarrow 0}

\renewcommand{\P}{\mathbb{P}}
\newcommand{\Z}{\mathbb{Z}}
\newcommand{\R}{\mathbb{R}}
\newcommand{\C}{\mathbb{C}}
\newcommand{\bbA}{\mathbb{A}}
\newcommand{\CB}{\mathcal{B}}
\newcommand{\CC}{\mathcal{C}}
\newcommand{\CD}{\mathcal{D}}
\newcommand{\rmC}{\mathrm{C}}
\newcommand{\calX}{\mathcal{X}}

\newcommand{\id}{\operatorname{id}}
\newcommand{\soc}{\operatorname{soc}\nolimits}
\newcommand{\rad}{\operatorname{rad}\nolimits}
\renewcommand{\top}{\operatorname{top}\nolimits}
\newcommand{\Ext}{\operatorname{Ext}\nolimits}
\newcommand{\Tor}{\operatorname{Tor}\nolimits}
\newcommand{\Hom}{\operatorname{Hom}\nolimits}
\newcommand{\End}{\operatorname{End}\nolimits}
\newcommand{\Aut}{\operatorname{Aut}\nolimits}
\newcommand{\op}{\mathrm{op}}
\newcommand{\en}{\mathrm{en}}
\newcommand{\grstar}[1][]{(*#1)}
\renewcommand{\Im}{\operatorname{Im}\nolimits}
\newcommand{\Ker}{\operatorname{Ker}\nolimits}
\newcommand{\Cok}{\operatorname{Cok}\nolimits}

\DeclareMathOperator{\moduleCategory}{\mathrm{mod}} \renewcommand{\mod}{\moduleCategory}
\DeclareMathOperator{\Mod}{\mathrm{Mod}}
\DeclareMathOperator{\proj}{\mathrm{proj}}
\DeclareMathOperator{\inj}{\mathrm{inj}}
\DeclareMathOperator{\ind}{\mathrm{ind}}
\DeclareMathOperator{\Grmod}{\mathrm{GrMod}}
\DeclareMathOperator{\grmod}{\mathrm{grmod}}

\DeclareMathOperator{\Char}{char}
\newcommand{\triv}{\operatorname{triv}\nolimits}
\newcommand{\Mat}{\operatorname{Mat}\nolimits}
\newcommand{\GL}{\operatorname{GL}\nolimits}
\DeclareMathOperator{\Sn}{\mathfrak{S}}
\newcommand{\Res}{\operatorname{Res}\nolimits}
\newcommand{\Ind}{\operatorname{Ind}\nolimits}
\newcommand{\Symm}{\operatorname{Sym}\nolimits}
\newcommand{\SymAlg}{\mathsf{S}}
\newcommand{\Stab}{\operatorname{Stab}\nolimits}
\newcommand{\Tr}{\operatorname{Tr}\nolimits}

\newcommand{\Ann}{\operatorname{Ann}\nolimits}
\newcommand{\pdim}{\operatorname{pdim}}
\newcommand{\idim}{\operatorname{idim}}
\newcommand{\gldim}{\operatorname{gldim}\nolimits}
\newcommand{\Rej}{\operatorname{Rej}\nolimits}
\newcommand{\std}{\Delta}
\newcommand{\costd}{\nabla}
\newcommand{\filt}{\mathcal{F}}
\newcommand{\add}{\operatorname{add}}
\newcommand{\dimvec}{\mathbf{dim}}
\newcommand{\nunlhd}{\ntrianglelefteq}
\newcommand{\nlhd}{\ntriangleleft}
\newcommand{\nunrhd}{\ntrianglerighteq}
\newcommand{\nrhd}{\ntriangleright}
\newcommand{\LL}{\mathrm{LL}}
\newcommand{\Ab}{\operatorname{\mathcal{A}b}}
\newcommand{\bbD}{\mathbb{D}}


\begin{document}
\begin{center}
\textsc{{\LARGE Topics in Mathematics Science VII}\\
Spring 2026}\\
\textsc{\Large Introduction to Koszul Duality}

\textsc{Aaron Chan}
\end{center}

Last update: 6 October 2026

\setcounter{tocdepth}{1}
\tableofcontents

\section*{Convention}

Throughout the course, $\Bbbk$ is a field.
All rings are unital and associative.
All modules are right modules unless otherwise specified.
When we discuss dimensions, we mean dimensions as $\Bbbk$-vector spaces.

There are a lot of contents not covered in lecture:
\begin{supplementary}
This is supplementary material; often for `unnecessarily long pedantic check of definitions' or content of similar nature -- things that are only there for rigour
\end{supplementary}
\begin{advanced}
This is advanced material; things that I personally would like to add, from personal experience in the theory, but definitely not necessary for lecture purpose.
\end{advanced}
Neither of these are rigourously checked since they are mostly produced by AI, so the reader should read these contents with extra cares themselves.  All other contents are drafted by AI and QA'ed (heavily edited) by myself.


% !TEX root = Lecture_Koszul.tex
\section{Graded algebras, modules, and duals}

\subsection{Algebras and modules}

\begin{definition}
A \defn{$\Bbbk$-algebra} is a $\Bbbk$-vector space with an associative unital multiplication that is \defn{bilinear}, meaning linear in each variable.
\end{definition}

A right $A$-module is understood to be a $\Bbbk$-vector space $M$ with a bilinear action satisfying $(ma)b=m(ab)$ and $m1_A=m$ for $m\in M$ and $a,b\in A$.
A module homomorphism $f:M\to N$ is linear and satisfies $f(ma)=f(m)a$.
A homomorphism of $\Bbbk$-algebras is a linear map preserving multiplication and the identity.

\begin{example}\label{ex:triangular-algebra}
(1) The matrix ring $\Mat_n(\Bbbk)$ is a $\Bbbk$-algebra.
The space of $1\times n$ row vectors is a right $\Mat_n(\Bbbk)$-module by matrix multiplication.

(2) Let $\Bbbk A_n$ be the \defn{upper triangular matrix ring}, which is also a $\Bbbk$-algebra.
\[
\Bbbk A_n = \begin{pmatrix}
\Bbbk & \Bbbk & \cdots & \Bbbk \\
0 & \Bbbk & \cdots & \Bbbk \\
0 & 0 & \ddots & \vdots \\
0 & \cdots & 0 & \Bbbk
\end{pmatrix}
\]
For those familiar with path algebras, $\Bbbk A_n$ is the path algebra of a \defn{linearly oriented $A_n$ quiver}, which explains the notation. We use matrix calculations, so no knowledge of quivers is required.
Write $e_{ij}$ for the matrix with $1$ in entry $(i,j)$ and zero elsewhere.
Our convention composes paths from left to right: $e_{ij}$ represents $i\to j$, and $e_{ij}e_{j\ell}=e_{i\ell}$.
\end{example}

\subsection{Categories and functors}\label{subsec:category-language}

We organise modules together with their homomorphisms. A category records both: specifying the objects alone does not specify a category.

\begin{definition}\label{def:category}
A \defn{category} $\CC$ consists of objects, a set $\Hom_{\CC}(X,Y)$ of \defn{morphisms} $f:X\to Y$ for each pair of objects, and composition sending $f:X\to Y$ and $g:Y\to Z$ to $g\circ f:X\to Z$.
Each object $X$ has an identity $\id_X:X\to X$, and
\[
h\circ(g\circ f)=(h\circ g)\circ f,
\qquad \id_Y\circ f=f=f\circ\id_X
\quad\text{for all composable }f,g,h.
\]
\end{definition}

A morphism $f:X\to Y$ is an \defn{isomorphism} if some $g:Y\to X$ satisfies $g\circ f=\id_X$ and $f\circ g=\id_Y$; we then write $X\cong Y$.
A \defn{full subcategory} of $\CC$ consists of chosen objects of $\CC$, with all morphisms between them and the same composition.

\begin{example}\label{ex:module-categories}
For a $\Bbbk$-algebra $A$, we use the notation
\[
\begin{aligned}
\Mod A&:=\text{category of right $A$-modules and $A$-linear maps},\\
\mod A&:=\text{full subcategory of finitely generated right $A$-modules}.
\end{aligned}
\]
In either category, $\Hom_A(M,N)$ is the space of $A$-linear maps $M\to N$; composition is the usual composition of maps.
For $A=\Bbbk$, these are the categories of all and of finite-dimensional vector spaces, respectively.

Note that, if $\dim_\Bbbk A<\infty$, then $M\in \mod A \;\Leftrightarrow\; M$ is finite-dimensional ($\because$ finite generation implies that $M$ is a quotient of $A^{\oplus n}$ some $n$, and a finite $\Bbbk$-basis of $M$ is clearly a finite generating set.)
\end{example}

\begin{example}\label{ex:add-category}
For $M\in\mod A$, let $\add M:=\add_{\mod A}(M)$ be the full subcategory with
\[
N\in\add M
\quad\Longleftrightarrow\quad
N\oplus N'\cong M^{\oplus r}
\quad\text{for some }N'\in\mod A\text{ and integer }r\geq0.
\]
Thus $\add M$ consists of direct summands of finite direct sums of $M$.
\end{example}

In modern representation theory, remembering $A$ and $\add A$ holds the same categorical information is crucial in utilising cateogircal technique in representation theory; this allows one to push many results to a more abstract framework.

\begin{definition}\label{def:functor}
A \defn{functor} $F:\CC\to\CD$ assigns an object $F(X)$ to each object $X$ and a morphism $F(f):F(X)\to F(Y)$ to each $f:X\to Y$, such that
\[
F(\id_X)=\id_{F(X)},\qquad F(g\circ f)=F(g)\circ F(f)
\quad\text{for all composable }f,g.
\]
\end{definition}

For example, forgetting the $A$-action gives a functor $\Mod A\to\Mod\Bbbk$: it keeps the underlying vector spaces and linear maps.
The \defn{identity functor} $\id_{\CC}$ leaves every object and morphism unchanged; functors compose by applying one after the other.
Every functor preserves isomorphisms, since $F(f^{-1})$ is inverse to $F(f)$.

\phantomsection\label{not:opposite-category}
The \defn{opposite category} $\CC^\op$ has the same objects as $\CC$, with $\Hom_{\CC^\op}(X,Y)=\Hom_{\CC}(Y,X)$ and composition reversed.
A \defn{contravariant functor} from $\CC$ to $\CD$ means a functor $\CC^\op\to\CD$.
Thus it sends $f:X\to Y$ in $\CC$ to $F(f):F(Y)\to F(X)$ and satisfies $F(g\circ f)=F(f)\circ F(g)$.
Duals will provide our main examples.

\begin{definition}\label{def:natural-transformation}
For functors $F,G:\CC\to\CD$, a \defn{natural transformation} $\eta:F\to G$ is a family of morphisms $\eta_X:F(X)\to G(X)$ such that
\begin{equation}\label{eq:naturality-square}
\xymatrix@C=45pt@R=20pt{
F(X) \ar[r]^{\eta_X} \ar[d]_{F(f)} & G(X) \ar[d]^{G(f)}\\
F(Y) \ar[r]_{\eta_Y} & G(Y)
}
\qquad
G(f)\circ\eta_X=\eta_Y\circ F(f)
\quad\text{for every }f:X\to Y.
\end{equation}
If every $\eta_X$ is an isomorphism, $\eta$ is a \defn{natural isomorphism}, written $F\cong G$.
\end{definition}

The square in \eqref{eq:naturality-square} is said to \defn{commute}: its two routes from $F(X)$ to $G(Y)$ give the same morphism.
We often say a family of morphisms (resp. isomorphisms) in a category $\CC$ is \defn{natural in $M\in\CC$} to mean that it defines a natural transformation (resp. isomorphism).

For example, for fixed $c\in\Bbbk$, the maps $c\id_M:M\to M$ define a natural transformation $\id_{\Mod A}\to\id_{\Mod A}$ because $f(cm)=cf(m)$ for every module homomorphism $f:M\to N$.
It is a natural isomorphism when $c\neq0$.

A functor $F:\CC\to\CD$ is an \defn{equivalence} if there is a functor $G:\CD\to\CC$ with natural isomorphisms $GF\cong\id_{\CC}$ and $FG\cong\id_{\CD}$.
A \defn{duality} between $\CC$ and $\CD$ is an equivalence $\CC^\op\to\CD$.
Thus an equivalence recovers objects and maps up to natural isomorphism; a duality also reverses the direction of maps.

\begin{definition}\label{def:exact-sequence}
Let $A$ be a $\Bbbk$-algebra, and let $L,M,N\in\Mod A$, with $i:L\to M$ and $p:M\to N$ module homomorphisms.
The sequence $L\xrightarrow{i}M\xrightarrow{p}N$ is \defn{exact at $M$} if $\Im i=\Ker p$.  
A sequence of module homomorphisms is \defn{exact} if it is exact at every term with an incoming and an outgoing map.
A \defn{short exact sequence} is an exact sequence
\begin{equation}\label{eq:short-exact-sequence}
0\longrightarrow L\xrightarrow{i}M\xrightarrow{p}N\longrightarrow0,
\qquad
\Ker i=0,\quad\Im i=\Ker p,\quad\Im p=N.
\end{equation}
\end{definition}


\begin{definition}\label{def:exact-functor}\label{def:exact-contravariant}
Let $A,B$ be $\Bbbk$-algebras.
A functor $F:\Mod A\to\Mod B$ is \defn{exact} if $F(0)=0$ and every short exact sequence \eqref{eq:short-exact-sequence} gives a short exact sequence
\[
0\longrightarrow F(L)\xrightarrow{F(i)}F(M)\xrightarrow{F(p)}F(N)\longrightarrow0.
\]
\end{definition}

An exact functor preserves injective and surjective maps: apply it to
$0\to L\xrightarrow{i}M\to M/i(L)\to0$ and $0\to\Ker p\to M\xrightarrow{p}N\to0$, respectively.
For a functor $F:(\Mod A)^\op\to\Mod B$, exactness instead means that \eqref{eq:short-exact-sequence} gives
\[
0\longrightarrow F(N)\xrightarrow{F(p)}F(M)\xrightarrow{F(i)}F(L)\longrightarrow0.
\]
Thus an exact contravariant functor exchanges injective and surjective maps.

\begin{example}[$\Bbbk$-linear duality]\label{ex:linear-duality}
For a vector space $V$, define its \defn{$\Bbbk$-linear dual} by $D(V):=\Hom_\Bbbk(V,\Bbbk)$.
For a linear map $u:V\to W$, set
\[
D(u):D(W)\longrightarrow D(V),\qquad \lambda\longmapsto\lambda\circ u.
\]
The identities $D(\id_V)=\id_{D(V)}$ and $D(v\circ u)=D(u)\circ D(v)$ make $D:(\Mod\Bbbk)^\op\to\Mod\Bbbk$ a functor.
The functor $D$ is exact.
\begin{supplementary}[nobreak=true]
For a short exact sequence \eqref{eq:short-exact-sequence} of vector spaces, $D(p)$ is injective because $p$ is surjective, and $D(i)$ is surjective because a functional on a subspace extends to the whole vector space.
If $\lambda\in D(M)$ satisfies $\lambda\circ i=0$, then
\[
\mu:N\longrightarrow\Bbbk,\qquad p(m)\longmapsto\lambda(m)
\]
is well-defined since $\Ker p=i(L)$, and $\lambda=\mu\circ p=D(p)(\mu)$.
Conversely, $p\circ i=0$ gives $D(i)\circ D(p)=0$, so $\Ker D(i)=\Im D(p)$.
Thus $D$ exchanges injective and surjective linear maps.
\end{supplementary}

For every vector space $V$, evaluation defines a natural injective map
\begin{equation}\label{eq:vector-bidual-evaluation}
\mathrm{ev}_V:V\longrightarrow D(D(V)),\qquad
v\longmapsto(\lambda\mapsto\lambda(v)).
\end{equation}
If $V$ is finite-dimensional, choose a basis $v_1,\ldots,v_n$ and its \defn{dual basis} $v_1^*,\ldots,v_n^*$, defined by $v_i^*(v_j)=\delta_{ij}$.
Then
\[
\mathrm{ev}_V^{-1}:D(D(V))\longrightarrow V,\qquad
\Phi\longmapsto\sum_{i=1}^n\Phi(v_i^*)v_i.
\]
Consequently, $D:(\mod\Bbbk)^\op\to\mod\Bbbk$ is a duality.
\begin{supplementary}
If $v\neq0$, choose $\lambda\in D(V)$ with $\lambda(v)\neq0$; hence $\mathrm{ev}_V$ is injective.
For $u:V\to W$, $v\in V$, and $\lambda\in D(W)$,
\[
D(D(u))(\mathrm{ev}_V(v))(\lambda)=\lambda(u(v))
=\mathrm{ev}_W(u(v))(\lambda),
\]
which proves naturality.
\begin{sloppypar}
In finite dimension, writing $\lambda=\sum_{i=1}^n\lambda(v_i)v_i^*$ gives
$\Phi(\lambda)=\lambda(\sum_{i=1}^n\Phi(v_i^*)v_i)$.
Thus the displayed formula is the inverse of \eqref{eq:vector-bidual-evaluation}, and is independent of the chosen basis.
\end{sloppypar}
\end{supplementary}
\end{example}

\subsection{Semisimple algebras}

\begin{definition}\label{def:semisimple}
A module $L$ is \defn{simple} if $L\neq 0$ and its only submodules are $0$ and $L$.
A module is \defn{semisimple} if it is a direct sum of simple modules.
An algebra $B$ is \defn{semisimple} if it is semisimple as a right module $B_B$ over itself.
The \defn{Jacobson radical} $\rad B$ is the intersection of all maximal right ideals of $B$.
For a right $B$-module $M$, the \defn{radical} $\rad M$ is the intersection of all maximal submodules of $M$, with $\rad M=M$ if there are none.
\end{definition}

We use the following standard ring-theoretic facts without proof.
For a finite-dimensional $\Bbbk$-algebra $B$ and any right $B$-module $M$,
\[
\rad M=M\rad B.
\]
Recall that a \defn{division $\Bbbk$-algebra} is a non-zero $\Bbbk$-algebra whose non-zero elements are invertible.

\begin{theorem}[Artin--Wedderburn]\label{thm:artin-wedderburn}
For a non-zero finite-dimensional $\Bbbk$-algebra $B$,
\[
B\text{ is semisimple}
\quad\Longleftrightarrow\quad
\rad B=0
\quad\Longleftrightarrow\quad
B\cong\prod_{i=1}^n\Mat_{n_i}(D_i),
\]
where $n,n_i\geq 1$ and each $D_i$ is a finite-dimensional division $\Bbbk$-algebra.
If $\Bbbk$ is algebraically closed, then $D_i=\Bbbk$ for every $i$.
For semisimple $B$, every left or right $B$-module is semisimple, and every submodule has a complement. In particular, the definition is independent of the chosen side.
\end{theorem}

\begin{advanced}
For proofs of $\rad M=M\rad B$ and of Theorem~\ref{thm:artin-wedderburn}, see for example \cite{DK94}; \cite{Bez23} gives a short account of semisimple modules and Wedderburn's theorem.
\end{advanced}

\begin{corollary}\label{cor:semisimple-summands}
Every finite-dimensional left or right module over a finite-diemnsional semisimple algebra $S$ is a direct summand of $S^{\oplus n}$ for some $n$, i.e. $M\in \add(S)$ for all $M\in \mod S$.
\end{corollary}

\begin{proof}
Choose a surjection $p:S^{\oplus n}\onto M$.
By Theorem~\ref{thm:artin-wedderburn}, there is a submodule $L$ such that $S^{\oplus n}=\Ker p\oplus L$.
The restriction $p|_L:L\to M$ has kernel $L\cap\Ker p=0$ and image $p(L)=p(S^{\oplus n})=M$.
Hence $L\cong M$, so $S^{\oplus n}\cong\Ker p\oplus M$.
\end{proof}

\begin{example}\label{ex:semisimple-algebras}
The algebras $\Bbbk^{\oplus n}$ and $\Mat_n(\Bbbk)$ are semisimple by Theorem~\ref{thm:artin-wedderburn}.
Upper triangular matrices need not form a semisimple algebra: for $A=\Bbbk A_2$,
\[
\rad A=\Bbbk e_{12}\neq0,\qquad
A/\rad A\cong\Bbbk^{\oplus2}.
\]
\begin{supplementary}
Put $J=\Bbbk e_{12}$, so $J^2=0$.
For a simple right module $L$, simplicity gives $LJ=0$ or $LJ=L$; the second alternative implies $L=LJ^2=0$, a contradiction.
Hence $J$ is contained in every maximal right ideal.
The quotient $A/J\cong\Bbbk^{\oplus2}$ has exactly two maximal right ideals, whose inverse images have intersection
\[
(\Bbbk e_{11}\oplus\Bbbk e_{12})\cap
(\Bbbk e_{12}\oplus\Bbbk e_{22})=J.
\]
\end{supplementary}
\end{example}

\begin{definition}\label{def:idempotents}
An element $e\in A$ is \defn{idempotent} if $e^2=e$.
Idempotents $e,f$ are \defn{orthogonal} if $ef=fe=0$.
A non-zero idempotent is \defn{primitive} if it is not the sum of two non-zero orthogonal idempotents.
A finite family $e_1,\ldots,e_r$ of pairwise orthogonal idempotents is \defn{complete} if $\sum_{i=1}^r e_i=1_A$.
\end{definition}

\begin{remark}\label{rem:idempotent-units}
For orthogonal idempotents $e,f\in A$,
\[
(e+f)A=eA\oplus fA.
\]
In particular, we always have $A\cong eA\oplus (1-e)A$.  
This strengthen the viewpoint of idempotent-as-local-identity: left-multiplying $e$ defines a canonical projection onto the direct summand $eA$ of $A$, and yields a commutative diagram
\[
\xymatrix@C=42pt@R=18pt{
A \ar[rr]^{a\mapsto ea} \ar[dr]_{\text{projection }a\mapsto ea\qquad} && A\\
& eA \ar[ur]_{\text{inclusion}} &
}
\]
So multiplying $e$ means we work inside the direct summand $eA$.
The \defn{corner subalgebra} $eAe=\{eae\mid a\in A\}$ has identity $e$.\footnote{note that this is not a unital subalgebra since the identity element changes; a more proper way to think about this is as a subcategory $\add_{\mod A}(eA)\simeq \add_{\mod eAe}(eAe)$.}
For $A=\prod_{i=1}^r A_i$, the \defn{coordinate idempotents} $e_i=(0,\ldots,1_{A_i},\ldots,0)$ are a complete orthogonal family and are the identities of the individual factors.
In $A=\Bbbk A_n$, the diagonal matrix units $e_{11},\ldots,e_{nn}$ are primitive, since $e_{ii}Ae_{ii}=\Bbbk e_{ii}$.
For $n\geq2$, the sum $e=e_{11}+\cdots+e_{n-1,n-1}$ gives the corner
\[
eAe=\left\{\begin{pmatrix}B&0\\0&0\end{pmatrix}\;\middle|\;B\in\Bbbk A_{n-1}\right\}
\cong\Bbbk A_{n-1}.
\]
\end{remark}

Idempotents also describe maps out of the summands of $A_A$.
\begin{lemma}[Yoneda lemma]\label{lem:idempotent-hom}
For an idempotent $e\in A$ and a right $A$-module $M$, evaluation at $e$ gives a natural (in $M$) vector space isomorphism
\[
\Hom_A(eA,M)\xrightarrow{\sim}Me,\qquad f\longmapsto f(e),
\quad\text{with inverse}\quad m\longmapsto(ea\mapsto ma).
\]
\end{lemma}
\begin{proof}
We have $f(e)=f(e)e$ and $f(ea)=f(e)a$.
If $m=me$, the map $ea\mapsto ma$ is well-defined because $ea=ea'$ implies $ma=mea=mea'=ma'$.
These formulas prove that the maps are inverse; postcomposition by a module map $M\to N$ proves naturality.
\end{proof}




\subsection{Graded algebras}

\begin{definition}\label{def:group-grading}
Let $G$ be a group with identity $e$.
A \defn{$G$-graded $\Bbbk$-algebra} is a $\Bbbk$-algebra $A$ with $1_A\in A_e$ and a decomposition
\[
A=\bigoplus_{g\in G}A_g
\quad\text{such that}\quad
A_gA_h\subset A_{gh}
\quad\text{for all }g,h\in G.
\]
An element of $A_g$ is \defn{homogeneous of degree $g$}.
\end{definition}

\begin{convention}
Unless stated otherwise, a graded algebra is $\Z$-graded and \uline{concentrated in non-negative degrees}.
\end{convention}

Note that the decomposition is into a direct sum, so every element of a graded algebra $A$ can always be written as a \emph{finite} sum of homogeneous elements.

\begin{example}\label{ex:triangular-grading}
For the upper triangular matrix ring $\Bbbk A_n$ of Example~\ref{ex:triangular-algebra}, the decomposition
\[
(\Bbbk A_n)_d=\operatorname{span}_\Bbbk\{e_{ij}\mid j-i=d\}
\]
makes $\Bbbk A_n$ a graded algebra.
\end{example}

\begin{supplementary}
\begin{example}
The algebra $\Bbbk A_3$ also has a $\Z^{\oplus 2}$-grading with
\[
\deg e_{ii}=(0,0),\qquad
\deg e_{12}=(1,0),\qquad
\deg e_{23}=(0,1),\qquad
\deg e_{13}=(1,1).
\]
The homomorphism $\Z^{\oplus 2}\to\Z$, $(p,q)\mapsto p+q$, gives the grading of Example~\ref{ex:triangular-grading}.
Alternatively, setting $\deg e_{ij}=2(j-i)$ gives a $\Z$-grading with $A_1=0$ but $A_2\neq 0$.
\end{example}
\end{supplementary}


\begin{definition}
A \defn{homomorphism of graded algebras} $f:A\to B$ is a $\Bbbk$-algebra homomorphism with $f(A_i)\subset B_i$ for all $i\in\Z$.
\end{definition}

\begin{definition}
A graded algebra $A$ is \defn{locally finite-dimensional} (\defn{lfd}) if $\dim_\Bbbk A_i<\infty$ for every $i\in\Z$.
\end{definition}
%
%\begin{assumption}\label{ass:degree-zero}
%For quadratic and Koszul algebras, $S:=A_0$ is finite-dimensional semisimple.
%\end{assumption}

\subsection{Graded modules}

\begin{definition}\label{def:graded-module}
Let $A=\bigoplus_{g\in G}A_g$ be a $G$-graded algebra.
A \defn{$G$-graded right $A$-module} is a right $A$-module $M$ with a direct sum decomposition
\[
M=\bigoplus_{g\in G}M_g
\quad\text{such that}\quad
M_gA_h\subset M_{gh}
\quad\text{for all }g,h\in G.
\]
An element of $M_g$ is called \defn{homogeneous of degree $g$}.

A $G$-graded $A$-module $M=\bigoplus_{g\in G} M_g$ is \defn{locally finite-dimensional} (\defn{lfd}) if $\dim_{\Bbbk} M_g <\infty$ for all $g\in G$.
\end{definition}

Unless specified otherwise, graded modules are $\Z$-graded; their degrees need not be non-negative.

\phantomsection\label{not:opposite-algebra}
For an algebra $A$, the \defn{opposite algebra} $A^\op$ has the same vector space as $A$, with multiplication $a^\op b^\op=(ba)^\op$.
A left $A$-module is a right $A^\op$-module by $m\cdot a^\op=am$.
Thus $\Mod(A^\op)$ is the category of left $A$-modules and left $A$-linear maps; not to be confused with the opposite category $(\Mod A)^\op$ though.
For a $\Z$-graded algebra $A=\bigoplus_{i\ge 0}A_i$, 
\[
\text{the opposite algebra is also graded with $(A^\op)_i=A_i$.}
\]



\begin{definition}\label{def:shift}
Let $M$ be a graded $A$-module and let $d\in\Z$.
The \defn{degree shift}, or \defn{grading shift}, $M(d)$ is the graded $A$-module whose underlying $A$-module is $M$, but whose grading is
\[
M(d)_j=M_{j+d}.
\]
\end{definition}

By Definition~\ref{def:shift}, $1\in A_0=A(-d)_d$, so $A(-d)$ has a generator in degree $d$.
If $M$ is concentrated in degree $0$, then $M(d)$ is concentrated in degree $-d$.

\begin{example}\label{ex:shifted-row}
For $A=\Bbbk A_3$ with the grading of Example~\ref{ex:triangular-grading}, the module $P_1:=e_{11}A$ consists of matrices supported on the first row.
Its basis $e_{11},e_{12},e_{13}$ has degrees $0,1,2$, respectively; in $P_1(-2)$, the same basis has degrees $2,3,4$.
\end{example}

\begin{definition}
Let $M,N$ be graded $A$-modules.
A homomorphism $f:M\to N$ is \defn{homogeneous of degree $d$} if
$f(M_j)\subset N_{j+d}$ for all $j\in\Z$. A \defn{homomorphism of graded modules} means a homogeneous $A$-module homomorphism of degree $0$, also said to \defn{preserve degrees}.
\end{definition}

\begin{example}\label{ex:matrix-hom}
For $A=\Bbbk A_3$, left multiplication by $e_{12}$ yields a map
\[
\begin{aligned}
e_{22}A&\longrightarrow e_{11}A,\\[-2pt]
\begin{pmatrix}0&0&0\\0&u&v\\0&0&0\end{pmatrix}
&\longmapsto
\begin{pmatrix}0&u&v\\0&0&0\\0&0&0\end{pmatrix}
\qquad(u,v\in\Bbbk).
\end{aligned}
\]
The $u$ and $v$ entries have degrees $0,1$ in the source and $1,2$ in the target, and so the map has degree $1$.
\end{example}

\begin{notation}\label{not:graded-hom}
For graded right $A$-modules $M,N$, the space of homomorphisms preserving degrees is
\begin{equation}\label{eq:graded-hom}
\Hom_A^{\Z}(M,N)=\{f\in\Hom_A(M,N)\mid f(M_j)\subset N_j\text{ for all }j\in\Z\}.
\end{equation}
Thus $\Hom_A^{\Z}(M,N(d))$ consists of homomorphisms $M\to N$ of degree $d$, since $N(d)_j=N_{j+d}$.
Note that, like the grading shift functor, different literature has different notation for this (there does not seem to be any widely agreed notation).
\end{notation}

\phantomsection\label{not:graded-module-categories}
We write 
\begin{align*}
\Grmod A &:= \text{category of graded right $A$-modules with morphism spaces $\Hom_A^{\Z}(M,N)$},\\
\grmod A &:= \text{full subcategory of $\Grmod A$ consisting of locally finite-dimensional modules}.
\end{align*}
Thus $\Grmod(A^\op)$ is the category of graded left $A$-modules and degree-preserving maps.
Forgetting the grading gives a functor $\Grmod A\to\Mod A$.

\begin{remark}
$\grmod A$ is somewhat analogous to $\mod A$ but not really the same:
If $A$ is locally finite-dimensional, every finitely generated graded $A$-module $M$ is locally finite-dimensional. ($\because$ a finite homogeneous generating set of a graded $A$-module $M$ gives a surjection from a finite direct sum of shifts of $A$ to $M$, so every $M_i$ has to be finite-dimensional.)  However, a locally finite-dimensional $M$ need not be finitely generated.
\end{remark}
 
For $d\in\Z$, the shift of Definition~\ref{def:shift} is a functor \[(d):\Grmod A\to\Grmod A.\]
It sends $f:M\to N$ to the same underlying map $f(d):M(d)\to N(d)$.
Since $M(d)(-d)=M$ and $f(d)(-d)=f$, the shift is an equivalence with inverse $(-d)$; it restricts to an equivalence $\grmod A\to\grmod A$.

\begin{remark}\label{rem:graded-exactness}
The notion of exactness makes sense for graded module categories -- a sequence in $\Grmod A$ is exact if the underlying sequence in $\Mod A$ is exact.
In particular, we talk about functors invovle such categories being exact.
\begin{supplementary}
For $L,M,N\in\Grmod A$ with maps $i:L\to M$ and $p:M\to N$, the sequence $0\to L\xrightarrow{i}M\xrightarrow{p}N\to0$ is exact if and only if
\[
0\longrightarrow L_j\xrightarrow{i|_{L_j}}M_j\xrightarrow{p|_{M_j}}N_j\longrightarrow0
\quad\text{is exact for every }j\in\Z.
\]
Indeed, $i$ and $p$ preserve the direct sum decompositions, so their kernels and images decompose into the kernels and images of the restrictions.
\end{supplementary}
\end{remark}

\begin{example}\label{ex:short-exact-row}
Let $A=\Bbbk A_2$. 
Consider $P_1:=e_{11}A$ which has $\rad P_1 = \Bbbk e_{12}$ and $L_1:=P_1/\rad P_1$.
Then we have a short exact sequence in $\Grmod A$:
\[
0\longrightarrow \rad P_1\xrightarrow{i} P_1
\xrightarrow{p}P_1/\rad P_1\longrightarrow0,
\qquad
i(c e_{12})=c e_{12},\quad p(m)=\overline m.
\]
Note that $\deg \rad P_1 = 1$ and $\deg L_1=0$.  Let $L_2:=\Bbbk e_{22} = e_{22} A$, which we treat as a graded $A$-module concentrated in degree $0$.  Note that $L_2 \cong \rad P_1$ as ungraded modules but we only have $L_2\cong (\rad P_1) (1)$, which yields the following short exact sequence in $\Grmod A$:
\[
0\longrightarrow L_2(-1)\xrightarrow{} P_1
\xrightarrow{}P_1/\rad P_1\longrightarrow0,
\]
\end{example}

\begin{lemma}\label{lem:homogeneous-hom-sum}
For graded right $A$-modules $M,N$, there is an inclusion 
\[
\bigoplus_{d\in\Z}\Hom_A^{\Z}(M,N(d))\subset\Hom_A(M,N), \text{ which is equality when $M$ is finitely generated.}
\]
\end{lemma}

\begin{proof}
For $f\in\Hom_A(M,N)$, define $f_d(m)=f(m)_{j+d}$ when $m\in M_j$, and extend linearly.
For $a\in A_i$, we have $f_d(ma)=f(ma)_{j+i+d}=f(m)_{j+d}a$, so $f_d$ is $A$-linear of degree $d$.
The sum of the homogeneous Hom spaces is direct: on $M_j$, distinct degrees have values in distinct summands of $N$.

If $M$ is finitely generated, the homogeneous components of a finite generating set form a finite homogeneous generating set.
Only finitely many $f_d$ are non-zero on these homogeneous generators, hence on $M$, giving $f=\sum_d f_d$.
\end{proof}

\subsection{Homogeneous ideals and quotients}

\begin{definition}
Let $A=\bigoplus_{i\in\Z}A_i$ be a graded algebra.
A right ideal $I\subset A$ is \defn{homogeneous} if
\[
I=\bigoplus_{i\in\Z}(I\cap A_i).
\]
Equivalently, whenever $a\in I$ and $a=a_0+\cdots+a_N$ is its decomposition into homogeneous components, each $a_i$ also belongs to $I$.
\end{definition}

For quotient algebras we need two-sided ideals.
If $I$ is a homogeneous two-sided ideal, then
\[
A/I=\bigoplus_{i\in\Z} A_i/(I\cap A_i)
\]
is again a graded algebra.

\begin{example}\label{ex:quadratic-matrix-quotient}
In $A=\Bbbk A_3$ with the grading of Example~\ref{ex:triangular-grading}, the two-sided ideal $I=(e_{13})=\Bbbk e_{13}$ is homogeneous.
Writing $\overline e_{ij}$ for the image of $e_{ij}$,
\[
(A/I)_0\cong\Bbbk^{\oplus 3},\qquad
(A/I)_1=\Bbbk\overline e_{12}\oplus\Bbbk\overline e_{23},\qquad
(A/I)_d=0\quad(d\geq 2).
\]
In particular, $\overline e_{12}\overline e_{23}=0$.
\end{example}

\begin{example}\label{ex:nonhomogeneous-ideal}
Let $A=\Bbbk[x]$ with $\deg x=1$, and let $I=(x-x^2)$.
The ideal $I$ is two-sided because $A$ is commutative.
Every element of $I$ vanishes at $x=1$, whereas neither $x$ nor $x^2$ does.
Thus $x-x^2\in I$ but $x,x^2\notin I$, so $I$ is not homogeneous.
Writing $\overline f=f+I$, we have $0\neq\overline x=\overline{x^2}$.
The images of $A_1$ and $A_2$ therefore intersect non-trivially, so $A/I$ does not inherit the quotient grading.
As an ungraded algebra, the quotient is semisimple:
\[
\Bbbk[x]/(x-x^2)\xrightarrow{\sim}\Bbbk\oplus\Bbbk,
\qquad \overline f\longmapsto(f(0),f(1)).
\]
The inverse sends $(a,b)$ to $a+(b-a)\overline x$.
\end{example}

\begin{definition}\label{def:positive-ideal}
The \defn{positive degree ideal} of a graded algebra $A$ is $A_+:=\bigoplus_{i>0}A_i$.\footnote{It is also called the augmentation ideal; we avoid that terminology.}
\end{definition}

The subspace $A_+$ is a homogeneous two-sided ideal.

\subsection{Tensor products of vector spaces}

\begin{definition}\label{def:vector-tensor}
Choose a basis $\{u_i\mid i\in I\}$ of $U$ and a basis $\{v_j\mid j\in J\}$ of $V$.
The \defn{tensor product} $U\otimes_\Bbbk V$ is the vector space with basis $\{u_i\otimes v_j\mid i\in I,\ j\in J\}$.
\end{definition}

The assignment $(u_i,v_j)\mapsto u_i\otimes v_j$ extends uniquely to a bilinear map $U\times V\to U\otimes_\Bbbk V$ sending $(u,v)\mapsto u\otimes v$.  Note that this means that elements like $(au_1+bu_2)\otimes (cv_1+dv_2)$ in $U\otimes_\Bbbk V$ is by definition the linear combination $ac(u_1\otimes v_1) + ad(u_1\otimes v_2) + bc(u_2\otimes v_1) + bd(u_2\otimes v_2)$.


% An element of the form $u\otimes v$ is called a \defn{pure tensor}.
% \comm{do not introduce useless terminology!}
\begin{proposition}[Universal property]\label{prop:vector-tensor-universal}
Every bilinear map $b:U\times V\to W$ has a unique linear factorisation (commutative diagram)
\[
\xymatrix@C=60pt@R=22pt{
U\times V \ar[r]^{(u,v)\mapsto u\otimes v} \ar[dr]_{\forall \text{ bilinear } b\phantom{...}}
& U\otimes_\Bbbk V \ar@{-->}[d]^{\exists!\,\widetilde b}\\
& W.
}
\]
% File-specific editorial rule for 2026_Koszul.tex: preserve the author's
% manual spacing adjustments, including \phantom{...} in diagram labels.
% Do not remove or change those adjustments unless explicitly requested.
Consequently, $U\otimes_\Bbbk V$ is independent of the chosen bases up to a unique isomorphism preserving the bilinear map.
\end{proposition}


\begin{supplementary}
\begin{proof}
The formula $\widetilde b(u_i\otimes v_j)=b(u_i,v_j)$ determines a unique linear map $\widetilde b:U\otimes_\Bbbk V\to W$.
Bilinearity gives $\widetilde b(u\otimes v)=b(u,v)$ for all $u\in U$ and $v\in V$.

Suppose $t:U\times V\to T$ and $t':U\times V\to T'$ both satisfy the universal property.
They give unique linear maps $a:T\to T'$ and $b:T'\to T$ with $a\circ t=t'$ and $b\circ t'=t$.
Both $b\circ a:T\to T$ and $\id_T:T\to T$ factor $t$ through itself, since
\[
(b\circ a)\circ t=b\circ t'=t=\id_T\circ t.
\]
Uniqueness for $(T,t)$ gives $b\circ a=\id_T$.
Likewise, $(a\circ b)\circ t'=a\circ t=t'=\id_{T'}\circ t'$, so uniqueness for $(T',t')$ gives $a\circ b=\id_{T'}$.
\end{proof}
\end{supplementary}

For finite bases with $|I|=r$ and $|J|=s$, we obtain $\dim_\Bbbk(U\otimes_\Bbbk V)=rs$.
We use the canonical identifications
\[
\begin{aligned}
(U\otimes_\Bbbk V)\otimes_\Bbbk W&\xrightarrow{\sim}U\otimes_\Bbbk(V\otimes_\Bbbk W),
& (u\otimes v)\otimes w&\longmapsto u\otimes(v\otimes w),\\
\Bbbk\otimes_\Bbbk U&\xrightarrow{\sim}U,
& c\otimes u&\longmapsto cu,
\end{aligned}
\]
and set $U^{\otimes_\Bbbk 0}:=\Bbbk$.


\subsection{Direct sums and direct products}\label{subsec:sums-products}

For an arbitrary family $(M_i)_{i\in I}$ of right $A$-modules, indexed by a set $I$, the \defn{direct product} and \defn{direct sum} are
\[
\begin{aligned}
\prod_{i\in I}M_i&=\{(m_i)_{i\in I}\mid m_i\in M_i\text{ for all }i\in I\},\\
\bigoplus_{i\in I}M_i&=\{(m_i)_{i\in I}\in\prod_{i\in I}M_i\mid m_i=0\text{ for all but finitely many }i\}.
\end{aligned}
\]
Addition and the $A$-action are componentwise; the same definitions apply to vector spaces by taking $A=\Bbbk$.
Thus $\bigoplus_{i\in I} M_i\subset\prod_{i\in I} M_i$, with equality for finite $I$.
For example, $(1,1,\ldots)\in\prod_{i\geq0}\Bbbk$ does not belong to $\bigoplus_{i\geq0}\Bbbk$.
In a direct sum, an element is a finite sum of its components; no notion of convergence is involved.
For instance, ordinary polynomials use non-negative powers, whereas Laurent polynomials allow integer powers, but both are finite sums:
\[
\Bbbk[x]=\bigoplus_{n\geq0}\Bbbk x^n,\qquad
\Bbbk[x,x^{-1}]=\bigoplus_{n\in\Z}\Bbbk x^n.
\]
\phantomsection\label{not:formal-power-series}
The \defn{formal power series ring} allows arbitrary coefficients:
\[
\Bbbk[[x]]=\left\{\sum_{n\geq0}a_nx^n\;\middle|\;a_n\in\Bbbk\text{ for every }n\geq0\right\}
\cong\prod_{n\geq0}\Bbbk x^n
\quad\text{as vector spaces}.
\]
\begin{supplementary}
Addition is coefficientwise, and multiplication is
\[
\Bigl(\sum_{n\geq0}a_nx^n\Bigr)\Bigl(\sum_{n\geq0}b_nx^n\Bigr)
=\sum_{n\geq0}\Bigl(\sum_{j=0}^n a_jb_{n-j}\Bigr)x^n.
\]
\end{supplementary}


Tensor products distribute over direct sums in either variable.
For vector spaces $U,V$ and arbitrary families $(U_i)_{i\in I}$, $(V_j)_{j\in J}$ indexed by sets $I,J$, there are natural linear isomorphisms
\begin{equation}\label{eq:vector-tensor-sums}
\begin{aligned}
\Bigl(\bigoplus_{i\in I}U_i\Bigr)\otimes_\Bbbk V
&\xrightarrow{\sim}\bigoplus_{i\in I}(U_i\otimes_\Bbbk V),
& (u_i)_{i\in I}\otimes v&\longmapsto(u_i\otimes v)_{i\in I},\\
U\otimes_\Bbbk\Bigl(\bigoplus_{j\in J}V_j\Bigr)
&\xrightarrow{\sim}\bigoplus_{j\in J}(U\otimes_\Bbbk V_j),
& u\otimes(v_j)_{j\in J}&\longmapsto(u\otimes v_j)_{j\in J}.
\end{aligned}
\end{equation}
The displayed bilinear assignments induce these linear maps by Proposition~\ref{prop:vector-tensor-universal}.
On the $i$th summand, the inverse of the first map sends $u_i\otimes v$ to the element with $u_i$ in coordinate $i$, tensored with $v$.
On the $j$th summand, the inverse of the second sends $u\otimes v_j$ to $u$ tensored with the element having $v_j$ in coordinate $j$.
Extending over the direct sums gives the two inverses.

\begin{advanced}
For direct products, the corresponding map need not be an isomorphism.

\begin{lemma}\label{lem:vector-tensor-products}
For an arbitrary family $(U_i)_{i\in I}$ of vector spaces, indexed by a set $I$, and a vector space $V$, the coordinate projections define a linear map
\begin{equation}\label{eq:vector-tensor-products}
\Bigl(\prod_{i\in I}U_i\Bigr)\otimes_\Bbbk V
\longrightarrow\prod_{i\in I}(U_i\otimes_\Bbbk V),
\qquad (u_i)_i\otimes v\longmapsto(u_i\otimes v)_i.
\end{equation}
This map is injective, and is an isomorphism if $I$ is finite or $V$ is finite-dimensional.
The same statements hold with the two tensor factors reversed.
\end{lemma}

\begin{proof}
The bilinear map $\bigl(\prod_{i\in I}U_i\bigr)\times V\to\prod_{i\in I}(U_i\otimes_\Bbbk V)$, $((u_i)_i,v)\mapsto(u_i\otimes v)_i$, induces \eqref{eq:vector-tensor-products} by Proposition~\ref{prop:vector-tensor-universal}.
Every element can be written $\sum_{\ell=1}^r u^{(\ell)}\otimes v_\ell$ with the $v_\ell$ linearly independent: choose a basis of the span of its finitely many second factors.
If its image is zero, then $\sum_{\ell=1}^r u_i^{(\ell)}\otimes v_\ell=0$ for every $i$.
The tensor-product basis construction gives $u_i^{(\ell)}=0$ for all $i,\ell$, proving injectivity.
If $v_1,\ldots,v_r$ is a basis of $V$, the inverse is
\[
\prod_{i\in I}(U_i\otimes_\Bbbk V)\longrightarrow
\Bigl(\prod_{i\in I}U_i\Bigr)\otimes_\Bbbk V,\qquad
\Bigl(\sum_{\ell=1}^r u_i^{(\ell)}\otimes v_\ell\Bigr)_i
\longmapsto\sum_{\ell=1}^r(u_i^{(\ell)})_i\otimes v_\ell.
\]
For finite $I$, products equal sums.
Reversing the tensor factors gives the other statements, since $U\otimes_\Bbbk V\to V\otimes_\Bbbk U$, $u\otimes v\mapsto v\otimes u$, is an isomorphism.
\end{proof}

\begin{example}\label{ex:tensor-product-product-failure}
The comparison in \eqref{eq:vector-tensor-products} need not be surjective when $V$ is infinite-dimensional.
Take $U_i=\Bbbk$ for $i\geq1$ and $V=\bigoplus_{j\geq1}\Bbbk v_j$.
Under $\Bbbk\otimes_\Bbbk V\cong V$, the image of the product comparison in \eqref{eq:vector-tensor-products} is
\[
\left\{(w_i)_i\in\prod_{i\geq1}V\;\middle|\;
\dim_\Bbbk\operatorname{span}_\Bbbk\{w_i\mid i\geq1\}<\infty\right\}.
\]
Indeed, the image of $\sum_{\ell=1}^r(c_i^{(\ell)})_i\otimes w_\ell$ has every entry in $\operatorname{span}_\Bbbk\{w_1,\ldots,w_r\}$.
Conversely, expand all entries in a basis of their common finite-dimensional span.
Since the $v_i$ are linearly independent, $(v_i)_i$ is not in the image.
\end{example}
\end{advanced}

\subsection{Bimodules and enveloping algebras}


\begin{definition}\label{def:bimodule}
For $\Bbbk$-algebras $B,C$, a \defn{bimodule} $ {}_BM_C$ is a vector space with $\Bbbk$-bilinear left $B$- and right $C$-module actions satisfying $(bm)c=b(mc)$ for all $b\in B$, $c\in C$, and $m\in M$.
The subscripts indicate the left and right acting algebras; when $B=C=A$, we call $M$ an \defn{$A$-bimodule}.
A \defn{bimodule homomorphism} $f:{}_BM_C\to{}_BN_C$ is a linear map satisfying $f(bmc)=bf(m)c$ for all $b\in B$, $c\in C$, and $m\in M$.
For graded $B,C$, a bimodule $M$ is \defn{graded} if $M=\bigoplus_{i\in\Z}M_i$ and $B_iM_jC_k\subset M_{i+j+k}$ for all $i,j,k\in\Z$.
A \defn{graded bimodule homomorphism} is a bimodule homomorphism $f:M\to N$ satisfying $f(M_i)\subset N_i$ for all $i\in\Z$.
\end{definition}

\phantomsection\label{not:bimodule-notation}
\begin{convention}
\begin{itemize}
\item Bilinearity implies that $\Bbbk$ acts centrally: $(c1_B)m=cm=m(c1_C)$ for $c\in\Bbbk$.  
From the perspective of general module theory over general rings (algebras over non-field commutative rings), we should have said that our bimodules are $\Bbbk$-central bimodule.  Since we alawys work over a field, all bimodules we considered are always $\Bbbk$-central and we will not write such adjective every time.

\item  Instead of saying $(B,C)$-bimodule, we usually just write ${}_BM_C$.
It would be natural to write $B\Mod C$ for the category of $(B,C)$-bimodules, but since $\Bbbk$-centrality comes for free from the assumption of $\Bbbk$ being a field, we can use $\Mod C\otimes_\Bbbk B^\op$ instead.
\end{itemize}
\end{convention}

\begin{example}
\begin{itemize}
\item Multiplication makes $A_A$ a graded right module and ${}_AA_A$ a graded bimodule; each degree component is a bimodule ${}_{A_0}(A_i)_{A_0}$.

\item For a bimodule ${}_AM_A$ and idempotents $e,f\in A$, the space $eMf:=\{enf\mid n\in M\}$ is a bimodule ${}_{eAe}(eMf)_{fAf}$.
\end{itemize}
\end{example}


\begin{definition}\label{def:enveloping-algebra}
For a $\Bbbk$-algebra $A$, its \defn{enveloping algebra} is $A^\en:=A\otimes_\Bbbk A^\op$, with multiplication extended bilinearly from
\[
(a\otimes b^\op)(a'\otimes(b')^\op)=aa'\otimes(b'b)^\op,
\qquad 1_{A^\en}=1_A\otimes1_A^\op
\quad(a,b,a',b'\in A).
\]
\end{definition}

A bimodule ${}_AM_A$ becomes a right $A^\en$-module by
\begin{equation}\label{eq:enveloping-action}
m\cdot(a\otimes b^\op)=bma
\quad(a,b\in A,\ m\in M).
\end{equation}
\begin{supplementary}
The first tensor factor records the right action, and the second the left action:
\[
\bigl(m\cdot(a\otimes b^\op)\bigr)\cdot(a'\otimes(b')^\op)
=b'bmaa'
=m\cdot\bigl(aa'\otimes(b'b)^\op\bigr).
\]
Conversely, for $M\in\Mod A^\en$, define $ma=m\cdot(a\otimes1)$ and $bm=m\cdot(1\otimes b^\op)$.
The two actions commute because $a\otimes1$ and $1\otimes b^\op$ commute in $A^\en$.
For bimodules ${}_AM_A$, ${}_AN_A$, a linear map $u:M\to N$ satisfies
\[
u(bma)=bu(m)a\quad(a,b\in A,\ m\in M)
\quad\Longleftrightarrow\quad
u\in\Hom_{A^\en}(M,N).
\]
Thus $A\Mod A$ is identified with $\Mod A^\en$.
More generally, for $\Bbbk$-algebras $B,C$, the vector space $C\otimes_\Bbbk B^\op$ is an algebra with
\[
(c\otimes b^\op)(c'\otimes(b')^\op)=cc'\otimes(b'b)^\op,
\qquad 1=1_C\otimes1_B^\op.
\]
The same construction identifies the categories
\[
B\Mod C\simeq\Mod(C\otimes_\Bbbk B^\op),
\qquad m\cdot(c\otimes b^\op)=bmc.
\]
\end{supplementary}

\begin{advanced}
Grant--Iyama \cite[\S2, p.~2591]{GI20} identify bimodules with left modules over $A^\en$, using $(a\otimes b^\op)m=amb$.
Our right module convention is \eqref{eq:enveloping-action}: $m\cdot(a\otimes b^\op)=bma$.
Both describe the same bimodules; the side on which $A^\en$ acts determines the formula.

\begin{remark}\label{rem:enveloping-semisimple}
For a finite-dimensional semisimple $\Bbbk$-algebra $S$, semisimplicity of $S^\en$ requires an additional condition.
For a field extension $K/\Bbbk$, give $K\otimes_\Bbbk S$ the multiplication $(c\otimes s)(d\otimes t)=cd\otimes st$.
The algebra $S$ is called \defn{separable over $\Bbbk$} if $K\otimes_\Bbbk S$ is semisimple for every field extension $K/\Bbbk$.
By \cite[Theorem~6.1.2, p.~105]{DK94},
\[
S^\en\text{ is semisimple}
\quad\Longleftrightarrow\quad
S\text{ is separable over }\Bbbk.
\]
For a product of matrix algebras, every field extension $K/\Bbbk$ gives
\[
K\otimes_\Bbbk\prod_{j=1}^n\Mat_{n_j}(\Bbbk)
\xrightarrow{\sim}\prod_{j=1}^n\Mat_{n_j}(K),
\qquad c\otimes(a_j)_j\longmapsto(ca_j)_j.
\]
The target is semisimple by Theorem~\ref{thm:artin-wedderburn}.
Thus products of matrix algebras over $\Bbbk$ are separable; in particular, every finite-dimensional semisimple algebra over an algebraically closed field is separable.
Semisimplicity of $S$ alone does not imply semisimplicity of $S^\en$ over an arbitrary field.
\end{remark}

\begin{example}\label{ex:inseparable-envelope}
Let $\Bbbk=\mathbb F_2(t)$, where $t$ is an indeterminate, and let $S=\Bbbk[u]/(u^2-t)$, also writing $u$ for its residue class.
The element $t$ is not a square in $\Bbbk$: an equality $p(t)^2=tq(t)^2$ for non-zero polynomials would give $2\deg p=1+2\deg q$.
Thus $u^2-t$ is irreducible and $S$ is a field, hence semisimple.
However, since $S$ is commutative and has characteristic $2$,
\[
S^\en=S\otimes_\Bbbk S
\cong S[v]/(v^2-t)
=S[v]/((v-u)^2)
\cong S[\varepsilon]/(\varepsilon^2).
\]
Here the first tensor factor supplies the coefficients, $v$ corresponds to $1\otimes u$, and $\varepsilon=1\otimes u-u\otimes1$.
The ideal $S\varepsilon$ is non-zero, has square zero, and is the unique maximal ideal, so $\rad(S^\en)=S\varepsilon\neq0$.
Thus $S^\en$ is not semisimple, and the field extension $S/\Bbbk$ already shows that $S$ is not separable over $\Bbbk$.
\end{example}
\end{advanced}


%%%%%%%%%%%%%%%%%%%%%%%%%%%%%%%%%%%%%%%%%%%%%%%%%%%%%%%%%%%%%%%%%%

\subsection{\texorpdfstring{Graded $\Bbbk$-linear duals}{Graded k-linear duals}}

A linear functional on a graded vector space can be non-zero on infinitely many degrees.
To obtain a graded space, retain only finite sums of homogeneous functionals.

\begin{definition}\label{def:linear-duals}
For a graded vector space $V=\bigoplus_{i\in\Z}V_i$, the \defn{graded $\Bbbk$-linear dual} is the graded vector space
\[
\bbD(V)=\bigoplus_{i\in\Z}\bbD(V)_i,
\qquad \bbD(V)_i=D(V_{-i})=\Hom_\Bbbk^\Z(V,\Bbbk(i)).
\]
Note that here $\Bbbk$ is regarded as a graded module concentrated in degree $0$.
\end{definition}

Linear duality reverses the side of a module action.
For a right $A$-module $M$ and a left $A$-module $N$, set
\begin{equation}\label{eq:linear-dual-actions}
(af)(m)=f(ma),\qquad (ga)(n)=g(an)
\quad\text{for }a\in A,\ f\in D(M),\ g\in D(N),\ m\in M,\ n\in N.
\end{equation}
For graded modules these actions restrict to the graded linear duals: if $a\in A_p$ and $f\in\bbD(M)_q$, then $af$ vanishes outside $M_{-p-q}$.


\begin{example}\label{ex:row-linear-dual}
Let $A=\Bbbk A_3$ and $P_1=e_{11}A$ as in Example~\ref{ex:shifted-row}.
The functionals $\varepsilon_{1j}(e_{1\ell})=\delta_{j\ell}$ for $1\leq j,\ell\leq 3$ form a basis of $\bbD(P_1)$, with
\[
\deg\varepsilon_{1j}=1-j,\qquad
e_{12}\varepsilon_{12}=\varepsilon_{11},\qquad
e_{23}\varepsilon_{13}=\varepsilon_{12},\qquad
e_{jj}\varepsilon_{1j}=\varepsilon_{1j}.
\]
The actions follow from \eqref{eq:linear-dual-actions}.
%For comparison with an infinite-dimensional space, take $V=\bigoplus_{i\geq 0}\Bbbk v_i$ with $\deg v_i=i$.
%The functional $\lambda(v_i)=1$ for every $i$ belongs to $D(V)$ but not to $\bbD(V)$.
\end{example}


For any $u\in \Hom_A(M,N)$, we have a dual map
\[D(u):D(N)\to D(M), \qquad f\mapsto f\circ u,
\] 
which is a left $A$-linear map, i.e. $D(u)\in \Hom_{A^\op}(D(N), D(M))$.
This gives rise to four exact (contravariant) functors (without any finiteness assumption):
\begin{equation}\label{eq:linear-dual-functors}
\begin{aligned}
D:(\Mod A)^\op&\longrightarrow\Mod(A^\op),
&D:(\Mod(A^\op))^\op&\longrightarrow\Mod A,\\
\bbD:(\Grmod A)^\op&\longrightarrow\Grmod(A^\op),
&\bbD:(\Grmod(A^\op))^\op&\longrightarrow\Grmod A.
\end{aligned}
\end{equation}
Exactness follows from Example~\ref{ex:linear-duality}, degree by degree for $\bbD$.

For these to be duality, we need to recover a module from its bidual; this requires some finiteness conditions.

\begin{lemma}\label{lem:linear-bidual}
Let $A$ be a $\Bbbk$-algebra.
\begin{enumerate}
\item The functor $D$ defines a duality between finite-dimensional right and left $A$-modules.  In particular, if $A$ is finite-dimensional, then we have a duality on finitely generated modules
\[
D: (\mod A)^\op \xrightarrow{\sim}\mod (A^\op), \qquad D: (\mod (A^\op))^\op \xrightarrow{\sim}\mod A.
\]

\item Suppose that $A$ is graded.  For locally finite-dimensional $M\in \grmod A$, evaluation is a natural graded module isomorphism
\begin{equation}\label{eq:linear-bidual-evaluation}
\mathrm{ev}_M:M\xrightarrow{\sim}\bbD(\bbD(M)),\qquad
m\longmapsto(f\mapsto f(m)).
\end{equation}
In particular, we have a duality between locally finite-dimensional graded right and left $A$-modules:
\begin{equation}\label{eq:graded-linear-dual}
\bbD:(\grmod A)^\op\xrightarrow{\sim}\grmod(A^\op),\qquad
\bbD:(\grmod(A^\op))^\op\xrightarrow{\sim}\grmod A.
\end{equation}
\end{enumerate}
\end{lemma}

\begin{proof}
First we show that $\mathrm{ev}_M$ is $A$-linear.
Indeed, 
\[
\mathrm{ev}_M(ma)(f)=f(ma)
=\mathrm{ev}_M(m)(af)
=(\mathrm{ev}_M(m)a)(f)
\quad(m\in M,\ a\in A,\ f\in D(M)).
\]

For (1), the claim follows from Example~\ref{ex:linear-duality}.
For (2), apply (1) degree-wise.
\end{proof}

\begin{advanced}
For a locally finite-dimensional graded vector space $V$, restriction gives
\[
\bbD(V)=\bigoplus_{i\in\Z}D(V_{-i})\subset
D(V)\cong\prod_{i\in\Z}D(V_i).
\]
The inclusion is equality exactly when only finitely many $V_i$ are non-zero.
For an arbitrary vector space, evaluation $V\to D(D(V))$ is injective but need not be surjective; hence the functors in \eqref{eq:linear-dual-functors} are not dualities on all modules.

\begin{example}
Graded linear duality need not preserve finite generation, even for a polynomial ring.
Let $A=\Bbbk[x]$ with $\deg x=1$ and let $\varepsilon_n(x^m)=\delta_{nm}$ for $n,m\geq0$.
Then $\deg\varepsilon_n=-n$, $x\varepsilon_0=0$, and $x\varepsilon_n=\varepsilon_{n-1}$ for $n\geq1$.
Any finite set of elements of $\bbD(A)$ lies in the span of $\varepsilon_0,\ldots,\varepsilon_N$ for some $N$, and multiplication by $x$ preserves that span.
Thus $\bbD(A)$ is locally finite-dimensional but not finitely generated.
\end{example}
\end{advanced}

For vector spaces $V,W$, a functional $f\in D(V)$ and a vector $w\in W$ define a linear map $V\to W$, $v\mapsto f(v)w$.
Finite sums of these maps give the following description of Hom.

\begin{proposition}\label{prop:linear-Hom-tensor}
For vector spaces $V,W$, there is a natural linear map
\begin{equation}\label{eq:linear-Hom-tensor}
W\otimes_\Bbbk D(V)\longrightarrow\Hom_\Bbbk(V,W),
\qquad w\otimes f\longmapsto\bigl(v\mapsto f(v)w\bigr).
\end{equation}
It is an isomorphism if either $V$ or $W$ is finite-dimensional.
\end{proposition}

\begin{supplementary}
\begin{proof}
If $V$ has a finite basis $\{v_i\mid 1\leq i\leq n\}$ with dual basis $\{v_i^*\mid 1\leq i\leq n\}$, the inverse of \eqref{eq:linear-Hom-tensor} is
\[
\Hom_\Bbbk(V,W)\longrightarrow W\otimes_\Bbbk D(V),
\qquad F\longmapsto\sum_{i=1}^nF(v_i)\otimes v_i^*,
\]
since $v=\sum_{i=1}^n v_i^*(v)v_i$.
If instead $W$ has a finite basis $\{w_j\mid 1\leq j\leq m\}$, write $F(v)=\sum_{j=1}^m f_j(v)w_j$; the uniquely determined coordinate functionals $f_j\in D(V)$ give the inverse $F\mapsto\sum_{j=1}^m w_j\otimes f_j$.
The formulas commute with precomposition and postcomposition, so the comparison is natural.
\end{proof}
\end{supplementary}

\begin{assumption}\label{ass:basic-semisimple}
For simplicity, from now on and unless stated otherwise, we assume the degree $0$ part $A_0$ of a graded algebra to be semisimple of the form
\[
S=\bigoplus_{i=1}^n\Bbbk e_i
\]
with coordinate idempotents $e_1,\ldots,e_n$.
\end{assumption}

The coordinate idempotents are the primitive idempotents of $S$, so they involve no choice.
Statements already proved for an arbitrary finite-dimensional semisimple $S$ keep their generality.
A remark headed \emph{General $S$} marks each place where an arbitrary finite-dimensional semisimple $S$ needs extra care; usually this means replacing the $\Bbbk$-linear dual by the one-sided $S$-duals of the advanced material below, or choosing a symmetrising form.
%To dualise these components while retaining both actions, we distinguish maps to $S$ that are right or left $S$-linear.
%For duality of arbitrary $A$-modules we already have the canonical functors in \eqref{eq:linear-dual-functors}.

For a bimodule ${}_SM_S$, the $\Bbbk$-linear dual retains both actions:
\begin{equation}\label{eq:linear-bimodule-dual-action}
(s\lambda t)(m)=\lambda(tms)\qquad(s,t\in S,\ \lambda\in D(M),\ m\in M).
\end{equation}
With precomposition on morphisms, we obtain functors
\[
D:(\Mod S^\en)^\op\longrightarrow\Mod S^\en,\qquad
\bbD:(\Grmod S^\en)^\op\longrightarrow\Grmod S^\en,
\]
whose restrictions to finite-dimensional bimodules, and to locally finite-dimensional graded bimodules, are dualities by Lemma~\ref{lem:linear-bidual}.
\begin{supplementary}
Evaluation preserves both actions: for $m\in M$, $\lambda\in D(M)$, and $s,t\in S$,
\[
(s\,\mathrm{ev}_M(m)\,t)(\lambda)=\mathrm{ev}_M(m)(t\lambda s)=(t\lambda s)(m)=\lambda(smt)=\mathrm{ev}_M(smt)(\lambda).
\]
Applying this in each degree gives the graded statement.
\end{supplementary}

Dualising reverses the two idempotents: for $1\leq i,j\leq n$,
\begin{equation}\label{eq:dual-blocks}
e_jD(M)e_i=\{\lambda\in D(M)\mid\lambda(e_kMe_\ell)=0\text{ unless }(k,\ell)=(i,j)\}\cong D(e_iMe_j),
\end{equation}
since $(e_j\lambda e_i)(m)=\lambda(e_ime_j)$ by \eqref{eq:linear-bimodule-dual-action}.


\begin{advanced}[suppressfirstparskip=true,innertopmargin=0pt]
\subsection{$S$-dual}

Under Assumption~\ref{ass:basic-semisimple}, the $\Bbbk$-linear dual is all we need.
For an arbitrary finite-dimensional semisimple $S$, the choice-free substitutes are the one-sided $S$-duals below; they are the duals used in \cite[\S2.7]{BGS96}.
Under Assumption~\ref{ass:basic-semisimple} they are canonically isomorphic to $D$; see Remark~\ref{rem:standard-form}.

\emph{One-sided $S$-duals.}

\begin{definition}\label{def:S-duals}
For a right $S$-module $M$ and a left $S$-module $N$, define their \defn{$S$-duals} by
\[
\begin{aligned}
M^*&:=\Hom_S(M,S)\in \Mod(S^\op),&
(sf)(m)&:=sf(m),\\
{}^*N&:=\Hom_{S^\op}(N,S) \in \Mod S,&
(gs)(n)&:=g(n)s.
\end{aligned}
\]
Here $s\in S$, $m\in M$, $n\in N$, $f\in M^*$, and $g\in{}^*N$.
\end{definition}

%The position of the star records the side of the original action: $f(ms)=f(m)s$ in $M^*$, whereas $g(sn)=sg(n)$ in $ {}^*N$.
%\comm{why repeat what is already written????}

Similar to $\Bbbk$-lienar duality, this definition extends to exact contravariant functors: for a right module map $u:M\to N$, define $u^*:N^*\to M^*$ by precomposition $u^*(f)=f\circ u$; for a left module map $u:M\to N$, the same recipe defines ${}^*u:{}^*N\to{}^*M$.
This yields the desired functors
\[
(-)^*:(\Mod S)^\op\longrightarrow\Mod(S^\op),
\qquad
{}^*(-):(\Mod(S^\op))^\op\longrightarrow\Mod S.
\]

We also have graded analogue similar to Definition~\ref{def:linear-duals}.

\begin{definition}\label{def:graded-S-duals}
For graded right and left $S$-modules $M$, respectively, define their \defn{graded $S$-duals} by
\[
M^{\grstar}:=\bigoplus_{i\in\Z}(M_{-i})^*
=\bigoplus_{i\in\Z}\Hom_S^{\Z}(M,S(i)),\qquad
{}^{\grstar}M:=\bigoplus_{i\in\Z}{}^*(M_{-i})=\bigoplus_{i\in\Z}\Hom_{S^\op}^{\Z}(M,S(i)),
\]

where each indicated summand has degree $i$.
%The shift is $S(i)$, since $S(i)_{-i}=S_0$.
\end{definition}

This also extends to a functor: for a degree $0$ right $S$-module map $u:M\to N$, define $u^{\grstar}:N^{\grstar}\to M^{\grstar}$ by
\[
(u^{\grstar})_i:(N_{-i})^*\longrightarrow(M_{-i})^*,
\qquad f\longmapsto f\circ u|_{M_{-i}}
\quad(i\in\Z).
\]
For a degree $0$ left module map $u:M\to N$, the same formula defines ${}^{\grstar}u:{}^{\grstar}N\to{}^{\grstar}M$.
With the actions from Definition~\ref{def:S-duals}, we obtain functors
\begin{equation}\label{eq:graded-S-dual-functors}
(-)^{\grstar}:(\Grmod S)^\op\longrightarrow\Grmod(S^\op),\qquad 
{}^{\grstar}(-):(\Grmod(S^\op))^\op\longrightarrow\Grmod S.
\end{equation}

Each functor in \eqref{eq:graded-S-dual-functors} restricts to locally finite-dimensional modules: replace $\Grmod$ by $\grmod$ in its source and target.
Indeed, for $R=S$ or $S^\op$ and $M\in\grmod R$,
\begin{equation}\label{eq:graded-S-dual-finiteness}
\dim_\Bbbk\Hom_R(M_{-i},R)
\leq(\dim_\Bbbk M_{-i})(\dim_\Bbbk R)<\infty
\quad(i\in\Z).
\end{equation}

\begin{lemma}\label{lem:S-bidual}
We have dualities
\[
\begin{aligned}
(-)^*:(\mod S)^\op & \xrightarrow{\sim}\mod(S^\op), & {}^*(-):(\mod(S^\op))^\op &\xrightarrow{\sim}\mod S,\\
(-)^{\grstar}:(\grmod S)^\op&\xrightarrow{\sim}\grmod(S^\op), &
{}^{\grstar}(-):(\grmod(S^\op))^\op&\xrightarrow{\sim}\grmod S.
\end{aligned}
\]
\end{lemma}

\begin{proof}
We show that evaluation gives natural isomorphisms of functors
\[
\id_{\mod S}\xrightarrow{\sim}{}^*((-)^*),
\qquad
\id_{\mod(S^\op)}\xrightarrow{\sim}({}^*(-))^*,
\]
with components
\[
\begin{aligned}
\mathrm{ev}_M:M&\xrightarrow{\sim}{}^*(M^*),
& m&\longmapsto(f\mapsto f(m))
&& (M\in\mod S),\\
\mathrm{ev}_N:N&\xrightarrow{\sim}({}^*N)^*,
& n&\longmapsto(g\mapsto g(n))
&& (N\in\mod(S^\op)).
\end{aligned}
\]
For $u:M\to N$ in $\mod S$, $m\in M$, and $f\in N^*$,
\[
\bigl({}^*(u^*)(\mathrm{ev}_M(m))\bigr)(f)
=(u^*f)(m)=f(u(m))=\mathrm{ev}_N(u(m))(f).
\]
Hence $ {}^*(u^*)\circ\mathrm{ev}_M=\mathrm{ev}_N\circ u$.
For $M=S$, evaluation is an isomorphism: its inverse sends $F\in{}^*(S^*)$ to $F(\id_S)$.
Indeed, every $f\in S^*$ satisfies $f(s)=f(1)s$, so $F(f)=f(1)F(\id_S)=f(F(\id_S))$.
Restriction to summands identifies the dual of a finite direct sum with the direct sum of the duals; evaluation therefore is an isomorphism for $S^{\oplus n}$.
By Corollary~\ref{cor:semisimple-summands}, $M\oplus L\cong S^{\oplus n}$ for some $L$.
Under the direct sum identifications, $\mathrm{ev}_{M\oplus L}=\mathrm{ev}_M\oplus\mathrm{ev}_L$, so $\mathrm{ev}_M$ is an isomorphism.
Replacing $S$ by $S^\op$ proves the left module statement.

For the graded cases, apply the ungraded case degree-wise.  Note that since $S$ is finite-dimensional, finitely generated $S$-module is the same as finite-dimensional $S$-modules.
\end{proof}

For a bimodule ${}_SM_S$, the linear and one-sided duals retain both actions:
\begin{equation}\label{eq:bimodule-dual-actions}
(s\lambda t)(m)=\lambda(tms),\qquad
(sft)(m)=sf(tm),\qquad
(sgt)(m)=g(ms)t,
\end{equation}
where $s,t\in S$, $m\in M$, $\lambda\in D(M)$, $f\in\Hom_S(M,S)$, and $g\in\Hom_{S^\op}(M,S)$.
The functionals still take values in $\Bbbk$ or $S$.
With precomposition on morphisms, we obtain
\begin{equation}\label{eq:one-sided-bimodule-dual-functors}
\begin{aligned}
D,\ (-)^*,\ {}^*(-)&:(\Mod S^\en)^\op\longrightarrow\Mod S^\en,\\
\bbD,\ (-)^{\grstar},\ {}^{\grstar}(-)&:(\Grmod S^\en)^\op\longrightarrow\Grmod S^\en.
\end{aligned}
\end{equation}
The restrictions to finite-dimensional modules, and to locally finite-dimensional graded modules, are dualities by Lemmas~\ref{lem:linear-bidual} and~\ref{lem:S-bidual}.
Evaluation preserves both actions: for $m\in M$, $f\in M^*$, and $s,t\in S$,
\[
(s\,\mathrm{ev}_M(m)\,t)(f)
=f(sm)t=f(smt)=\mathrm{ev}_M(smt)(f).
\]
The corresponding checks for $D$ and the left dual use the same evaluation formula; applying them in each degree gives the graded statements.

Recall from Definition~\ref{def:enveloping-algebra} and \eqref{eq:enveloping-action} that
$S^\en=S\otimes_\Bbbk S^\op$ and $S$-bimodules are the same thing as $S^\en$-modules via $m\cdot(s\otimes t^\op)=tms$.

\begin{definition}\label{def:en-dual}
The \defn{two-sided dual} of a bimodule ${}_SM_S$ is
\[
M^{*_{\en}}:=\Hom_{S^{\en}}(M,S^{\en}),
\qquad (sFt)(m):=(s\otimes t^\op)F(m)
\quad(s,t\in S,\ F\in M^{*_{\en}},\ m\in M).
\]
\end{definition}

In Definition~\ref{def:en-dual}, $M$ has the right action \eqref{eq:enveloping-action}, so $F(tms)=F(m)(s\otimes t^\op)$.
The output is a left $S^{\en}$-module by multiplication on the values, which gives the displayed bimodule action.
The two-sided dual takes values in $S^{\en}$, not in $S$.
On bimodule morphisms it acts by precomposition, as do $D$, $(-)^*$, and ${}^*(-)$; each is a contravariant functor on the category of $S$-bimodules.
The two-sided dual records both module actions in the target algebra $S^\en$.

\phantomsection\label{not:graded-en-dual}
For a graded bimodule ${}_SM_S$, define its \defn{graded two-sided dual} by
\[
M^{\grstar[_{\en}]}:=\bigoplus_{i\in\Z}(M_{-i})^{*_{\en}},
\qquad (M^{\grstar[_{\en}]})_i=(M_{-i})^{*_{\en}}.
\]
Definition~\ref{def:en-dual} gives a functor
\begin{equation}\label{eq:graded-bimodule-dual-functors}
(-)^{\grstar[_{\en}]}:(\Grmod S^\en)^\op\longrightarrow\Grmod S^\en.
\end{equation}
On a degree $0$ bimodule map $u:M\to N$, its degree $i$ component is
\[
(u^{\grstar[_{\en}]})_i:(N_{-i})^{*_{\en}}\longrightarrow(M_{-i})^{*_{\en}},
\qquad F\longmapsto F\circ u|_{M_{-i}}
\quad(i\in\Z).
\]
The functor restricts to locally finite-dimensional bimodules by \eqref{eq:graded-S-dual-finiteness} with $R=S^\en$.

The four duals above follow the distinction in Grant--Iyama \cite[\S2.1, pp.~2592--2593]{GI20}, with the following notation changes:
\[
\begin{array}{c|c|c}
\text{construction}&\text{Grant--Iyama}&\text{these notes}\\ \hline
\Bbbk\text{-linear dual}&M^*&D(M)\\
\text{left }S\text{-linear dual}&M^{*_\ell}&{}^*M\\
\text{right }S\text{-linear dual}&M^{*_r}&M^*\\
\text{two-sided dual}&M^\vee&M^{*_{\en}}
\end{array}
\]
The last row also changes the convention for enveloping modules.
Their $F:M\to S^\en$ is left $S^\en$-linear.
Define
\[
\sigma:S^\en\longrightarrow S^\en,\qquad
a\otimes b^\op\longmapsto b\otimes a^\op.
\]
Since $\sigma(xy)=\sigma(y)\sigma(x)$ for $x,y\in S^\en$, the map $\sigma\circ F:M\to S^\en$ satisfies
\[
(\sigma\circ F)(tms)=(\sigma\circ F)(m)(s\otimes t^\op)
\quad(s,t\in S,\ m\in M),
\]
and belongs to our $M^{*_{\en}}$.
Thus $M^\vee\to M^{*_{\en}}$, $F\mapsto\sigma\circ F$, translates their two-sided dual into our convention.
In particular, their star denotes the linear dual; our star denotes the right $S$-dual.

\subsection{Comparing the duals}

The one-sided $S$-duals and the $\Bbbk$-linear dual are naturally isomorphic after choosing a symmetrising form on $S$.



\begin{definition}\label{def:symmetrising-form}
A \defn{symmetrising form} on a finite-dimensional $\Bbbk$-algebra $S$ is a linear map $\tau:S\to\Bbbk$ satisfying
\[
\tau(st)=\tau(ts)\quad\text{for all }s,t\in S,
\qquad
\bigl(\tau(st)=0\text{ for all }t\in S\bigr)\ \Longrightarrow\ s=0.
\]
An algebra admitting such a form is \defn{symmetric}.
\end{definition}

The second condition in Definition~\ref{def:symmetrising-form} says that the bilinear form $(s,t)\mapsto\tau(st)$ is \defn{non-degenerate}: every non-zero $s\in S$ pairs non-trivially with some $t\in S$.

Every finite-dimensional semisimple $\Bbbk$-algebra admits a symmetrising form.
For products of matrix algebras over $\Bbbk$, the forms are explicit:

\begin{example}\label{ex:symmetrising-form}
The following maps are symmetrising forms:
\[
\begin{aligned}
\tau:\Bbbk^{\oplus n}&\longrightarrow\Bbbk,
& (a_1,\ldots,a_n)&\longmapsto\sum_{j=1}^n a_j,\\
\tau:\Mat_n(\Bbbk)&\longrightarrow\Bbbk,
& a&\longmapsto\Tr(a):=\sum_{i=1}^n a_{ii},\\
\text{and more generally,}\quad \tau:\prod_{j=1}^n\Mat_{n_j}(\Bbbk)&\longrightarrow\Bbbk,
& (a_1,\ldots,a_n)&\longmapsto\sum_{j=1}^n\Tr(a_j).
\end{aligned}
\]
\end{example}

Generally, any finite-dimensional semisimple algebra has a symmetrising form, by Eilenberg--Nakayama.\footnote{S.~Eilenberg and T.~Nakayama, \emph{On the dimension of modules and algebras, II. Frobenius algebras and quasi-Frobenius rings}, Nagoya Math. J. \textbf{9} (1955), 1--16, \S2, Propositions 3--5, pp.~4--5; \href{https://doi.org/10.1017/S0027763000023229}{doi:10.1017/S0027763000023229}.}
For the graded analogue, see \cite{DNN17}.
For example, let $\lambda_j:\Delta_j\to\Bbbk$ be a symmetrising form on a finite-dimensional division $\Bbbk$-algebra $\Delta_j$ for $j=1,\ldots,n$.
Then the semisimple algebra $S=\prod_{j=1}^n\Mat_{n_j}(\Delta_j)$ has a symmetrising form given by
\[
\tau:S\longrightarrow\Bbbk,\qquad
(a_1,\ldots,a_n)\longmapsto
\sum_{j=1}^n\sum_{i=1}^{n_j}\lambda_j((a_j)_{ii}).
\]

\begin{theorem}[Comparison of duals]\label{prop:dual-comparison}\label{cor:graded-dual-comparison}
Choose a symmetrising form $\tau:S\to\Bbbk$.
For $M\in\mod S$ and $N\in\mod(S^\op)$, postcomposition with $\tau$ gives $S$-module isomorphisms
\begin{equation}\label{eq:dual-comparison}
\alpha_M:M^*\longrightarrow D(M),\quad f\longmapsto\tau\circ f,
\qquad
\beta_N:{}^*N\longrightarrow D(N),\quad g\longmapsto\tau\circ g,
\end{equation}
natural in $M$ and $N$.
Applied in each degree, the same maps give natural isomorphisms $M^{\grstar}\cong\bbD(M)$ and ${}^{\grstar}N\cong\bbD(N)$ for $M\in\grmod S$ and $N\in\grmod(S^\op)$; we again write $\alpha_M$ and $\beta_N$.
\end{theorem}

We can spell this out more precisely:


\underline{\emph{Ungraded duals.}}
The isomorphisms \eqref{eq:dual-comparison} give natural isomorphisms
\begin{equation}\label{eq:dual-comparison-functors}
\begin{aligned}
\alpha:(-)^*&\xrightarrow{\sim}D
&&\text{of functors }(\mod S)^\op\longrightarrow\mod(S^\op),\\
\beta:{}^*(-)&\xrightarrow{\sim}D
&&\text{of functors }(\mod(S^\op))^\op\longrightarrow\mod S.
\end{aligned}
\end{equation}

\underline{\emph{Graded duals.}}
Composition with $\tau$ in each degree gives natural isomorphisms
\begin{equation}\label{eq:graded-dual-comparison}
\begin{aligned}
\alpha:(-)^{\grstar}&\xrightarrow{\sim}\bbD
&&\text{of functors }(\grmod S)^\op\longrightarrow\grmod(S^\op),\\
\beta:{}^{\grstar}(-)&\xrightarrow{\sim}\bbD
&&\text{of functors }(\grmod(S^\op))^\op\longrightarrow\grmod S.
\end{aligned}
\end{equation}
For $M\in\grmod S$ and $N\in\grmod(S^\op)$, the degree $i$ components are $\alpha_{M_{-i}}$ and $\beta_{N_{-i}}$, respectively.

\begin{proof}
We prove the right module statement, i.e. taking $\alpha_M(f):= \tau\circ f$ yields 
\[
\text{ungraded case } \alpha: (-)^*\xrightarrow{\sim} D, \quad \text{and}\quad \text{graded case }\alpha: (-)^{(*)}\xrightarrow{\sim} \bbD.
\]
First, non-degeneracy of $\tau$ gives an isomorphism $S\to D(S)$, $a\mapsto(s\mapsto\tau(as))$.
Hence, for $\lambda\in D(M)$, there is a unique linear map $f:M\to S$ satisfying
\begin{equation}\label{eq:dual-inverses}
\tau(f(m)s)=\lambda(ms)\quad\text{for all }m\in M,\ s\in S.
\end{equation}
For $m\in M$ and $s,t\in S$,
\[
\tau(f(mt)s)=\lambda(mts)=\tau(f(m)ts),
\]
so non-degeneracy implies $f(mt)=f(m)t$.
Taking $s=1$ in \eqref{eq:dual-inverses} gives $\tau\circ f=\lambda$; conversely, every right $S$-linear preimage of $\lambda$ must satisfy \eqref{eq:dual-inverses}.
Thus $\alpha_M$ is bijective.

Both $M^*$ and $D(M)$ are left $S$-modules, so the component $\alpha_M:M^*\to D(M)$ must preserve the left action:
\[
\alpha_M(sf)(m)=\tau(sf(m))
=\tau(f(m)s)=\tau(f(ms))
=(s\alpha_M(f))(m)
\quad(s\in S,\ f\in M^*,\ m\in M).
\]
For a right module map $u:M\to N$ and $f\in N^*$, associativity of composition gives
$\alpha_M(f\circ u)=\alpha_N(f)\circ u$, proving naturality.

For the graded comparisons, apply the ungraded natural isomorphism degree-wise. 

Same argument applies dually to $S^\op$ and $\beta$.
\end{proof}

\begin{remark}\label{rem:standard-form}
The natural isomorphisms in Theorem~\ref{prop:dual-comparison} depend on the chosen symmetrising form.
For example, every choice $c_1,\ldots,c_n\in\Bbbk\setminus\{0\}$ gives a form
$\tau(a_1,\ldots,a_n)=\sum_{i=1}^n c_i a_i$ on $\Bbbk^{\oplus n}$.
Under Assumption~\ref{ass:basic-semisimple} we always use the \defn{standard symmetrising form} $\tau_S(\sum_{i=1}^na_ie_i)=\sum_{i=1}^na_i$, which involves no choice; then $\alpha_M$ and $\beta_N$ are canonical isomorphisms.
\end{remark}

\begin{remark}[General $S$]\label{rem:general-S-comparison}
For an arbitrary finite-dimensional semisimple $S$, identifying an $S$-dual with a $\Bbbk$-linear dual requires choosing a symmetrising form.
If $\tau'$ is another one, then $\tau'(s)=\tau(zs)$ for a central unit $z\in S$, and the two identifications differ by multiplication by $z$; see Exercise~\ref{ex:transported-dual-action}(3).
\end{remark}

\begin{remark}\label{rem:transported-dual-action}
Let $A$ be a graded algebra with $A_0=S$.
The space $\Hom_S(M,S)$ has no natural left $A$-module structure on its own, but acquires one after \emph{picking a symmetrising form $\tau$ of $S$}; under Assumption~\ref{ass:basic-semisimple} we pick $\tau=\tau_S$.
Indeed, the graded comparison $\alpha_M$ of Theorem~\ref{prop:dual-comparison} transports the canonical left $A$-action on $\bbD(M)$ to the graded $S$-dual of a locally finite-dimensional module:
\[
af=\alpha_M^{-1}(a\alpha_M(f))\quad\text{for right }M,
\qquad ga=\beta_N^{-1}(\beta_N(g)a)\quad\text{for left }N,
\]
where $a\in A$, $f\in M^{\grstar}$, and $g\in{}^{\grstar}N$.
Similarly for right action on left dual.
%For comparison, a bimodule ${}_AX_S$ gives a bimodule ${}_S\bigl(\Hom_S(X,S)\bigr)_A$ by$(sfa)(x)=s f(ax)$.
%This is the construction used in \cite[\S2.2]{BGS96}; it uses commuting actions on opposite sides, not merely restriction of a right $A$-module to $S$.
%Exercise~\ref{ex:transported-dual-action} compares two transported actions on one graded $S$-dual.
\end{remark}


The comparison theorem has a bimodule version:

\begin{theorem}
With the actions in \eqref{eq:bimodule-dual-actions}, we obtain natural isomorphisms of functors
\begin{equation}\label{eq:bimodule-dual-comparison}
\begin{aligned}
(-)^*\xrightarrow[\alpha]{\sim}D\xleftarrow[\beta]{\sim}{}^*(-)
&\quad:(\mod S^\en)^\op\longrightarrow\mod S^\en,\\
(-)^{\grstar}\xrightarrow[\alpha]{\sim}\bbD\xleftarrow[\beta]{\sim}{}^{\grstar}(-)
&\quad:(\grmod S^\en)^\op\longrightarrow\grmod S^\en.
\end{aligned}
\end{equation}

There is also a natural isomorphism $\omega:(-)^{*_\en}\xrightarrow{\sim}(-)^*$ given by
\begin{equation}\label{eq:en-dual-comparison}
\omega_M:M^{*_{\en}}\to M^*, \qquad F\mapsto \big(m\mapsto \sum_j a_j\tau(b_j)\big)
\quad\text{for}\quad
F(m)=\sum_j a_j\otimes b_j^\op,
\end{equation}
which also depends on $\tau$, but does not require $S^\en$ to be semisimple.

Consequently, on finite-dimensional $S$-bimodules and lfd graded $S$-bimodules we have natural isomorphisms:
\begin{align*}
\text{Ungraded:}\quad &(-)^{*_{\en}}\xrightarrow[\omega]{\sim}(-)^*
\xrightarrow[\alpha]{\sim}D\xleftarrow[\beta]{\sim}{}^*(-) \text{ as functors } (\mod S^\en)^\op\to\mod S^\en, \\
\text{Graded:}\quad  &(-)^{\grstar[_{\en}]}\xrightarrow[\omega]{\sim}(-)^{\grstar}
\xrightarrow[\alpha]{\sim}\bbD\xleftarrow[\beta]{\sim}{}^{\grstar}(-) \text{ as functors } (\grmod S^\en)^\op\to\grmod S^\en.
\end{align*}
\end{theorem}
\begin{proof}
The form $\tau_{\en}(a\otimes b^\op)=\tau(a)\tau(b)$ is symmetrising on $S^{\en}$.
Indeed, $\tau_{\en}((a\otimes b^\op)(c\otimes d^\op))=\tau(ac)\tau(db)=\tau(ca)\tau(bd)$, and its pairing is the tensor product of two non-degenerate pairings.
The formula \eqref{eq:en-dual-comparison} satisfies
\[
\omega_M(F)(ms)=\omega_M(F)(m)s,\qquad
\omega_M(sFt)=s\omega_M(F)t,\qquad
\alpha_M(\omega_M(F))=\tau_{\en}\circ F.
\]
The argument \eqref{eq:dual-inverses}, applied to $S^{\en}$ and $\tau_{\en}$, makes $F\mapsto\tau_{\en}\circ F$ an isomorphism.
Thus $\omega_M$ is an isomorphism.
For a bimodule map $u:M\to N$ and $F\in N^{*_{\en}}$, the formula \eqref{eq:en-dual-comparison} gives $\omega_M(F\circ u)=\omega_N(F)\circ u$, proving naturality.
Taking the direct sum of $\omega_{M_{-i}}$ over $i\in\Z$ gives the graded two-sided comparison.
Each component is a bimodule isomorphism, and naturality follows in each degree from the ungraded comparison.
\end{proof}
\end{advanced}


\subsection{Exercises}

\begin{exercise}[Matrix calculations]
Let $n\geq2$ and $A=\Bbbk A_n$.
Verify the grading in Example~\ref{ex:triangular-grading} using $e_{ij}e_{pq}=\delta_{jp}e_{iq}$.
Compute $\rad A$ from its definition as the intersection of all maximal right ideals.

\emph{Hint.} The strictly upper triangular ideal $J$ satisfies $J^n=0$.
Show that $LJ=0$ for every simple right $A$-module $L$, and use the diagonal maps $A\to\Bbbk$, $X\mapsto X_{ii}$, for $1\leq i\leq n$.
\end{exercise}

\begin{exercise}[Shifts and duals]
Let $A=\Bbbk A_3$ and $M=e_{22}A$ with the inherited grading.
\begin{enumerate}
\item Find a homogeneous basis of $\bbD(M)$, its degrees, and all non-zero actions of matrix units on the basis vectors.
\item Show that $\Bbbk e_{13}\subset Ae_{33}$ is a graded left submodule and that
\[
\bbD(M)\cong(Ae_{33}/\Bbbk e_{13})(1)
\]
as graded left $A$-modules.
\item Deduce that $\bbD(M)\not\cong(Ae_{jj})(r)$ for every $j\in\{1,2,3\}$ and $r\in\Z$.
\end{enumerate}
\emph{Hint for (3).} Compare dimensions and the actions of $e_{11},e_{22},e_{33}$.
\end{exercise}

\begin{exercise}[Homogeneous maps]\label{ex:infinite-hom}
The finite-generation hypothesis in Lemma~\ref{lem:homogeneous-hom-sum} is needed.
Let $S=\Bbbk e_1\oplus\Bbbk e_2$ with componentwise multiplication, let $N=e_1S$ have degree $0$, and put
\[
M=\bigoplus_{i\geq0}(e_1S)(-i)=\bigoplus_{i\geq0}\Bbbk m_i,
\qquad \deg m_i=i,\quad m_i e_1=m_i,\quad m_i e_2=0.
\]
For the $S$-linear map $f:M\to N$ with $f(m_i)=e_1$ for every $i\geq0$, determine all homogeneous components and prove that $f$ is not a finite sum of homogeneous maps.
\end{exercise}

\begin{hidden}
\begin{proof}[Solution]
For a finitely supported family $(c_i)_{i\geq0}$,
\[
f\Bigl(\sum_{i\geq0}c_i m_i\Bigr)=\sum_{c_i\neq0}c_i e_1.
\]
On the summand $(e_1S)(-i)$, the restriction $f|_{(e_1S)(-i)}:(e_1S)(-i)\to e_1S$ is the underlying identity map, of degree $-i$.
Explicitly, the homogeneous components are
\[
f_{-r}:M\longrightarrow N,\qquad f_{-r}(m_i)=\delta_{ri}e_1
\quad\text{for }r,i\geq0,
\]
and $f_d=0$ for $d>0$.
A homogeneous map $M\to N$ of degree $d$ vanishes on $m_i$ unless $d=-i$.
Every finite sum of homogeneous maps therefore vanishes on all but finitely many $m_i$, whereas $f(m_i)=e_1$ for every $i\geq0$.
Hence $f$ is not a finite sum of homogeneous maps.
\end{proof}
\end{hidden}

\begin{advanced}
\begin{exercise}[Dependence on the form]\label{ex:dual-blocks}
Let $S=(\Bbbk A_2)_0=\Bbbk e_{11}\oplus\Bbbk e_{22}$, $V=(\Bbbk A_2)_1=\Bbbk e_{12}$, and $c_1,c_2\in\Bbbk\setminus\{0\}$.
Set $\tau(ae_{11}+be_{22})=c_1a+c_2b$ and let $\varepsilon:V\to\Bbbk$ satisfy $\varepsilon(e_{12})=1$.
Find the right and left $S$-linear maps corresponding to $\varepsilon$ under \eqref{eq:dual-comparison}, and determine how they depend on $c_1,c_2$.
\end{exercise}

\begin{hidden}
\begin{proof}[Solution]
Right $S$-linearity forces $f(e_{12})=f(e_{12})e_{22}\in\Bbbk e_{22}$, and left $S$-linearity forces $g(e_{12})=e_{11}g(e_{12})\in\Bbbk e_{11}$.
Thus $\tau\circ f=\varepsilon=\tau\circ g$ requires
\[
f(e_{12})=c_2^{-1}e_{22},\qquad g(e_{12})=c_1^{-1}e_{11}.
\]
Conversely, the linear maps with these values satisfy, for $s=ae_{11}+be_{22}\in S$,
\[
f(e_{12}s)=bc_2^{-1}e_{22}=f(e_{12})s,\qquad
g(se_{12})=ac_1^{-1}e_{11}=sg(e_{12}).
\]
Their compositions with $\tau$ take $e_{12}$ to $1$, as required.
The original worked example is the case $c_1=c_2=1$.
\end{proof}
\end{hidden}

\begin{exercise}[Dependence of the action on the form]\label{ex:transported-dual-action}
The transported $A$-action can depend on the symmetrising form, but its graded isomorphism class does not.
Let $A=\Bbbk A_2$, $S=\Bbbk e_{11}\oplus\Bbbk e_{22}$, and $M=e_{11}A$, with the inherited gradings.
For $c=(c_1,c_2)\in(\Bbbk\setminus\{0\})^2$, define
\begin{equation}\label{eq:transported-dual-form}
\tau_c:S\longrightarrow\Bbbk,\qquad
a e_{11}+b e_{22}\longmapsto c_1a+c_2b
\quad\text{for all }a,b\in\Bbbk.
\end{equation}
Write $\alpha_c:M^{\grstar}\to\bbD(M)$, $f\mapsto\tau_c\circ f$, and let $a\cdot_c f=\alpha_c^{-1}(a\alpha_c(f))$ be the transported action of Remark~\ref{rem:transported-dual-action}.
Consider the right $S$-linear maps
\begin{equation}\label{eq:transported-dual-functionals}
\begin{aligned}
\varphi_0:M&\longrightarrow S,& a e_{11}+b e_{12}&\longmapsto a e_{11},\\
\varphi_1:M&\longrightarrow S,& a e_{11}+b e_{12}&\longmapsto b e_{22}
\quad\text{for all }a,b\in\Bbbk.
\end{aligned}
\end{equation}
\begin{enumerate}
\item Verify that $\tau_c$ is symmetrising and that $\varphi_0,\varphi_1$ form a homogeneous basis of $M^{\grstar}$. Determine their degrees.
\item Calculate the actions of $e_{11},e_{22},e_{12}$ on $\varphi_0,\varphi_1$.
For $c,c'\in(\Bbbk\setminus\{0\})^2$, prove that $\cdot_c=\cdot_{c'}$ if and only if $c'=\rho c$ for some $\rho\in\Bbbk\setminus\{0\}$.
\item More generally, let $A$ be a locally finite-dimensional graded algebra with semisimple $S=A_0$, and let $M$ be a locally finite-dimensional graded right $A$-module.
Let $\tau,\tau'$ be symmetrising forms on $S$, with comparison maps $\alpha_\tau,\alpha_{\tau'}:M^{\grstar}\to\bbD(M)$ given by composition with $\tau,\tau'$, respectively, and transported actions $\cdot_\tau,\cdot_{\tau'}$.
Prove that there is an invertible $z\in S$ satisfying $\tau'(s)=\tau(zs)$ for every $s\in S$, and that
\[
\gamma:(M^{\grstar},\cdot_{\tau'})\longrightarrow(M^{\grstar},\cdot_\tau),
\qquad (\gamma f)(m)=z f(m)
\]
is an isomorphism of graded left $A$-modules.
Deduce that the modules in (2) are isomorphic for every $c,c'$, and give the isomorphism on $\varphi_0,\varphi_1$.
\end{enumerate}
\emph{Hint for (3).} Compare $\alpha_{\tau'}$ with $\alpha_\tau$.
\end{exercise}

\begin{hidden}
\begin{proof}[Solution]
\emph{The form and the basis.}
The algebra $S$ is commutative, so $\tau_c(st)=\tau_c(ts)$.
If $s=ae_{11}+be_{22}$ satisfies $\tau_c(se_{11})=\tau_c(se_{22})=0$, then $c_1a=c_2b=0$, hence $s=0$.
Thus $\tau_c$ is non-degenerate.

A right $S$-linear map $f:M\to S$ must satisfy $f(e_{11})\in Se_{11}=\Bbbk e_{11}$ and $f(e_{12})\in Se_{22}=\Bbbk e_{22}$.
Conversely, any choices of these two values define such a map.
Consequently, $\varphi_0,\varphi_1$ form a basis, with degrees $0,-1$, respectively.

\emph{The actions.}
For $m=ae_{11}+be_{12}$,
\[
\alpha_c(\varphi_0)(m)=c_1a,\qquad
\alpha_c(\varphi_1)(m)=c_2b,\qquad
\alpha_c(e_{12}\cdot_c\varphi_1)(m)
=\alpha_c(\varphi_1)(me_{12})=c_2a.
\]
Using also $me_{11}=ae_{11}$ and $me_{22}=be_{12}$ gives all non-zero actions:
\[
e_{11}\cdot_c\varphi_0=\varphi_0,\qquad
e_{22}\cdot_c\varphi_1=\varphi_1,\qquad
e_{12}\cdot_c\varphi_1=\frac{c_2}{c_1}\varphi_0.
\]
All other actions of the three matrix units on the two basis vectors are zero.
Consequently, $\cdot_c=\cdot_{c'}$ if and only if $c_2/c_1=c'_2/c'_1$.
This equality is equivalent to $c'=\rho c$ with $\rho=c'_1/c_1\neq0$.

\emph{The general isomorphism.}
Non-degeneracy of $\tau$ makes
\[
S\longrightarrow D(S),\qquad z\longmapsto\bigl(s\mapsto\tau(zs)\bigr)
\]
injective, hence bijective since $S$ is finite-dimensional.
Choose $z$ mapping to $\tau'$.
If $zs=0$, then $\tau'(st)=\tau(zst)=0$ for every $t\in S$, so $s=0$ by non-degeneracy of $\tau'$.
Thus multiplication by $z$ on the left is bijective.
Choose $u\in S$ with $zu=1$; injectivity applied to $z(uz-1)=0$ gives $uz=1$.

For $f\in M^{\grstar}$, $m\in M$, and $s\in S$, one has $zf(ms)=zf(m)s$, so $\gamma f$ is right $S$-linear.
Multiplication by $z$ preserves degrees, and multiplication by $z^{-1}$ gives the inverse of $\gamma$.
Moreover, $\alpha_\tau\gamma=\alpha_{\tau'}$, so for every $a\in A$ and $f\in M^{\grstar}$,
\[
\gamma(a\cdot_{\tau'}f)
=\alpha_\tau^{-1}\bigl(a\alpha_{\tau'}(f)\bigr)
=a\cdot_\tau\gamma(f).
\]
Hence $\gamma$ is a graded left $A$-module isomorphism.
For the forms $\tau_c,\tau_{c'}$, take
\[
z=\frac{c'_1}{c_1}e_{11}+\frac{c'_2}{c_2}e_{22}.
\]
The resulting isomorphism $(M^{\grstar},\cdot_{c'})\to(M^{\grstar},\cdot_c)$ sends
$\varphi_0\mapsto(c'_1/c_1)\varphi_0$ and $\varphi_1\mapsto(c'_2/c_2)\varphi_1$.
\end{proof}
\end{hidden}
\end{advanced}

\pagebreak


% !TEX root = Lecture_Koszul.tex
\section{Tensor and quadratic algebras}

% TODO (ChangeProposal9, Lecture 2 follow-up): Introduce tensor products over
% arbitrary algebras, with the defining identities explicit. References now
% use prop:vector-tensor-universal for the universal property and
% def:vector-tensor for the basis construction. Keep routine
% verification supplementary; retain the existing editorial comments below.

\subsection{Tensor products over an algebra}

Recall the tensor product $M\otimes_A N$ from Definition~\ref{def:module-tensor}, for a right $A$-module $M$ and a left $A$-module $N$.
Its defining identity is $ma\otimes_A n=m\otimes_A an$ for all $m\in M$, $n\in N$, and $a\in A$.

%No semisimplicity or finiteness assumption on $S,M,N$ is needed for the following properties.

\begin{lemma}[Basic properties of tensor products]\label{lem:tensor-properties}
Let $A$ be any $\Bbbk$-algebra, $M$ a right $A$-module, and $N$ a left $A$-module.
\begin{enumerate}
\item For a vector space $W$ and a bilinear map $b:M\times N\to W$ satisfying $b(ma,n)=b(m,an)$ for all $m\in M$, $a\in A$, and $n\in N$, there is a unique linear factorisation (commutative diagram)
\[
\xymatrix@C=75pt@R=22pt{
M\times N \ar[r]^{(m,n)\mapsto m\otimes_A n} \ar[dr]_{\forall \text{ bilinear }b \text{ s.t. } b(-\cdot a,-)=b(-,a\cdot-)\phantom{aaaaaaaaaa}}
& M\otimes_A N \ar@{-->}[d]^{\exists!\,\widetilde b}\\
& W.
}
\]
%Conversely, every linear map $M\otimes_A N\to W$ defines such a bilinear map.
%not sure what this line wants to say really. just ignore it

\item If $B,C$ are $\Bbbk$-algebras and ${}_BM_A$, ${}_AN_C$ are bimodules, then $M\otimes_A N$ is a bimodule ${}_B(M\otimes_A N)_C$ with
\[
b(m\otimes_A n)c=bm\otimes_A nc
\qquad(b\in B,\ c\in C,\ m\in M,\ n\in N).
\]
For right $A$-module maps $u:M\to M'$ and left $A$-module maps $v:N\to N'$, the assignment $(m,n)\mapsto u(m)\otimes v(n)$ induces
\[
u\otimes_A v:M\otimes_A N\longrightarrow M'\otimes_A N',
\qquad m\otimes n\longmapsto u(m)\otimes v(n).
\]
These define functors
\[
(-)\otimes_A N:\Mod A\longrightarrow\Mod C,\qquad
M\otimes_A(-):\Mod(A^\op)\longrightarrow\Mod(B^\op).
\]
On bimodule maps, they retain the outer actions:
\[
(-)\otimes_A N: B\Mod A \longrightarrow B\Mod C,\qquad
M\otimes_A(-):A\Mod C\longrightarrow B\Mod C.
\]
\item For $\Bbbk$-algebras $A,B$, a right $A$-module $U$, a bimodule ${}_AV_B$, and a left $B$-module $W$, there is a natural linear isomorphism
\[
(U\otimes_A V)\otimes_B W\xrightarrow{\sim}U\otimes_A(V\otimes_B W),
\qquad (u\otimes v)\otimes w\longmapsto u\otimes(v\otimes w).
\]
There are also natural module isomorphisms
\[
M\otimes_A A\xrightarrow{\sim}M,\quad m\otimes a\longmapsto ma,
\qquad
A\otimes_A N\xrightarrow{\sim}N,\quad a\otimes n\longmapsto an.
\]
Moreover, these maps are linear with respect to any additional bimodule action on $U, W, M, N$.

\item Tensor products distribute over direct sums in either variable. For arbitrary families $(M_i)_{i\in I}$ of right $A$-modules and $(N_j)_{j\in J}$ of left $A$-modules, indexed by sets $I,J$,
\[
\begin{aligned}
\Bigl(\bigoplus_{i\in I}M_i\Bigr)\otimes_A N
&\xrightarrow{\sim}\bigoplus_{i\in I}(M_i\otimes_A N),
& (m_i)_{i\in I}\otimes n&\longmapsto(m_i\otimes n)_{i\in I},\\
M\otimes_A\Bigl(\bigoplus_{j\in J}N_j\Bigr)
&\xrightarrow{\sim}\bigoplus_{j\in J}(M\otimes_A N_j),
& m\otimes(n_j)_{j\in J}&\longmapsto(m\otimes n_j)_{j\in J}.
\end{aligned}
\]

\item For finite sets $I,J$ and right $A$-modules $M_j,N_i$, there is a natural isomorphism
\begin{equation}\label{eq:hom-finite-sums}
\begin{aligned}
\Hom_A\Bigl(\bigoplus_{j\in J}M_j,\bigoplus_{i\in I}N_i\Bigr)
&\xrightarrow{\sim}\bigoplus_{\substack{i\in I\\j\in J}}\Hom_A(M_j,N_i),\\
f&\longmapsto(\pi_i f\iota_j)_{i\in I,\ j\in J},
\end{aligned}
\end{equation}
where $\iota_j:M_j\to\bigoplus_{j'\in J}M_{j'}$ and $\pi_i:\bigoplus_{i'\in I}N_{i'}\to N_i$ are the inclusions and projections.
\end{enumerate}
\end{lemma}

\begin{supplementary}
\begin{proof}
For (1), Proposition~\ref{prop:vector-tensor-universal} gives a unique linear map $M\otimes_\Bbbk N\to W$, $m\otimes n\mapsto b(m,n)$.
The map vanishes on the defining subspace precisely when $b(ma,n)=b(m,an)$.

For (2), the outer actions preserve that subspace: for $b\in B$, $c\in C$, and $a\in A$,
\[
b(ma\otimes n-m\otimes an)c
=(bm)a\otimes nc-bm\otimes a(nc).
\]
The module axioms follow from those on $M$ and $N$.
For module maps $u,v$,
\[
u(ma)\otimes v(n)=u(m)a\otimes v(n)
=u(m)\otimes av(n)=u(m)\otimes v(an),
\]
so the displayed formula induces $u\otimes_A v$.
The functor identities follow from
$(u'\otimes v')(u(m)\otimes v(n))=u'u(m)\otimes v'v(n)$ and
$(\id_M\otimes\id_N)(m\otimes n)=m\otimes n$.
If $u,v$ preserve the outer actions, then $u\otimes_A v$ does as well.

For (3), both parenthesisations impose
\[
(ua)\otimes v\otimes w=u\otimes(av)\otimes w,\qquad
u\otimes(vb)\otimes w=u\otimes v\otimes(bw)
\quad(a\in A,\ b\in B).
\]
Thus the displayed associativity map and the reverse assignment are well-defined and inverse.
The unit inverses are $M\to M\otimes_A A$, $m\mapsto m\otimes1$, and $N\to A\otimes_A N$, $n\mapsto1\otimes n$.

For (4), the bilinear map $((m_i)_{i\in I},n)\mapsto(m_i\otimes n)_{i\in I}$ satisfies the identity in (1), so it induces the displayed linear map.
For each $i\in I$, let $\iota_i:M_i\to\bigoplus_{j\in I}M_j$ be the inclusion.
The inverse sends $m\otimes n$ in the $i$th summand to $\iota_i(m)\otimes n$ and extends over the direct sum.
Both composites fix every element $m\otimes n$ from a summand.
The other variable is handled in the same way.
\end{proof}
\end{supplementary}


\begin{advanced}
\begin{lemma}[Hom, direct sums, and direct products]\label{lem:hom-sums-products}
Let $M,N$ be right $A$-modules, and let $(M_i)_{i\in I}$ and $(N_i)_{i\in I}$ be arbitrary families of right $A$-modules, indexed by a set $I$.
Write $\iota_i:M_i\to\bigoplus_{j\in I}M_j$ for the inclusions and $\pi_i:\prod_{j\in I}N_j\to N_i$ for the projections.
There are natural linear isomorphisms
\[
\begin{aligned}
\Hom_A\Bigl(\bigoplus_{i\in I}M_i,N\Bigr)&\xrightarrow{\sim}\prod_{i\in I}\Hom_A(M_i,N),
& f&\longmapsto(f\circ\iota_i)_{i\in I},\\
\Hom_A\Bigl(M,\prod_{i\in I}N_i\Bigr)&\xrightarrow{\sim}\prod_{i\in I}\Hom_A(M,N_i),
& f&\longmapsto(\pi_i\circ f)_{i\in I}.
\end{aligned}
\]
There is also a natural injection
\[
\bigoplus_{i\in I}\Hom_A(M,N_i)\longrightarrow
\Hom_A\Bigl(M,\bigoplus_{i\in I}N_i\Bigr),
\qquad (f_i)_{i\in I}\longmapsto\bigl(m\mapsto(f_i(m))_{i\in I}\bigr),
\]
which is an isomorphism when $M$ is finitely generated.
\end{lemma}

\begin{proof}
The inverse of the first map sends $(f_i)_{i\in I}$ to $(m_i)_{i\in I}\mapsto\sum_{i\in I}f_i(m_i)$; the sum is finite because $(m_i)_{i\in I}$ has finite support.
The inverse of the second sends $(f_i)_{i\in I}$ to $m\mapsto(f_i(m))_{i\in I}$; a direct product imposes no support condition.
For the third map, coordinate projections recover each $f_i$, proving injectivity.
If $m_1,\ldots,m_r$ generate $M$, the images of these generators under $f:M\to\bigoplus_{i\in I}N_i$ are supported on a common finite subset $J\subset I$.
Then $f(M)\subset\bigoplus_{i\in J}N_i$, so $f$ comes from a finite family of maps.
All three constructions commute with precomposition and postcomposition.
\end{proof}

\begin{proposition}[Tensor products and direct products]\label{prop:tensor-direct-products}
Let $A$ be a $\Bbbk$-algebra, $M,M_i$ right $A$-modules, and $N,N_i$ left $A$-modules, with $i\in I$.
The coordinate projections define linear maps
\begin{equation}\label{eq:module-tensor-products}
\begin{aligned}
\kappa_N:\Bigl(\prod_iM_i\Bigr)\otimes_A N
&\longrightarrow\prod_i(M_i\otimes_A N),
& (m_i)_i\otimes n&\longmapsto(m_i\otimes n)_i,\\
\lambda_M:M\otimes_A\Bigl(\prod_iN_i\Bigr)
&\longrightarrow\prod_i(M\otimes_A N_i),
& m\otimes(n_i)_i&\longmapsto(m\otimes n_i)_i.
\end{aligned}
\end{equation}
\begin{enumerate}
\item If $I$ is finite, both maps are isomorphisms.
\item The map $\kappa_N$ is an isomorphism if $N$ is a direct summand of $A^{\oplus r}$ as a left module for some finite $r$.
The corresponding condition on the right module $M$ makes $\lambda_M$ an isomorphism.
\item If $A$ is finite-dimensional semisimple, both maps are injective.
Moreover, $\kappa_N$ is an isomorphism when $N$ is finite-dimensional, and $\lambda_M$ is an isomorphism when $M$ is finite-dimensional.
\end{enumerate}
\end{proposition}

\begin{proof}
The formulas in \eqref{eq:module-tensor-products} are bilinear. For $s\in A$, the component identities
$(m_is)\otimes n=m_i\otimes sn$ and $(ms)\otimes n_i=m\otimes sn_i$
show that both formulas give well-defined maps from the indicated tensor products.
For finite $I$, products equal sums, and Lemma~\ref{lem:tensor-properties}(4) proves (1).

For $N=A$, the map $\kappa_A$ is the identity on $\prod_iM_i$ under $m\otimes s\mapsto ms$.
For $N=A^{\oplus r}$, both sides are $(\prod_iM_i)^{\oplus r}$, since a finite direct sum commutes with a direct product.
If $N\oplus N'\cong A^{\oplus r}$, the comparison for $N\oplus N'$ is $\kappa_N\oplus\kappa_{N'}$.
Its bijectivity implies the bijectivity of $\kappa_N$.
This proves the assertion about $\kappa_N$ in (2).

Now let $A$ be finite-dimensional semisimple.
For finite-dimensional $N$, Corollary~\ref{cor:semisimple-summands} gives (2).
For arbitrary $N$, take $z=\sum_{\ell=1}^r m^{(\ell)}\otimes n_\ell$ in $\Ker\kappa_N$ and set $N_0=\sum_{\ell=1}^r An_\ell$.
The submodule $N_0$ is finite-dimensional and has a complement in $N$ by Theorem~\ref{thm:artin-wedderburn}.
Thus the inclusions
\[
\Bigl(\prod_iM_i\Bigr)\otimes_A N_0
\longrightarrow\Bigl(\prod_iM_i\Bigr)\otimes_A N,
\qquad
M_i\otimes_A N_0\longrightarrow M_i\otimes_A N
\]
are injective, by the direct sum formula in Lemma~\ref{lem:tensor-properties}(4).
Regard $z$ as an element with second factors in $N_0$.
The equality $\kappa_N(z)=0$ then gives $\kappa_{N_0}(z)=0$.
Aince $\kappa_{N_0}$ is an isomorphism, $z=0$.

For $\lambda_M$, apply the same arguments over $A^\op$, using
$N\otimes_{A^\op}M\cong M\otimes_A N$, $n\otimes m\mapsto m\otimes n$.
\end{proof}

The finite-dimensional hypothesis on the fixed factor cannot be dropped from the isomorphism assertion: Example~\ref{ex:tensor-product-product-failure} gives a non-surjective comparison even for vector spaces.
Thus tensor products always distribute over direct sums, but not over arbitrary direct products.

\begin{remark}\label{rem:infinite-products}
For infinite families, $\Hom$ need not preserve direct sums in its second variable.
For $M=\bigoplus_{i\geq1}\Bbbk v_i$, the identity $M\to M$ does not come from $\bigoplus_i\Hom_\Bbbk(M,\Bbbk v_i)$: its projection onto every $\Bbbk v_i$ is non-zero.
Thus the injection in Lemma~\ref{lem:hom-sums-products} need not be surjective.

A product in the source of $\Hom$ behaves differently from a product in the target.
For example, a linear map $\prod_{i\geq1}\Bbbk\to\Bbbk$ need not be determined by its restrictions to the coordinate copies of $\Bbbk$.
Indeed, $(1,1,\ldots)\notin\bigoplus_{i\geq1}\Bbbk$.
Define a functional on $\bigoplus_{i\geq1}\Bbbk+\Bbbk(1,1,\ldots)$ to vanish on the direct sum and take value $1$ at $(1,1,\ldots)$, and extend it to the product by extending a basis.
The resulting non-zero functional vanishes on every coordinate copy.
\end{remark}
\end{advanced}



\phantomsection\label{not:tensor-powers}
For a bimodule ${}_AV_A$, define
\[
V^{\otimes_A 0}:=A,\qquad
V^{\otimes_A i}:=
\underbrace{V\otimes_A\cdots\otimes_A V}_{i\text{ factors}}
\quad(i\geq1).
\]
Note that there is no parenthesis in the definition of $V^{\otimes_A i}$ since Lemma~\ref{lem:tensor-properties}(3) says that tenso product is an associatitive operation.

\begin{advanced}
Let $A,B,C$ be $\Bbbk$-algebras.
For bimodules ${}_AM_B$, ${}_CX_B$, we have $\Hom_B(M,X)\in C\Mod A$, with
\begin{equation}\label{eq:general-tensor-dual-actions}
(cfa)(m)=cf(am)
\quad(f\in\Hom_B(M,X),\ c\in C,\ a\in A,\ m\in M).
\end{equation}
If $X\in\Mod B$ has no specified left action, the same formula defines the right $A$-action $(fa)(m)=f(am)$.
The left-module version follows by passing to opposite algebras.

\begin{proposition}[Tensor and Hom]\label{prop:tensor-duals}
% TODO (ChangeProposal9, Lecture 2 follow-up): Explain the general case for
% an (A,B)-bimodule M and a (B,C)-bimodule N. For right duals, require M_B
% finitely generated projective; for left duals, require _B N finitely generated
% projective. Finite k-dimension alone does not suffice over a general algebra.
% Keep the idempotent-matching calculation in the main exposition and move
% routine verification to supplementary material. Use the one-sided bimodule
% actions from eq:bimodule-dual-actions, independently of a symmetrising form.
Consider bimodules ${}_AM_B$, ${}_BN_C$, a right $C$-module $X$, and a left $A$-module $Y$.
There are maps, natural in $M,N$,
\begin{equation}\label{eq:general-tensor-duals}
\begin{aligned}
\theta_r:\Hom_C(N,X)\otimes_B\Hom_B(M,B)
&\longrightarrow\Hom_C(M\otimes_B N,X),\\
\theta_r(g\otimes f)(m\otimes n)&=g(f(m)n),\\
\theta_\ell:\Hom_{B^\op}(N,B)\otimes_B\Hom_{A^\op}(M,Y)
&\longrightarrow\Hom_{A^\op}(M\otimes_B N,Y),\\
\theta_\ell(g\otimes f)(m\otimes n)&=f(mg(n)).
\end{aligned}
\end{equation}
The maps $\theta_r$ and $\theta_\ell$ are right $A$-linear and left $C$-linear, respectively.
\begin{enumerate}
\item The map $\theta_r$ is an isomorphism if $M\in\add B$ as a right $B$-module.
\item The map $\theta_\ell$ is an isomorphism if ${}_BN\in\add_{\mod B^\op}(B)$.
\end{enumerate}
Here $\add B=\add_{\mod B}(B)$, as in Example~\ref{ex:add-category}.
If ${}_CX_C$ and ${}_AY_A$ are bimodules, the respective maps also preserve the additional actions and are morphisms in $C\Mod A$.
\end{proposition}

\begin{proof}
\emph{Well-definedness of $\theta_r$.}
For $f\in\Hom_B(M,B)$, $g\in\Hom_C(N,X)$, and $b\in B$,
\[
g(f(mb)n)=g(f(m)bn),\qquad
(gb)(f(m)n)=g(bf(m)n)=g((bf)(m)n).
\]
These identities respect the defining identifications in both tensor products.
Right $C$-linearity follows from $g(f(m)nc)=g(f(m)n)c$ for $c\in C$.
For $a\in A$,
\[
\theta_r(g\otimes fa)(m\otimes n)
=g(f(am)n)=(\theta_r(g\otimes f)a)(m\otimes n).
\]
If $X$ also has a commuting left $C$-action, replacing $g$ by $cg$ multiplies the value by $c$, so $\theta_r$ preserves that action as well.

\emph{Bijectivity of $\theta_r$.}
Regard $M$ only as a right $B$-module.
For $M=B$, evaluation at $1$ identifies $\Hom_B(B,B)\cong B$, while multiplication identifies $B\otimes_B N\cong N$.
Under these isomorphisms, $\theta_r$ becomes
\[
\Hom_C(N,X)\otimes_B B\longrightarrow\Hom_C(N,X),
\qquad g\otimes b\longmapsto gb,
\]
which is an isomorphism by Lemma~\ref{lem:tensor-properties}(3).
Finite direct sums commute with both constructions by Lemma~\ref{lem:tensor-properties}(4) and \eqref{eq:hom-finite-sums}.
Thus $\theta_r$ is bijective for $M=B^{\oplus r}$, where $r\geq0$ is finite.
If $M\in\add B$, choose $M'$ with $M\oplus M'\cong B^{\oplus r}$.
The comparison for $M\oplus M'$ is the direct sum of the comparisons for $M$ and $M'$, so both summand maps are bijective.

\emph{The left comparison.}
For $f\in\Hom_{A^\op}(M,Y)$ and $g\in\Hom_{B^\op}(N,B)$, well-definedness follows from
\[
f(mb\,g(n))=f(mg(bn)),\qquad
f(m(gb)(n))=f(mg(n)b)=(bf)(mg(n))
\quad(b\in B).
\]
Left $A$-linearity of $f$ makes the resulting functional left $A$-linear.
For $c\in C$,
\[
\theta_\ell(cg\otimes f)(m\otimes n)
=f(mg(nc))=(c\theta_\ell(g\otimes f))(m\otimes n).
\]
If $Y$ also has a commuting right $A$-action, replacing $f$ by $fa$ multiplies the value by $a$, so $\theta_\ell$ preserves that action as well.
For $N=B$ as a left module, evaluation at $1$ identifies $\Hom_{B^\op}(B,B)\cong B$ and turns $\theta_\ell$ into
\[
B\otimes_B\Hom_{A^\op}(M,Y)\longrightarrow\Hom_{A^\op}(M,Y),
\qquad b\otimes f\longmapsto bf.
\]
The same finite-sum and direct-summand argument proves (2).

\emph{Naturality.}
Let $p:M\to M'$ and $q:N\to N'$ be bimodule maps.
For $f\in\Hom_B(M',B)$ and $g\in\Hom_C(N',X)$, applying $\theta_r$ before or after precomposition gives $g(f(p(m))q(n))$ at $m\otimes n$.
For $f\in\Hom_{A^\op}(M',Y)$ and $g\in\Hom_{B^\op}(N',B)$, the two orders give $f(p(m)g(q(n)))$ for $\theta_\ell$.
Thus both comparisons are natural.
\end{proof}
\end{advanced}

Let $S$ be finite-dimensional semisimple, and consider bimodules ${}_SV_S$, ${}_SW_S$ finitely generated on the right and left, respectively.
Corollary~\ref{cor:semisimple-summands} gives $V_S\in\add S$ and ${}_SW\in\add_{\mod S^\op}(S)$.
Thus the semisimple case of tensor duality, Proposition~\ref{prop:tensor-duals}, gives natural bimodule isomorphisms \comm{you always call this tensor-dual isomorphism, but this name is very ambiguous... it is more like anti-distributivity of $S$-dual over tensor-over-$S$.  you can use emphasise on this phrase to make it more easy to remember, if you need to use word, then say ``we refer to this as anti-distributivity of $S$-dual" from now on; then apply changes for any further use of this (whenever you say `tensor-dual isom', it should be clarified)}
\begin{equation}\label{eq:tensor-duals}
\theta_r:W^*\otimes_S V^*\xrightarrow{\sim}(V\otimes_S W)^*,
\qquad
\theta_\ell:{}^*W\otimes_S{}^*V\xrightarrow{\sim}{}^*(V\otimes_S W),
\end{equation}
\begin{equation}\label{eq:tensor-evaluations}
\theta_r(g\otimes f)(v\otimes w)=g(f(v)w),
\qquad
\theta_\ell(g\otimes f)(v\otimes w)=f(vg(w)).
\end{equation}
For $M\in\{V,W\}$, the actions from \eqref{eq:bimodule-dual-actions} are
\begin{equation}\label{eq:tensor-dual-actions}
(sft)(m)=sf(tm),\qquad (sgt)(m)=g(ms)t
\quad\text{for all }s,t\in S,\ f\in M^*,\ g\in{}^*M,\ m\in M.
\end{equation}

\begin{hidden}
\begin{convention}[Tensor duals]
For $S,V,W$ as in \eqref{eq:tensor-duals}, identify the dual of a tensor product with the tensor product of the corresponding duals in reverse order, using the maps in \eqref{eq:tensor-evaluations}.
No basis or symmetrising form is needed.
Comparisons of $M^*$, $ {}^*M$, and $D(M)$ instead use a chosen symmetrising form, as in \eqref{eq:dual-comparison}.
\end{convention}
\end{hidden}

\begin{advanced}
\begin{example}\label{ex:tensor-dual-finite-generation}
Finite generation alone does not suffice in Proposition~\ref{prop:tensor-duals} over a general algebra.
Take $A=C=\Bbbk$, $B=\Bbbk[x]$, and $X=Y=\Bbbk$.
Consider bimodules ${}_\Bbbk M_B$, ${}_BN_\Bbbk$ with $M=N=B/(x)\cong\Bbbk$ and the outer $\Bbbk$-actions given by scalar multiplication.
Every $f\in\Hom_B(M,B)$ satisfies $f(\overline1)x=f(\overline1x)=0$, hence $f=0$.
Consequently
\[
\Hom_\Bbbk(N,\Bbbk)\otimes_B\Hom_B(M,B)=0,\qquad
\Hom_\Bbbk(M\otimes_B N,\Bbbk)\cong\Bbbk.
\]
Both $M_B$ and $ {}_B N$ are generated by $\overline1$, so neither direct-summand hypothesis can be replaced by finite generation alone.
For the left comparison, $\Hom_{B^\op}(N,B)=0$, whereas $\Hom_\Bbbk(M\otimes_B N,\Bbbk)\cong\Bbbk$.
\end{example}

For $S,V,W$ as in \eqref{eq:tensor-duals}, take $A=B=S$, $C=\Bbbk$, $M=V$, $N=W$, and $X=\Bbbk$ in Proposition~\ref{prop:tensor-duals}(1).
Since $\Hom_\Bbbk(W,\Bbbk)=D(W)$, we obtain the natural bimodule isomorphism \cite[\S2.1]{GI20}
\[
\begin{aligned}
\theta_D:D(W)\otimes_S V^*&\xrightarrow{\sim}D(V\otimes_S W),\\
\theta_D(\lambda\otimes f)(v\otimes w)&=\lambda(f(v)w).
\end{aligned}
\]
The actions on $D(W)$ are those in \eqref{eq:bimodule-dual-actions}.
For $s\in S$, naturality with respect to $W\to W$, $w\mapsto ws$, gives left $S$-linearity of $\theta_D$.

Suppose now that $S$ is equipped with a symmetrising form $\tau:S\to\Bbbk$, as in Definition~\ref{def:symmetrising-form}.
The two-sided dual of Definition~\ref{def:en-dual} satisfies
\begin{equation}\label{eq:tensor-two-sided-dual}
\begin{aligned}
\theta_{\en}:W^*\otimes_S V^{*_{\en}}
&\xrightarrow{\sim}(V\otimes_S W)^{*_{\en}},\\
\theta_{\en}(g\otimes F)(v\otimes w)
&=\sum_j g(a_jw)\otimes b_j^\op,
\qquad F(v)=\sum_j a_j\otimes b_j^\op.
\end{aligned}
\end{equation}
We check that the formula defines a map on the two indicated tensor products and preserves both outer actions.
For example, $F(vs)=F(v)(s\otimes1)$ replaces $a_j$ by $a_js$, while
$(Fs)(v)=(1\otimes s^\op)F(v)$ replaces $b_j$ by $b_js$.
These give, respectively, compatibility with $vs\otimes w=v\otimes sw$ and with the right outer action.
Compatibility with $(gs)\otimes F=g\otimes(sF)$ follows from $(gs)(a_jw)=g(sa_jw)$.
The other actions follow from right $S$-linearity of $g$ and $F(tv)=F(v)(1\otimes t^\op)$.

The comparison maps $\gamma_M$ in \eqref{eq:en-dual-comparison} give
\[
\gamma_{V\otimes_S W}\circ\theta_{\en}
=\theta_r\circ(\id_{W^*}\otimes\gamma_V),
\]
since both sides send $g\otimes F$, evaluated at $v\otimes w$, to $\sum_j g(a_jw)\tau(b_j)$.
Thus $\theta_{\en}$ is an isomorphism by Theorem~\ref{prop:dual-comparison}.
Precomposition shows that $\theta_D$ and $\theta_{\en}$ are natural.
No assumption that $S^\en$ is semisimple is used; compare Remark~\ref{rem:enveloping-semisimple}.

For the maps $\alpha_M(f)=\tau\circ f$ in \eqref{eq:dual-comparison},
\[
\alpha_{V\otimes_S W}\circ\theta_r
=\theta_D\circ(\alpha_W\otimes\id_{V^*}),
\]
because both sides send $g\otimes f$ to the functional $v\otimes w\mapsto\tau(g(f(v)w))$.
\end{advanced}


\subsection{Generated algebras}

\begin{definition}
Let $A$ be graded with $S=A_0$.
Multiplication makes $A_1$ an $S$-bimodule and gives
\begin{equation}\label{eq:multiplication-maps}
\mu_i:A_1^{\otimes_S i}\longrightarrow A_i,\qquad
a_1\otimes\cdots\otimes a_i\longmapsto a_1\cdots a_i
\quad(i\geq 1),\qquad \mu_0=\id_S.
\end{equation}
The algebra $A$ is \defn{generated in degree $1$ (over $S$)} if every $\mu_i$ in \eqref{eq:multiplication-maps} is surjective.
\end{definition}

Equivalently, $S$ and $A_1$ generate $A$ as an algebra.

\begin{example}
The algebra $\Bbbk A_n$ of Example~\ref{ex:triangular-grading} is generated in degree $1$ over $S=\bigoplus_i\Bbbk e_{ii}$.
%\[
%e_{ij}=e_{i,i+1}e_{i+1,i+2}\cdots e_{j-1,j}\quad(1\leq i<j\leq n).
%\]
%For the grading $\deg e_{ij}=2(j-i)$, the same algebra has $A_1=0$ and $A_2\neq 0$, so it is not generated in degree $1$.
\end{example}

\begin{example}\label{ex:polynomial-generators}
The polynomial algebra $\Bbbk[x_1,\ldots,x_n]$, with each $\deg x_i=1$ for all $i$, is generated in degree $1$ over $\Bbbk$.
Indeed, every monomial of total degree $d$ is a product of $d$ variables, and these monomials span the degree $d$ component.
\end{example}

\begin{supplementary}
\begin{example}\label{ex:triangular-polynomials}
\emergencystretch=1em
Let
\[
A=\begin{pmatrix}\Bbbk[x]&\Bbbk[x]\\0&\Bbbk[x]\end{pmatrix},
\qquad \deg(x^m e_{ij})=m+j-i
\quad(m\geq0,\ 1\leq i\leq j\leq2).
\]
Then
\[
A_0=\Bbbk e_{11}\oplus\Bbbk e_{22},\qquad
A_d=\Bbbk x^d e_{11}\oplus\Bbbk x^d e_{22}\oplus\Bbbk x^{d-1}e_{12}
\quad(d\geq1).
\]
It is generated in degree $1$ over $A_0$, since
\[
x^d e_{11}=(xe_{11})^d,\qquad
x^d e_{22}=(xe_{22})^d,\qquad
x^{d-1}e_{12}=(xe_{11})^{d-1}e_{12}
\quad(d\geq1).
\]
Here $A_0$ has two simple summands, $A_d\neq0$ for every $d\geq0$, and $A$ is non-commutative: $e_{11}e_{12}=e_{12}$ whereas $e_{12}e_{11}=0$.
\end{example}
\end{supplementary}

\subsection{Tensor algebras and the universal property}

% TODO (ChangeProposal9, Lecture 2 follow-up): Keep the tensor-algebra universal
% property and its explicit formula in the main text; shorten its proof or put
% routine verification in supplementary material. Address the remaining local
% editorial comments when revising Lecture 2, not as part of the Lecture 1 edit.

\begin{definition}\label{def:tensor-algebra}
Let $S$ be a $\Bbbk$-algebra and consider a bimodule ${}_SV_S$.
The \defn{tensor algebra of $V$ over $S$} is
\[
T_S(V):=\bigoplus_{i\in\Z}T_S(V)_i,
\quad\text{where}\quad
T_S(V)_i=
\begin{cases}
V^{\otimes_S i},&i\geq1,\\
S,&i=0,\\
0,&i<0,
\end{cases}.
\]
For $v_1,\ldots,v_{p+q}\in V$, multiplication is determined $\Bbbk$-bilinearly by
\[
(v_1\otimes_S\cdots\otimes_S v_p)(v_{p+1}\otimes_S\cdots\otimes_S v_{p+q})
=v_1\otimes_S\cdots\otimes_S v_{p+q}
\qquad(p,q\geq1),
\]
and by the left and right $S$-actions when a factor has degree $0$.
\end{definition}

We also call $T_S(V)$ the tensor algebra generated by $V$ over $S$.
Every element has finitely many non-zero homogeneous components.
The unit is $1_S$, and associativity follows from Lemma~\ref{lem:tensor-properties}(3).
Write $v_1\cdots v_d$ for $v_1\otimes_S\cdots\otimes_S v_d$ in $T_S(V)$.
In a quotient algebra (including $T_S(V)$ itself) of $T_S(V)$, the notation $v_1\cdots v_d$ denotes the image of $v_1\otimes_S\cdots\otimes_S v_d$.

\begin{example}[Polynomial and exterior algebras]\label{ex:polynomial-exterior}
Let $V,W$ be finite-dimensional vector spaces with bases $\{x_i\mid1\leq i\leq n\}$ and $\{y_j\mid1\leq j\leq m\}$, respectively.
The \defn{symmetric algebra} of $V$ and the \defn{exterior algebra} of $W$ are
\begin{equation}\label{eq:polynomial-exterior}
\begin{aligned}
\SymAlg(V)&=T_\Bbbk(V)/(x_ix_j-x_jx_i\mid1\leq i<j\leq n)
            \cong\Bbbk[x_1,\ldots,x_n],\\
\bigwedge W&=T_\Bbbk(W)/(w^2\mid w\in W) = T_\Bbbk(W)/(y_j^2, y_iy_j+y_jy_i \mid 1\le i<j\le m).
\end{aligned}
\end{equation}
They are generated in degree $1$, with
\[
\SymAlg(V)_0=\Bbbk=(\bigwedge W)_0,\qquad
\SymAlg(V)_1=V,\qquad(\bigwedge W)_1=W.
\]
Their respective bases are
\[
\begin{gathered}
\{x_1^{a_1}\cdots x_n^{a_n}\mid a_1,\ldots,a_n\in\Z,\ a_1,\ldots,a_n\geq0\},\\
\{y_{j_1}\cdots y_{j_d}\mid0\leq d\leq m,\ 1\leq j_1<\cdots<j_d\leq m\}.
\end{gathered}
\]
\end{example}

\begin{supplementary}
The bases in Example~\ref{ex:polynomial-exterior} can be checked directly.
For $\SymAlg(V)$, commuting adjacent generators gives the ordered monomials, which are independent under the map to the polynomial ring.
For $\bigwedge W$, the relations reduce every product of basis elements to a multiple of a strictly increasing product or to zero.
For $1\leq d\leq m$, fix $j_1<\cdots<j_d$ and consider the linear functional $W^{\otimes_\Bbbk d}\to\Bbbk$ sending $w_1\otimes\cdots\otimes w_d$ to the determinant of their coordinates in rows $j_1,\ldots,j_d$ relative to the basis $\{y_j\mid1\leq j\leq m\}$.
The determinant vanishes whenever two adjacent factors agree, hence on the degree $d$ component of $(w^2\mid w\in W)$.
On increasing products it is $1$ on $y_{j_1}\cdots y_{j_d}$ and $0$ on the others, proving independence in every characteristic.
\end{supplementary}

\begin{lemma}\label{lem:idempotent-tensors}
Let $S=\bigoplus_{i=1}^r\Bbbk e_i$ with coordinate idempotents, and consider bimodules ${}_SV_S$, ${}_SW_S$.
There is a bimodule isomorphism
\begin{equation}\label{eq:idempotent-tensor-decomposition}
V\otimes_S W\cong\bigoplus_{i,j,k=1}^r e_iVe_j\otimes_\Bbbk e_jWe_k.
\end{equation}
Consequently, after choosing a basis of each $e_iVe_j$, a basis of $V^{\otimes_S d}$ for $d\geq1$ consists of products of $d$ chosen basis elements whose adjacent idempotents match.
\end{lemma}
\begin{supplementary}
\begin{proof}
The orthogonal idempotents sum to $1$, so
\begin{equation}\label{eq:composable-tensors}
V=\bigoplus_{i,j=1}^r e_iVe_j,\qquad
(e_iVe_j)\otimes_S(e_pWe_k)=0
\quad(1\leq i,j,p,k\leq r,\ j\neq p).
\end{equation}
Indeed, $v\otimes_S w=ve_j\otimes_S w=v\otimes_S e_jw=0$ for $v\in e_iVe_j$, $w\in e_pWe_k$, and $j\neq p$.
The map from right to left in \eqref{eq:idempotent-tensor-decomposition} sends $v\otimes_\Bbbk w$ to $v\otimes_S w$ on each summand.
Its inverse is
\[
V\otimes_S W\longrightarrow\bigoplus_{i,j,k=1}^r e_iVe_j\otimes_\Bbbk e_jWe_k,
\qquad
v\otimes_S w\longmapsto\sum_{i,j,k=1}^r e_ive_j\otimes_\Bbbk e_jwe_k.
\]
Checking that these two assignments induce well-defined, mutually inverse maps is left to the reader.
Iterating \eqref{eq:idempotent-tensor-decomposition} and using Definition~\ref{def:vector-tensor} gives the stated basis.
\end{proof}
\end{supplementary}


\begin{example}\label{ex:tensor-A3}
We apply Lemma~\ref{lem:idempotent-tensors} to realise $\Bbbk A_3$ as a tensor algebra.
Let $S=\Bbbk^{\oplus3}$ with coordinate idempotents $e_1,e_2,e_3$, and take the bimodule ${}_SV_S$ with $V=\Bbbk a\oplus\Bbbk b$, $e_1ae_2=a$, and $e_2be_3=b$; all other idempotent actions are zero.
By Lemma~\ref{lem:idempotent-tensors},
\[
V\otimes_S V=\Bbbk(a\otimes b),\qquad
V^{\otimes_S3}=0,\qquad
T_S(V)=S\oplus\Bbbk a\oplus\Bbbk b\oplus\Bbbk ab.
\]
Here $ab=a\otimes b$.
The linear map $T_S(V)\to\Bbbk A_3$ given by $e_i\mapsto e_{ii}$, $a\mapsto e_{12}$, $b\mapsto e_{23}$, and $ab\mapsto e_{13}$ preserves multiplication and sends a basis to a basis, hence is an algebra isomorphism.
\end{example}

\comm{at some point before universal property prop, it should be clarified that tensor algebra is a graded algebra.  the current definition just says something of degree 0, without saying degree 0 in where.  If necessary, we will refer to the canonical grading structure $\deg S=0, \deg V=1$ as tensor grading.}

\begin{proposition}[Universal property of the tensor algebra]\label{prop:tensor-universal}
Let $S$ be a $\Bbbk$-algebra, consider a bimodule ${}_SV_S$, and let $\iota_S:S\to T_S(V)$ and $\iota_V:V\to T_S(V)$ be the canonical inclusions.
Given a unital $\Bbbk$-algebra homomorphism $\varphi:S\to B$ and a bimodule map $f:{}_SV_S\to{}_SB_S$, where $sbt=\varphi(s)b\varphi(t)$ on $B$, there is a unique $\Bbbk$-algebra homomorphism $\widetilde f:T_S(V)\to B$ making the diagram commute:
\begin{equation}\label{eq:tensor-universal}
\xymatrix@C=42pt@R=28pt{
V \ar@{^{(}->}[r]^{\iota_V} \ar[dr]_{\forall\,\text{bimodule map }f\phantom{...}}
& T_S(V) \ar@{-->}[d]^{\exists!\,\widetilde f}
& S \ar@{_{(}->}[l]_{\iota_S} \ar[dl]^{\phantom{...}\forall\,\text{algebra map }\varphi}\\
& B. &
}
\end{equation}
\end{proposition}
\begin{supplementary}
\begin{proof}
Define $\widetilde f$ on $S$ and each tensor power by
\begin{equation}\label{eq:tensor-extension}
\widetilde f(s)=\varphi(s),\qquad
\widetilde f(v_1\otimes_S\cdots\otimes_S v_i)=f(v_1)\cdots f(v_i)
\quad(s\in S,\ i\geq1).
\end{equation}
The identity $f(vs)f(w)=f(v)\varphi(s)f(w)=f(v)f(sw)$ shows that replacing adjacent factors $vs\otimes w$ by $v\otimes sw$ does not change the value in \eqref{eq:tensor-extension}.
The formulas therefore induce linear maps on the tensor powers, which extend to their direct sum.
Multiplication is preserved by \eqref{eq:tensor-extension}, including the degree $0$ factors because $f(sv)=\varphi(s)f(v)$ and $f(vs)=f(v)\varphi(s)$.
Also, $\widetilde f(1_S)=1_B$.
Every algebra homomorphism making \eqref{eq:tensor-universal} commute must have the values in \eqref{eq:tensor-extension}, proving uniqueness.
\end{proof}

If $B$ is graded, $\varphi(S)\subset B_0$, and $f(V)\subset B_1$, then \eqref{eq:tensor-extension} gives $\widetilde f(V^{\otimes_S i})\subset B_i$ for every $i\geq0$.
Thus $\widetilde f$ is a graded algebra homomorphism.
\end{supplementary}


\begin{corollary}\label{cor:tensor-presentation}
If $A$ is generated in degree $1$ over $S=A_0$, then the multiplication map defines a surjective graded algebra homomorphism
\[
\mu:T_S(A_1)\onto A,\qquad \mu|_S=\id_S, \qquad 
a_1\otimes_S\cdots\otimes_S a_i\longmapsto a_1\cdots a_i
\quad \forall i\geq1.
\]
Its degree $i$ component is $\mu_i$ from \eqref{eq:multiplication-maps}.
\end{corollary}

\begin{supplementary}
\begin{proof}
In Proposition~\ref{prop:tensor-universal}, take $V=A_1$, $B=A$, and the inclusions
\[
\varphi:S\hookrightarrow A,\quad s\longmapsto s,
\qquad f:A_1\hookrightarrow A,\quad a\longmapsto a.
\]
Then \eqref{eq:tensor-extension} gives $\widetilde f=\mu$.
The map $\mu$ preserves degrees, and each component $\mu_i$ is surjective by the assumption that $A$ is generated in degree $1$.
\end{proof}
\end{supplementary}

\begin{definition}\label{def:relations}
For the map $\mu$ in Corollary~\ref{cor:tensor-presentation}, the homogeneous two-sided ideal $I:=\Ker\mu$ is the \defn{ideal of relations}.
\end{definition}

\comm{this paragraph, up to `Homegentiy is essential...', should be inside the corollary??}
Note that $I_0=I_1=0$ since $\mu_0=\id_S$ and $\mu_1=\id_{V=A_1}$.
Since $\mu$ preserves degrees, its kernel is homogeneous, and
\[
T_S(V)/I\xrightarrow{\sim}A,\qquad z+I\longmapsto\mu(z)
\]
is a graded algebra isomorphism.
The bimodule $V$ specifies the generators and their left and right $S$-actions; $I$ specifies the relations between their products.
Homogeneity is essential for the quotient grading: the relation $x-x^2$ in Example~\ref{ex:nonhomogeneous-ideal} does not satisfy this condition.

\subsection{Quadratic algebras}

We now consider algebras whose ideal of relations is generated in degree $2$.

\begin{definition}\label{def:quadratic} 
A graded algebra $A$ is \defn{quadratic} if there exist a finite-dimensional semisimple $\Bbbk$-algebra $S$, a finite-dimensional $S$-bimodule $V$, and an $S$-sub-bimodule $R\subset V\otimes_S V$ such that $A\cong T_S(V)/(R)$ as graded algebras, where $(R)$ is the two-sided ideal of $T_S(V)$ generated by $R$.
\end{definition}

An isomorphism $A\cong T_S(V)/(R)$ as in Definition~\ref{def:quadratic} is called a \defn{quadratic presentation} of $A$.

\begin{proposition}\label{prop:quadratic-basic}
A quadratic algebra $A=T_S(V)/(R)$ has $A_0=S$, $A_1=V$, and $A_2=(V\otimes_S V)/R$, is locally finite-dimensional, and is generated in degree $1$ over $S$.
\end{proposition}

\begin{proof}
The ideal $(R)$ has no components in degrees $0$ and $1$.
Its degree $2$ component is exactly $R$.
Each $A_i$ is a quotient of the finite-dimensional space $V^{\otimes_S i}$, and multiplication maps $V^{\otimes_S i}$ onto $A_i$.
\end{proof}

Thus a quadratic presentation recovers $S\cong A_0$, $V\cong A_1$, and $R\cong\Ker(A_1\otimes_{A_0}A_1\to A_2)$.
In particular, the kernel of the canonical map $T_{A_0}(A_1)\onto A$ is generated by its degree $2$ component.
Quadraticity depends on the grading: for $n\geq2$, the grading $\deg e_{ij}=2(j-i)$ on $\Bbbk A_n$ is not quadratic, since it has no degree $1$ generators.

\begin{example}\label{ex:tensor-An}
Let $S=\bigoplus_{i=1}^n\Bbbk e_{ii}$ and $V=\bigoplus_{i=1}^{n-1}\Bbbk e_{i,i+1}$ inside $\Bbbk A_n$.
Then we have $T_S(V)\cong\Bbbk A_n$ given by\[
e_{i,i+1}\otimes_S e_{i+1,i+2}\otimes_S\cdots\otimes_S e_{j-1,j}
\longmapsto e_{ij}
\quad(1\leq i<j\leq n).
\]
\end{example}

\begin{example}\label{ex:zero-composition}
For $S,V,a,b$ from Example~\ref{ex:tensor-A3}, the quotient
\[
T_S(V)/(ab)\cong\Bbbk A_3/(e_{13})
\]
is quadratic with $R=\Bbbk(a\otimes b)$.
Its basis is $e_1,e_2,e_3,a,b$, in degrees $0,0,0,1,1$, as in Example~\ref{ex:quadratic-matrix-quotient}.
\end{example}

\begin{example}\label{ex:cubic-quotient}
Let $B=\Bbbk A_4$ and $A=B/(e_{14})$, with the grading of Example~\ref{ex:triangular-grading}.
The ideal $I=(e_{14})=\Bbbk e_{14}$ has $I_2=0$ but $I_3\neq 0$.
Since $e_{14}=e_{12}e_{23}e_{34}$, the algebra $A$ is generated in degree $1$ with a relation of degree $3$ and no relations of degree $2$.
By Proposition~\ref{prop:quadratic-basic}, any quadratic presentation would have the same generators and zero quadratic relation space, hence would give $\Bbbk A_4$ itself. Thus $A$ is not quadratic.
\end{example}

\begin{example}
The symmetric and exterior algebras of Example~\ref{ex:polynomial-exterior} are quadratic: the defining relations in \eqref{eq:polynomial-exterior} belong to $V\otimes_\Bbbk V$ and $W\otimes_\Bbbk W$, respectively.
\end{example}

\subsection{Hilbert series}

\comm{TODO: throughout, use $n$ instead of $r$ for number of simples?  Avoid using $P_i$, and just stick with $e_iA$.}

Let $A$ be a locally finite-dimensional graded algebra with $S=A_0$ finite-dimensional semisimple.
Then each degree piece of a graded right $A$-module is semisimple as $A$-module.  We call the simple direct summands in a direct sum decomposition of $M_d$  \defn{composition factors}.

\begin{notation}
By Theorem~\ref{thm:artin-wedderburn}, choose primitive idempotents $e_1,\ldots,e_r\in S$.  Then  
\[L_j:=e_jS\cong e_j A/e_jA_+\]
gives the full list of isomorphism class of (graded) simple right $S$-modules.
\phantomsection\label{not:hilbert-simples-summands}
\end{notation}

\begin{definition}\label{def:dimension-vector}\label{def:hilbert-series}
For a locally finite-dimensional graded right $A$-module $M$ with primitive idempotent $e_1, \ldots e_r$, write $[M_d:L_j]$ for the multiplicity of the composition factor $L_j$ in $M_d$.
Its degree $d$ \defn{dimension vector} is the row vector\footnote{In most representation theory literature, it is common to regard dimension vector as a column vector.  But in the context of Hilbert series, we follow Beilinson--Ginzburg--Soergel convention to use row vector here.}
\[
\dimvec M_d:=\bigl([M_d:L_j]\bigr)_{1\leq j\leq r}.
\]
The module $M$ is \defn{bounded below} if $M_d=0$ for all sufficiently negative $d$.
For such $M$, its \defn{Hilbert series} is
\[
h_M(t):=\sum_{d\in\Z}(\dimvec M_d)t^d
=\sum_{d\in\Z}\bigl([M_d:L_j]\bigr)_{1\leq j\leq r}t^d.
\]
The \defn{Hilbert series} of the algebra $A$ collects the right-module series $h_{e_iA}(t)$ as rows:
\[
h_A(t)_{ij}:=\bigl(h_{e_iA}(t)\bigr)_j
=\sum_{d\geq0}[e_iA_d:L_j]t^d
\quad\text{for all }1\leq i,j\leq r.
\]
\end{definition}

We think of the Hilbert series as ``graded dimension matrix" with one indecomposable per row. 
The entries of the Hilbert series are formal Laurent series; for algebras they are formal power series.
Finite support gives Laurent polynomials or polynomials, respectively.
Since $e_iA_0=L_i$, we have $h_A(0)=I_r$, the $r\times r$ identity array.
For a finite-dimensional right $S$-module $N$, use the same multiplicity notation and write $\dimvec N=([N:L_j])_{1\leq j\leq r}$.

The multiplicities can be computed using Hom spaces:
\begin{equation}\label{eq:composition-Hom}
[M_d:L_j]=\dim_{\End_A(L_j)}\Hom_A(e_jA,M_d)
\quad\text{for all }d\in\Z,\ 1\leq j\leq r.
\end{equation}
Indeed, $M_dA_+=0$, so every map $e_jA\to M_d$ factors uniquely through $e_jA/e_jA_+\cong L_j$.
This identifies $\Hom_A(e_jA,M_d)$ with $\Hom_A(L_j,M_d)$; precomposition gives its right vector-space structure over the division algebra $\End_A(L_j)$.
By the Schur's Lemma, A non-zero map between simple modules is an isomorphism, so each copy of $L_j$ contributes one copy of $\End_A(L_j)$, and the other simples contribute zero.
This also shows that the multiplicities are independent of the chosen decomposition.

\begin{example}\label{ex:dimension-vector-coordinate}
If $S=\bigoplus_{j=1}^r\Bbbk e_j$ with coordinate idempotents, then $\End_A(L_j)=\Bbbk$.
The Yoneda lemma~\ref{lem:idempotent-hom} says that 
\[
[M_d:L_j]=\dim_\Bbbk(M_de_j),
\qquad \dimvec M_d=(\dim_\Bbbk M_de_j)_{1\leq j\leq r}.
\]
\end{example}

\begin{lemma}\label{lem:hilbert-operations}
For locally finite-dimensional graded right $A$-modules bounded below:
\begin{enumerate}
\item A graded isomorphism $M\cong N$ implies $h_M(t)=h_N(t)$.
\item A short exact sequence $0\to L\to M\to N\to0$ gives $h_M(t)=h_L(t)+h_N(t)$.
\item For every $a\in\Z$, $h_{M(a)}(t)=t^{-a}h_M(t)$.
\end{enumerate}
\end{lemma}

\begin{proof}
A graded isomorphism preserves each degree and its simple-module multiplicities.
In each degree, a graded short exact sequence gives $0\to L_d\to M_d\to N_d\to0$, so $[M_d:L_j]=[L_d:L_j]+[N_d:L_j]$ for all $1\leq j\leq r$.
Finally, $M(a)_d=M_{d+a}$ gives
\[
h_{M(a)}(t)=\sum_d(\dimvec M_{d+a})t^d
=t^{-a}\sum_j(\dimvec M_j)t^j.\qedhere
\]
\end{proof}

\comm{From here}

\begin{definition}\label{def:total-tensor-grading}
For graded right and left $S$-modules $M,N$, their tensor product has \defn{total grading}
\[
(M\otimes_S N)_d=\bigoplus_{p+q=d}M_p\otimes_S N_q.
\]
The same definition applies to bimodules, and to vector spaces by taking $S=\Bbbk$.
\end{definition}

The decomposition follows by distributing tensor products over direct sums, as in Lemma~\ref{lem:tensor-properties}(4).
When $M,N$ are locally finite-dimensional and bounded below, only finitely many pairs $(p,q)$ contribute in a fixed degree, so the tensor product is locally finite-dimensional and bounded below.

For a finite-dimensional graded bimodule $ {}_SV_S$, use $V^{\grstar}$ and $ {}^{\grstar}V$ in Proposition~\ref{prop:S-evaluation-coevaluation}.
With the total grading, evaluation and coevaluation preserve degree: evaluation pairs degrees $i$ and $-i$, and $\id_V$ has degree $0$.

To multiply two arrays $C=(c_{ij})$ and $D=(d_{ij})$, or a row vector $m=(m_i)$ and an array $C$, use
\[
(CD)_{i\ell}=\sum_{j=1}^r c_{ij}d_{j\ell},
\qquad (mC)_j=\sum_{i=1}^r m_i c_{ij}.
\]
Thus multiplication sums over the matching middle index, rather than multiplying individual entries independently.
For series, use the same products on coefficients and collect terms with the same total power of $t$.

\begin{lemma}\label{lem:dimension-tensor}\label{lem:hilbert-tensor}
Let $N$ be a finite-dimensional right $S$-module and ${}_SU_S$ a finite-dimensional bimodule.
Then
\[
[N\otimes_S U:L_j]
=\sum_{i=1}^r[N:L_i][e_iU:L_j]
\quad\text{for all }1\leq j\leq r.
\]
For a locally finite-dimensional graded right $A$-module $M$ bounded below,
\[
h_{M\otimes_S A}(t)=h_M(t)h_A(t).
\]
\end{lemma}

\begin{proof}
Decompose $N\cong\bigoplus_{i=1}^r(e_iS)^{\oplus[N:L_i]}$.
Multiplication $e_iS\otimes_S U\to e_iU$, $e_is\otimes u\mapsto e_isu$, is an isomorphism with inverse $u\mapsto e_i\otimes u$.
Distribution over direct sums therefore gives
\[
N\otimes_S U\cong\bigoplus_{i=1}^r(e_iU)^{\oplus[N:L_i]}.
\]
Counting copies of $L_j$ proves the first formula.
Apply the first formula to $M_p\otimes_S A_q$ and sum over all $p,q\in\Z$ with $p+q=d$.
Since $e_iA_q=(P_i)_q$, the resulting coefficient is the coefficient of $t^d$ in $h_M(t)h_A(t)$.
\end{proof}
\comm{up to here.  make as supplementary}

\begin{example}[Hilbert series of a tensor algebra]\label{ex:hilbert-tensor-algebra}
Consider the tensor algebra $A=T_S(V)$ and put 
\[C:=([e_iA_1:L_j])_{1\leq i,j\leq r}=([e_iV:L_j])_{1\leq i,j\leq r}.\]
Since $e_iA_0=L_i$ and $e_iA_{d+1}\cong(e_iA_d)\otimes_S A_1$, Lemma~\ref{lem:dimension-tensor} gives, by induction,
\[
[e_iA_d:L_j]=(C^d)_{ij},
\qquad C^0=I_r
\quad\text{for all }d\geq0,\ 1\leq i,j\leq r.
\]
Consequently,
\begin{equation}\label{eq:hilbert-tensor-algebra}
\boxed{\displaystyle
h_{T_S(V)}(t)=\sum_{d\geq0}t^dC^d=(I_r-tC)^{-1}.}
\end{equation}
Indeed, multiplying the series on either side by $I_r-tC$ gives constant coefficient $I_r$; the coefficient of $t^d$ for $d\geq1$ is $C^d-C^d=0$.
The inverse is therefore an identity of formal power series with the coefficient multiplication defined above.

When $S=\Mat_n(\Bbbk)$ and $\dim_\Bbbk V=m$, the formula specialises to $h_{T_\Bbbk(V)}(t)=1/(1-mt)$.
\end{example}

\comm{make an exercise to show that if $A$ is generated in degree 1, then $h_A=h_{T_S(V)}$ yields isom $A\cong T_S(V)$, and another part of the exercise show an example when this is false if there is generator in degree 2.  put this exercise into homework 1 Q4 after part (1) before the current (2)}

\begin{example}[Upper triangular algebras and a quotient]\label{ex:hilbert-triangular}
For $A=\Bbbk A_n$, order the coordinate idempotents as $e_{11},\ldots,e_{nn}$.
Let $C$ have $C_{i,i+1}=1$ for $1\leq i<n$ and all other entries zero.
The tensor presentation in Example~\ref{ex:tensor-An} and \eqref{eq:hilbert-tensor-algebra} give
\[
h_A(t)=I_n+tC+\cdots+t^{n-1}C^{n-1},
\qquad h_A(t)_{ij}=\begin{cases}t^{j-i},&i\leq j,\\0,&i>j.\end{cases}
\]
For $n=3$, killing $e_{13}$ removes the degree $2$ contribution in entry $(1,3)$:
\[
h_{\Bbbk A_3}(t)=\begin{pmatrix}1&t&t^2\\0&1&t\\0&0&1\end{pmatrix},
\qquad
h_{\Bbbk A_3/(e_{13})}(t)=\begin{pmatrix}1&t&0\\0&1&t\\0&0&1\end{pmatrix}.
\]
The quotient removes $\Bbbk e_{13}$ from degree $2$, and hence removes $t^2$ from entry $(1,3)$.
\end{example}

\begin{example}\label{ex:hilbert-polynomial-exterior}
For the symmetric algebra $\SymAlg(V)$ and exterior algebra $\bigwedge W$ in Example~\ref{ex:polynomial-exterior} with $n=\dim_\Bbbk V, m=\dim_\Bbbk W$. 
Using the description of their basis we get
\[
h_{\SymAlg(V)}(t)=\sum_{d\ge 0}\binom{d+n-1}{n-1}t^d=\frac{1}{(1-t)^n},\footnote{differential $1/(1-t)=\sum_{d\ge 0} t^d$ on both side $n$ times to see this}\qquad
h_{\bigwedge W}(t)=(1+t)^m.
\]
\end{example}

\begin{supplementary}
\begin{lemma}\label{lem:hilbert-field-tensor}
For locally finite-dimensional graded vector spaces $U,W$ bounded below,
\[
h_{U\otimes_\Bbbk W}(t)=h_U(t)h_W(t).
\]
Let $A,B$ be locally finite-dimensional graded algebras with $A_0=\Bbbk^{\oplus r}$ and $B_0=\Bbbk^{\oplus s}$.
Write $e_1,\ldots,e_r$ and $f_1,\ldots,f_s$ for the coordinate idempotents.
Give $A\otimes_\Bbbk B$ the total grading and multiplication
$(a\otimes b)(a'\otimes b')=aa'\otimes bb'$.
Then
\[
h_{A\otimes_\Bbbk B}(t)_{(i,p),(j,q)}
=h_A(t)_{ij}h_B(t)_{pq}
\quad(1\leq i,j\leq r,\ 1\leq p,q\leq s).
\]
\end{lemma}
\end{supplementary}

\begin{proof}
For vector spaces, take dimensions in each summand of the total grading.
For algebras,
\[
(e_i\otimes f_p)(A_a\otimes_\Bbbk B_b)(e_j\otimes f_q)
=(e_iA_ae_j)\otimes_\Bbbk(f_pB_bf_q).
\]
Take vector space dimensions and sum over $a+b=d$ to obtain the coefficient of $t^d$.
\end{proof}

\begin{advanced}[suppressfirstparskip=true,innertopmargin=0pt]
\subsection{The degree 2 relations of a quotient}

Let $S$ be finite-dimensional semisimple, $V$ a finite-dimensional $S$-bimodule, and $I\subset T_S(V)$ a homogeneous two-sided ideal with $I_0=I_1=0$.
Set $A=T_S(V)/I$.
In degree $2$,
\[
A_2\cong(V\otimes_S V)/I_2,
\quad\text{where}\quad
I_2=I\cap(V\otimes_S V).
\]

\begin{definition}\label{def:quadratic-part}
The \defn{quadratic part} of $A$ is $qA:=T_S(V)/(I_2)$.
\end{definition}

Since $(I_2)\subset I$, there is a surjective graded algebra homomorphism
\[
\pi_A:qA\longrightarrow A,\qquad z+(I_2)\longmapsto z+I,
\qquad \Ker\pi_A=I/(I_2).
\]
Consequently,
\[
\pi_A\text{ is an isomorphism}
\quad\Longleftrightarrow\quad I=(I_2)
\quad\Longleftrightarrow\quad A\text{ is quadratic}.
\]
For the last equivalence, $(I_2)$ gives a quadratic presentation, while Proposition~\ref{prop:quadratic-basic} identifies the generators and degree $2$ relations of any quadratic presentation with $A_1$ and $I_2$.

\begin{example}
For $A=\Bbbk A_4/(e_{14})$ from Example~\ref{ex:cubic-quotient}, we have $I_2=0$, so
\[
qA=\Bbbk A_4,\qquad
\pi_A:\Bbbk A_4\onto\Bbbk A_4/(e_{14}),\qquad
\Ker\pi_A=\Bbbk e_{14}.
\]
Thus $qA$ does not impose the cubic relation $e_{12}e_{23}e_{34}=0$.
\end{example}
\end{advanced}

\subsection{Orthogonal complements and quadratic duality}

Let $S$ be finite-dimensional semisimple.

\begin{definition}\label{def:orthogonals}
For an $S$-sub-bimodule $U$ of a finite-dimensional $S$-bimodule $M$, its \defn{right orthogonal complement} and \defn{left orthogonal complement} are
\[
\begin{aligned}
U^\perp&:=\{f\in M^*\mid f(u)=0\text{ for all }u\in U\},\\
{}^\perp U&:=\{g\in{}^*M\mid g(u)=0\text{ for all }u\in U\}.
\end{aligned}
\]
\end{definition}

Both orthogonal complements are $S$-sub-bimodules.
Indeed, for $s,t\in S$ and $u\in U$, \eqref{eq:tensor-dual-actions} gives
\[
(sft)(u)=sf(tu)=0\quad(f\in U^\perp),\qquad
(sgt)(u)=g(us)t=0\quad(g\in{}^\perp U),
\]
since $tu,us\in U$.
The orthogonal complements are defined without choosing a symmetrising form, but it is possible to use realise (through the natural isomorphism associated to a chosen symmeterising form) the orthogonal complements as a sub-bimodule of $\bbD(M)$.

\begin{advanced}
If $S$ is equipped with a symmetrising form $\tau:S\to\Bbbk$, the comparison maps $\alpha_M,\beta_M$ of \eqref{eq:dual-comparison} satisfy
\begin{equation}\label{eq:orthogonal-comparison}
\alpha_M(U^\perp)=\beta_M({}^\perp U)
=\{\lambda\in D(M)\mid\lambda(u)=0\text{ for all }u\in U\}.
\end{equation}
The forward inclusions follow by composing with $\tau$.
Conversely, $\alpha_M(f)(U)=0$ gives
\[
\tau(f(u)s)=\tau(f(us))=0
\quad\text{for all }s\in S,\ u\in U.
\]
Non-degeneracy of $\tau$ forces $f(u)=0$ for every $u\in U$.
For the left dual, $\tau(sg(u))=\tau(g(su))=0$ gives $g(u)=0$ instead.
Bijectivity of $\alpha_M,\beta_M$ in Theorem~\ref{prop:dual-comparison} proves \eqref{eq:orthogonal-comparison}.
\end{advanced}

\begin{example}\label{ex:symmetric-orthogonals}
We compute the orthogonal complements of the relation spaces of $\SymAlg(V)$ and $\bigwedge W$ for $\dim_{\Bbbk}V=2=\dim_{\Bbbk}W$.
Take $V=\Bbbk x_1\oplus\Bbbk x_2$ and $W=\Bbbk y_1\oplus\Bbbk y_2$ in Example~\ref{ex:polynomial-exterior}.
The relation spaces in $V\otimes_\Bbbk V$ and $W\otimes_\Bbbk W$, respectively, are
\[
\begin{aligned}
R_{\SymAlg(V)}&=\Bbbk(x_1\otimes x_2-x_2\otimes x_1),\\
R_{\bigwedge W}&=\Bbbk(y_1\otimes y_1)\oplus\Bbbk(y_2\otimes y_2)\oplus\Bbbk(y_1\otimes y_2+y_2\otimes y_1).
\end{aligned}
\]
Let $x_1^*,x_2^*$ and $y_1^*,y_2^*$ be the dual bases of $D(V)$ and $D(W)$, respectively.
Under the tensor-dual identification \eqref{eq:tensor-evaluations},
\[
\begin{aligned}
R_{\SymAlg(V)}^\perp&=\operatorname{span}_\Bbbk\{x_1^*\otimes x_1^*,\,x_2^*\otimes x_2^*,\,x_1^*\otimes x_2^*+x_2^*\otimes x_1^*\},\\
R_{\bigwedge W}^\perp&=\Bbbk(y_1^*\otimes y_2^*-y_2^*\otimes y_1^*).
\end{aligned}
\]
The first space consists of the exterior relations on $D(V)$; the second consists of the symmetric relations on $D(W)$.
\end{example}

\Needspace{12\baselineskip}
\subsection{Exercises}

\begin{exercise}[Interaction of tensor product with $S$-dual]\label{ex:tensor-duality-proof}
Let $S$ be finite-dimensional semisimple, and consider finite-dimensional bimodules ${}_SV_S$, ${}_SW_S$.
Prove the right $S$-dual isomorphism \eqref{eq:tensor-duals}, the semisimple case of Proposition~\ref{prop:tensor-duals}.
\begin{enumerate}
\item Check that the formula for $\theta_r$ in \eqref{eq:tensor-evaluations} is bilinear in each pair of variables, is unchanged by
\[
vs\otimes w\longleftrightarrow v\otimes sw,
\qquad gs\otimes f\longleftrightarrow g\otimes sf
\quad(s\in S),
\]
and gives a right $S$-linear map $V\otimes_S W\to S$ for each $g\otimes f$.
Verify $\theta_r(sg\otimes ft)=s\theta_r(g\otimes f)t$ for $s,t\in S$.
\item Prove that $\theta_r$ is bijective.
\item For bimodule maps $p:V\to V'$ and $q:W\to W'$, prove
\[
(p\otimes q)^*\circ\theta_r^{V',W'}
=\theta_r^{V,W}\circ(q^*\otimes p^*).
\]
Here the superscripts specify the pair of bimodules, and $p^*,q^*$, and $(p\otimes q)^*$ are given by precomposition with $p,q$, and $p\otimes q$, respectively.
\end{enumerate}
\emph{Hint for (2).} First take $V=S$, then use Corollary~\ref{cor:semisimple-summands}.

\emph{Remark.} The analogous assertions hold for left $S$-duals, with left and right interchanged.
\end{exercise}

\begin{hidden}
\begin{proof}
\emph{Well-definedness.}
For the right duals, right $S$-linearity gives
\[
g(f(vs)w)=g(f(v)sw),\qquad
(gs)(f(v)w)=g(sf(v)w)=g((sf)(v)w).
\]
The first identity respects the relation in $V\otimes_S W$, and the second respects the relation in $W^*\otimes_S V^*$.
The resulting functional is right $S$-linear because $g$ is.
For the left duals, the corresponding identities are
\[
f(vs\,g(w))=f(vg(sw)),\qquad
f(v(gs)(w))=f(vg(w)s)=(sf)(vg(w)).
\]
Left $S$-linearity of $f$ makes the resulting functional left $S$-linear.
The outer actions in \eqref{eq:tensor-dual-actions} give
\[
\begin{aligned}
\theta_r(sg\otimes ft)(v\otimes w)&=s\,g(f(tv)w),\\
\theta_\ell(sg\otimes ft)(v\otimes w)&=f(vg(ws))t,
\end{aligned}
\]
which are the required bimodule identities.

\emph{Bijectivity.}
To check $\theta_r$, forget the left action on $V$.
For $V=S$, identify $S^*\cong S$ by $f\mapsto f(1)$; then $\theta_r$ is $W^*\otimes_S S\to W^*$, $g\otimes s\mapsto gs$, an isomorphism by Lemma~\ref{lem:tensor-properties}(3).
The formula commutes with maps of right $S$-modules in $V$, so it is an isomorphism for finite direct sums and direct summands of $S$.
Corollary~\ref{cor:semisimple-summands} applies to $V$.
For $\theta_\ell$, use $W=S$ as a left module, where the map is $S\otimes_S{}^*V\to{}^*V$, $s\otimes f\mapsto sf$, and apply the same argument to $W$.

\emph{Naturality.}
For bimodule maps $p:V\to V'$ and $q:W\to W'$, the square
\[
\xymatrix@C=48pt@R=22pt{
(W')^*\otimes_S(V')^* \ar[r]^{\theta_r} \ar[d]_{q^*\otimes p^*}
& (V'\otimes_S W')^* \ar[d]^{(p\otimes q)^*}\\
W^*\otimes_S V^* \ar[r]_{\theta_r}
& (V\otimes_S W)^*
}
\]
commutes: both routes send $g\otimes f$, evaluated at $v\otimes w$, to $g(f(p(v))q(w))$.
For $\theta_\ell$, the corresponding two routes give $f(p(v)g(q(w)))$.
\end{proof}

\end{hidden}

\begin{exercise}
Let $A=\Bbbk A_4/(e_{13},e_{24})$ with the grading of Example~\ref{ex:triangular-grading}.
Find a homogeneous basis, compute $h_A(t)$, and write $A=T_S(V)/(R)$ with explicit $S,V,R$.
\emph{Hint.} Use Lemma~\ref{lem:idempotent-tensors}; the ideal contains $e_{14}$ as well as $e_{13},e_{24}$.
\end{exercise}

\begin{exercise}
For the bimodule $V$ in Example~\ref{ex:two-vertex-relations}, find a basis of each $T_S(V)_d$ and show that
\[
h_{T_S(V)}(t)=\frac{1}{1-t^2}\begin{pmatrix}1&t\\t&1\end{pmatrix}.
\]
Compare with the quotient $T_S(V)/(ab)$.
\emph{Hint.} Count alternating words separately in each degree, including degree $0$.
\end{exercise}

\begin{exercise}
For the algebras in Example~\ref{ex:polynomial-exterior}, prove
\[
\dim_\Bbbk\SymAlg(V)_2=\frac{n(n+1)}2,\qquad
\dim_\Bbbk(\bigwedge W)_2=\frac{m(m-1)}2
\quad\text{where }n=\dim_\Bbbk V,\ m=\dim_\Bbbk W.
\]
\emph{Hint.} Count the ordered pairs allowed in the two bases of Example~\ref{ex:polynomial-exterior}.
\end{exercise}

\pagebreak


% !TEX root = Lecture_Koszul.tex
\section{The quadratic dual}\label{sec:quadratic-dual}

\subsection{Generators and orthogonal relations}

Fix a quadratic presentation $A=T_S(V)/(R)$ as in Definition~\ref{def:quadratic}: $S$ is finite-dimensional semisimple, $V$ is a finite-dimensional $S$-bimodule, and $R\subset V\otimes_S V$ is an $S$-sub-bimodule.
We will replace the generators by an $S$-dual and the relations by the $S$-linear maps that vanish on $R$.
Use $V^*=\Hom_S(V,S)$ and ${}^*V=\Hom_{S^\op}(V,S)$ from Definition~\ref{def:S-duals}, with the bimodule actions in \eqref{eq:tensor-dual-actions}.
The tensor isomorphisms \eqref{eq:tensor-duals} have values
\begin{equation}\label{eq:quadratic-pairings}
\begin{aligned}
\theta_r(g\otimes f)(v\otimes w)&=g(f(v)w),
& f,g&\in V^*,\\
\theta_\ell(g\otimes f)(v\otimes w)&=f(vg(w)),
& f,g&\in{}^*V.
\end{aligned}
\end{equation}
Here $v,w\in V$. Both isomorphisms reverse the order of the factors.

\begin{convention}
We identify $V^*\otimes_S V^*$ with $(V\otimes_S V)^*$ via $\theta_r$, and ${}^*V\otimes_S{}^*V$ with ${}^*(V\otimes_S V)$ via $\theta_\ell$.
Thus $R^\perp$ and ${}^\perp R$ in tensor algebras mean the inverse images of the orthogonal complements in Definition~\ref{def:orthogonals} under these evaluations.
\end{convention}

\begin{definition}\label{def:quadratic-duals}
The \defn{right quadratic dual} and \defn{left quadratic dual} of $A$ are
\begin{equation}\label{eq:quadratic-duals}
\begin{aligned}
A^!&:=T_S(V^*)/(R^\perp),
&R^\perp&=\{F\in V^*\otimes_S V^*\mid\theta_r(F)(R)=0\},\\
{}^!A&:=T_S({}^*V)/({}^\perp R),
&{}^\perp R&=\{F\in{}^*V\otimes_S{}^*V\mid\theta_\ell(F)(R)=0\}.
\end{aligned}
\end{equation}
\end{definition}

We use $A^!$ unless a side is specified.
The ordinary $S$-duals $V^*$ and ${}^*V$ are placed in degree $1$ in \eqref{eq:quadratic-duals}.
By contrast, when $V$ has degree $1$, its graded $S$-dual $V^{\grstar}$ has degree $-1$ by Definition~\ref{def:graded-S-duals}.

By Proposition~\ref{prop:quadratic-basic}, a quadratic presentation identifies $S$ with $A_0$, $V$ with $A_1$, and $R$ with $\Ker\mu_2$, where $\mu_2:A_1\otimes_{A_0}A_1\to A_2$ is \eqref{eq:multiplication-maps}.
The natural tensor-dual isomorphisms \eqref{eq:tensor-duals} therefore make the dual algebras independent of the chosen presentation up to isomorphism.
An isomorphism \defn{over $S$} means one restricting to $\id_S$ on degree $0$.

\begin{advanced}
The constructions are those of \cite[\S\S2.7--2.8]{BGS96}, with different star conventions: their $V^*$ uses left $S$-linear maps and their ${}^*V$ uses right $S$-linear maps.
Thus our $A^!$ is their ${}^!A$, and our ${}^!A$ is their $A^!$.
\end{advanced}

\begin{proposition}\label{prop:dual-low-degrees}
The algebra $A^!$ is quadratic and
\[
A^!_0=S,\qquad A^!_1=V^*,\qquad A^!_2\cong R^*.
\]
In particular, $\dim_\Bbbk A^!_2=\dim_\Bbbk R$.
\end{proposition}

\begin{proof}
Restriction to $R$ gives an exact sequence of $S$-bimodules
\begin{equation}\label{eq:restriction-relations}
0\longrightarrow R^\perp\longrightarrow V^*\otimes_S V^*
\xrightarrow{F\mapsto\theta_r(F)|_R}R^*\longrightarrow0.
\end{equation}
Surjectivity follows from Theorem~\ref{thm:artin-wedderburn}: $R$ has a complement as a right $S$-module, so a right $S$-linear map $R\to S$ extends by zero on that complement.
Taking the quotient in \eqref{eq:restriction-relations} gives $A^!_2\cong R^*$; Proposition~\ref{prop:quadratic-basic} gives the remaining assertions about $A^!$.

\begin{advanced}[nobreak=true]
For the dimension assertion, Theorem~\ref{thm:artin-wedderburn} reduces to simple right modules over matrix rings.
For a finite-dimensional division algebra $\Delta$, the right $\Mat_n(\Delta)$-module $\Delta^{1\times n}$ has dual $\Delta^{n\times1}$: a column $c$ gives the homomorphism $x\mapsto cx$.
The row and column spaces have the same $\Bbbk$-dimension, and dimensions add over simple summands.
Hence $\dim_\Bbbk R^*=\dim_\Bbbk R$ without choosing a symmetrising form.\qedhere
\end{advanced}
\end{proof}

Adding relations has the opposite effect after dualising. For $S$-sub-bimodules $R\subset R'\subset V\otimes_S V$,
\[
R\subset R'\quad\Longrightarrow\quad(R')^\perp\subset R^\perp
\quad\Longrightarrow\quad
\bigl(T_S(V)/(R')\bigr)^!\onto\bigl(T_S(V)/(R)\bigr)^!.
\]
The displayed surjection fixes $S$ and is induced by $\id_{V^*}$ on the generators.

\subsection{Recovering the original algebra}

% ChangeProposal9 consistency update: forward bidual maps now use ev
% (evaluation), as in lem:S-bidual. The left/right dual conventions and
% mathematical arguments are unchanged.

\begin{lemma}\label{lem:double-orthogonal}
Let $U\subset M$ be an $S$-sub-bimodule of a finite-dimensional $S$-bimodule. Under the evaluation isomorphisms of Lemma~\ref{lem:S-bidual},
\[
{}^\perp(U^\perp)=U,\qquad ({}^\perp U)^\perp=U.
\]
\end{lemma}

\begin{proof}
For $u\in U$ and $f\in U^\perp$, evaluation gives $\mathrm{ev}_M(u)(f)=f(u)=0$, so $U\subset{}^\perp(U^\perp)$ under Lemma~\ref{lem:S-bidual}.
If $m\notin U$, Corollary~\ref{cor:semisimple-summands} embeds $M/U$ in $S^{\oplus n}$ as a right module for some $n$.
A coordinate of the image of $m+U$ is non-zero.
Composing $M\onto M/U\hookrightarrow S^{\oplus n}$ with that coordinate projection gives a right $S$-linear map $f:M\to S$ satisfying $f(U)=0$ and $f(m)\neq0$.
Thus $f\in U^\perp$ and $m\notin{}^\perp(U^\perp)$.
The proof of $({}^\perp U)^\perp=U$ uses left modules instead.
\end{proof}

\begin{theorem}\label{thm:quadratic-bidual}
Left and right quadratic duality are canonically inverse:
\[
{}^!(A^!)\cong A\cong({}^!A)^!
\quad\text{as graded algebras over }S.
\]
\end{theorem}

\begin{proof}
\emph{Generators.} The evaluation $\mathrm{ev}_V:V\xrightarrow{\sim}{}^*(V^*)$, $v\mapsto(f\mapsto f(v))$, of Lemma~\ref{lem:S-bidual} identifies the generators of $A$ and ${}^!(A^!)$.
It preserves both $S$-actions by \eqref{eq:tensor-dual-actions}.

\emph{Relations.} For $v,w\in V$ and $f,g\in V^*$, the same actions and \eqref{eq:quadratic-pairings} give
\begin{equation}\label{eq:bidual-tensor-square}
\begin{aligned}
\theta_\ell(\mathrm{ev}_V(v)\otimes\mathrm{ev}_V(w))(g\otimes f)
&=\mathrm{ev}_V(w)(g f(v))=g(f(v)w)\\
&=\mathrm{ev}_{V\otimes_S V}(v\otimes w)(\theta_r(g\otimes f)).
\end{aligned}
\end{equation}
\begin{sloppypar}
Thus $\mathrm{ev}_V\otimes\mathrm{ev}_V$ agrees with evaluation on $V\otimes_S V$ under the tensor-dual identifications \eqref{eq:tensor-duals}.
Lemma~\ref{lem:double-orthogonal}, applied to $R\subset V\otimes_S V$, identifies the relations ${}^\perp(R^\perp)$ of ${}^!(A^!)$ with $R$.
The presentations \eqref{eq:quadratic-duals} therefore give ${}^!(A^!)\cong A$ over $S$.
Interchanging left and right gives $({}^!A)^!\cong A$.
\end{sloppypar}
\end{proof}

\begin{proposition}\label{prop:quadratic-dual-sides}
A choice of symmetrising form $\tau:S\to\Bbbk$ gives isomorphisms $A^!\cong{}^!A$ and $(A^!)^!\cong A$ of graded algebras over $S$.
\end{proposition}

\begin{advanced}
\begin{proof}
The maps $\alpha_V(f)=\tau\circ f$ and $\beta_V(g)=\tau\circ g$ in \eqref{eq:dual-comparison} give a bimodule isomorphism $\beta_V^{-1}\alpha_V:V^*\to{}^*V$.
For $f,g\in V^*$, write $f'=\beta_V^{-1}\alpha_V(f)$ and $g'=\beta_V^{-1}\alpha_V(g)$. Then
\[
\begin{aligned}
\tau\bigl(f'(vg'(w))\bigr)
&=\tau\bigl(f(v)g'(w)\bigr)
=\tau\bigl(g'(f(v)w)\bigr)\\
&=\tau\bigl(g(f(v)w)\bigr)
\qquad(v,w\in V).
\end{aligned}
\]
By \eqref{eq:orthogonal-comparison}, $(\beta_V^{-1}\alpha_V)^{\otimes_S2}$ therefore identifies $R^\perp$ with ${}^\perp R$.
The bimodule map $\beta_V^{-1}\circ\alpha_V$ extends by Proposition~\ref{prop:tensor-universal} to an isomorphism $T_S(V^*)\to T_S({}^*V)$ and hence descends to $A^!\cong{}^!A$.
Apply the same comparison to $A^!$ and then Theorem~\ref{thm:quadratic-bidual} to obtain $(A^!)^!\cong{}^!(A^!)\cong A$.
\end{proof}
\end{advanced}

Theorem~\ref{thm:quadratic-bidual} needs no choice of form. Proposition~\ref{prop:quadratic-dual-sides} permits the usual statement that quadratic duality is an involution, but the identification of left and right duals uses $\tau$.

\subsection{Tensor algebras and monomial relations}

\begin{example}\label{ex:tensor-square-zero-dual}
Quadratic duality exchanges having no relations with setting every product of two generators equal to zero.
Taking $R=0$ or $R=V\otimes_S V$ in \eqref{eq:quadratic-duals} gives
\[
\begin{aligned}
\bigl(T_S(V)\bigr)^!
&=T_S(V^*)/(V^*\otimes_S V^*)\cong S\oplus V^*,\\
\bigl(T_S(V)/(V\otimes_S V)\bigr)^!&=T_S(V^*).
\end{aligned}
\]
In $S\oplus V^*$, $(V^*)^2=0$ and $(s,f)(t,g)=(st,sg+ft)$.
\end{example}

For the rest of this subsection, take $S=\bigoplus_{i=1}^r\Bbbk e_i$, with $e_ie_j=\delta_{ij}e_i$ and $\sum_i e_i=1$.
Choose a basis $\mathcal B_{ij}$ of each $e_iVe_j$ and let $\mathcal B=\bigcup_{i,j}\mathcal B_{ij}$.
A non-empty word $b_1\cdots b_d$ is \defn{composable} if $b_t\in\mathcal B_{i_{t-1},i_t}$ for $1\leq t\leq d$ and some $i_0,\ldots,i_d\in\{1,\ldots,r\}$.
Thus adjacent factors have matching local units.
By Lemma~\ref{lem:idempotent-tensors}, the composable words form a basis of the positive degree part of $T_S(V)$.

\begin{definition}
The presentation $A=T_S(V)/(R)$ is \defn{monomial quadratic} with respect to $\mathcal B$ if $R$ is spanned by a subset $\mathcal R$ of the composable words of length $2$ in $\mathcal B$.
\end{definition}

For $b\in\mathcal B_{ij}$, define the right $S$-linear map $b^*:V\to S$ by its values on $\mathcal B$:
\begin{equation}\label{eq:dual-basis-idempotents}
b^*(c)=\delta_{bc}e_j\quad(c\in\mathcal B),
\qquad e_jb^*e_i=b^*.
\end{equation}
The elements $b^*$ form a $\Bbbk$-basis of $V^*$; the left and right local units of each basis element are reversed.

\begin{proposition}\label{prop:monomial-dual}
For a monomial quadratic presentation:
\begin{enumerate}
\item A basis of $A$ consists of the $e_i$ and the non-empty composable words containing no adjacent pair from $\mathcal R$.
\item The quadratic dual is monomial, with
\begin{equation}\label{eq:monomial-dual}
R^\perp=\operatorname{span}_\Bbbk
\{c^*\otimes_S b^*\mid bc\text{ is composable and }bc\notin\mathcal R\}.
\end{equation}
\end{enumerate}
\end{proposition}

\begin{proof}
The ideal generated by $\mathcal R$ is spanned by the composable words containing an element of $\mathcal R$ as an adjacent pair. Removing those basis words proves (1).
For $b\in\mathcal B_{ij}$ and $c\in\mathcal B_{j\ell}$, \eqref{eq:quadratic-pairings} and \eqref{eq:dual-basis-idempotents} give
\[
\theta_r(c^*\otimes b^*)(b\otimes c)=c^*(e_jc)=e_\ell.
\]
The functional vanishes on every other basis tensor.
By Lemma~\ref{lem:idempotent-tensors}, the composable $c^*\otimes b^*$ form a basis of $V^*\otimes_S V^*$.
Evaluating a linear combination on each $b\otimes c\in\mathcal R$ forces precisely the coefficient of $c^*\otimes b^*$ to be zero, since $e_\ell\neq0$.
Thus the annihilator is the span of the complementary basis tensors, proving (2).
\end{proof}

\begin{example}\label{ex:An-quadratic-dual}
We apply Example~\ref{ex:tensor-square-zero-dual} to the tensor algebra description of $\Bbbk A_n$ in Example~\ref{ex:tensor-An}.
Take $A=\Bbbk A_n$, $S=\bigoplus_{i=1}^n\Bbbk e_{ii}$, and $V=\bigoplus_{i=1}^{n-1}\Bbbk e_{i,i+1}$.
There are no quadratic relations, so $A^!$ has the reversed generators $e_{i,i+1}^*\in e_{i+1,i+1}V^*e_{ii}$ and every composable product of two generators is zero.
With $C$ as in Example~\ref{ex:hilbert-triangular}, this gives
\[
h_{A^!}(t)=I_n+tC^{\mathsf T}.
\]
Here ${}^{\mathsf T}$ denotes transpose: $(C^{\mathsf T})_{ij}=C_{ji}$.
For $n=3$,
\[
(\Bbbk A_3)^!\cong\bigl(\Bbbk A_3/(e_{13})\bigr)^\op,
\qquad
\bigl(\Bbbk A_3/(e_{13})\bigr)^!\cong(\Bbbk A_3)^\op.
\]
These are graded isomorphisms fixing the diagonal idempotents.
\end{example}

\begin{example}\label{ex:hilbert-dual-An}
For $A=\Bbbk A_2$, Examples~\ref{ex:hilbert-triangular} and~\ref{ex:An-quadratic-dual} give
\[
h_A(t)=\begin{pmatrix}1&t\\0&1\end{pmatrix},\qquad
h_{A^!}(-t)^{\mathsf T}=\begin{pmatrix}1&-t\\0&1\end{pmatrix},\qquad
h_A(t)h_{A^!}(-t)^{\mathsf T}=I_2.
\]
The transpose aligns the indices because an $S$-dual reverses the two idempotents in \eqref{eq:dual-basis-idempotents}.
\end{example}



\begin{example}\label{ex:alternating-quadratic-dual}
Proposition~\ref{prop:monomial-dual} exchanges the forbidden and surviving composable pairs.
Take $A=T_S(V)/(ab)$ from Example~\ref{ex:two-vertex-relations}, where $S=\Bbbk e_1\oplus\Bbbk e_2$, $e_1ae_2=a$, and $e_2be_1=b$.
The only composable words of length $2$ are $ab$ and $ba$. Thus
\[
A^!=T_S(V^*)/(a^*b^*),\qquad
A^!=S\oplus\Bbbk a^*\oplus\Bbbk b^*\oplus\Bbbk b^*a^*.
\]
The original relation $ab=0$ corresponds to the surviving word $b^*a^*$; the surviving word $ba$ corresponds to the dual relation $a^*b^*=0$.
In the order $e_1,e_2$, their Hilbert series are
\[
h_A(t)=\begin{pmatrix}1&t\\t&1+t^2\end{pmatrix},\qquad
h_{A^!}(t)=\begin{pmatrix}1+t^2&t\\t&1\end{pmatrix}.
\]
The degree $2$ terms occur in different entries: $ba\in e_2Ae_2$, whereas $b^*a^*\in e_1A^!e_1$.
\end{example}

\subsection{Symmetric and exterior algebras}

\begin{proposition}\label{prop:symmetric-exterior-dual}
For finite-dimensional $\Bbbk$-vector spaces $V,W$,
\[
\SymAlg(V)^!\cong\bigwedge D(V),\qquad
\bigl(\bigwedge W\bigr)^!\cong\SymAlg(D(W)).
\]
\end{proposition}

\begin{proof}
For $\SymAlg(V)$, choose a basis $x_1,\ldots,x_n$ of $V$, with dual basis $x_1^*,\ldots,x_n^*$ of $D(V)$.
Here $S=\Bbbk$, and the evaluation \eqref{eq:quadratic-pairings} reads
\[
\theta_r(x_j^*\otimes x_i^*)(x_p\otimes x_q)=\delta_{ip}\delta_{jq}.
\]
A tensor $\sum_{i,j=1}^n c_{ij}x_i^*\otimes x_j^*$ vanishes on each $x_p\otimes x_q-x_q\otimes x_p$ precisely when $c_{pq}=c_{qp}$ for $1\leq p<q\leq n$.
Thus the orthogonal complement of the symmetric relations is spanned by $(x_i^*)^2$ for $1\leq i\leq n$ and $x_i^*x_j^*+x_j^*x_i^*$ for $1\leq i<j\leq n$, which are the exterior relations on $D(V)$.

For $\bigwedge W$, choose a basis $y_1,\ldots,y_m$ of $W$, with dual basis $y_1^*,\ldots,y_m^*$ of $D(W)$.
A tensor $\sum_{i,j=1}^m c_{ij}y_i^*\otimes y_j^*$ annihilates the exterior relations precisely when $c_{ii}=0$ for $1\leq i\leq m$ and $c_{pq}+c_{qp}=0$ for $1\leq p<q\leq m$.
The orthogonal complement is therefore spanned by $y_i^*y_j^*-y_j^*y_i^*$ for $1\leq i<j\leq m$, which are the symmetric relations on $D(W)$.
The presentations in \eqref{eq:polynomial-exterior} give both isomorphisms.
\end{proof}

Consequently, to pair $\SymAlg(V)$ with $\bigwedge W$ by quadratic duality, take $W=D(V)$ and $y_i=x_i^*$, so $m=n$.
Evaluation $D(D(V))\cong V$ from Example~\ref{ex:linear-duality} then identifies $(\bigwedge W)^!$ with $\SymAlg(V)$.

The proof works in characteristic $2$: the relations $(x_i^*)^2=0$ remain part of the exterior presentation even when the signs $+$ and $-$ coincide.
Example~\ref{ex:symmetric-orthogonals} is the calculation for two generators.
Quadratic duality need not preserve Hilbert series: if $\dim_\Bbbk V>0$, Example~\ref{ex:hilbert-polynomial-exterior} gives different series for $\SymAlg(V)$ and $\SymAlg(V)^!\cong\bigwedge D(V)$.

\begin{advanced}
The \emph{Koszul condition}, to be defined in Lecture~5, implies the identity
\begin{equation}\label{eq:matrix-koszul-preview}
h_A(t)h_{A^!}(-t)^{\mathsf T}=I_r
\qquad\text{when }A_0=\bigoplus_{i=1}^r\Bbbk e_i.
\end{equation}
This is \cite[Lemma~2.11.1]{BGS96}; the proof will be given in Lecture~8.
Equation~\eqref{eq:matrix-koszul-preview} is not asserted for arbitrary quadratic algebras.
For $\SymAlg(V)$ and its quadratic dual $\bigwedge D(V)$, with $n=\dim_\Bbbk V$, the case $S=\Bbbk$ gives $(1-t)^{-n}(1-t)^n=1$.
\end{advanced}

\begin{advanced}[suppressfirstparskip=true,innertopmargin=0pt]
\subsection{Tensor products and signs}

\begin{definition}\label{def:algebra-tensor-products}
For graded $\Bbbk$-algebras $A,B$, the \defn{ordinary tensor product} $A\otimes_\Bbbk B$ and the \defn{graded tensor product} $A\otimes_\Bbbk^{\mathrm{gr}}B$ have underlying vector space $A\otimes_\Bbbk B$ and the total grading of Definition~\ref{def:total-tensor-grading}, with products
\begin{equation}\label{eq:signed-tensor-product}
\begin{aligned}
(a\otimes b)(a'\otimes b')&=aa'\otimes bb'
&&\text{in }A\otimes_\Bbbk B,\\
(a\otimes b)(a'\otimes b')&=(-1)^{\deg b\,\deg a'}aa'\otimes bb'
&&\text{in }A\otimes_\Bbbk^{\mathrm{gr}}B
\end{aligned}
\end{equation}
for homogeneous $a,a'\in A$ and $b,b'\in B$, extended bilinearly.
\end{definition}

Both products have identity $1_A\otimes1_B$.
For the graded product, associativity follows from
$\deg b\,\deg a'+(\deg b+\deg b')\deg a''=\deg b'\,\deg a''+\deg b(\deg a'+\deg a'')$
for homogeneous $a',a''\in A$ and $b,b'\in B$.
For homogeneous $a\in A$ and $b\in B$, the elements $a\otimes1$ and $1\otimes b$ commute in the ordinary tensor product; in the graded tensor product they satisfy
\[
(1\otimes b)(a\otimes1)=(-1)^{\deg a\,\deg b}(a\otimes1)(1\otimes b).
\]
Suppose $A_0=S$ and $B_0=T$ are finite-dimensional semisimple and $S\otimes_\Bbbk T$ is semisimple.
The sign does not change the underlying graded $(S\otimes_\Bbbk T)$-bimodule, so the two tensor products have the same Hilbert series whenever $A,B$ are locally finite-dimensional.
For $S=\Bbbk^{\oplus r}$ and $T=\Bbbk^{\oplus s}$, Lemma~\ref{lem:hilbert-field-tensor} gives its entries:
\[
h_{A\otimes_\Bbbk B}(t)_{(i,p),(j,q)}
=h_A(t)_{ij}h_B(t)_{pq}.
\]
In particular, when $S=T=\Bbbk$, this is $h_A(t)h_B(t)$.

\begin{proposition}\label{prop:dual-tensor-products}
Let $A=T_S(V)/(R)$ and $B=T_T(W)/(Q)$ be quadratic algebras, and assume that $S\otimes_\Bbbk T$ is semisimple.
Then both tensor products are quadratic over $S\otimes_\Bbbk T$, and
\begin{equation}\label{eq:dual-tensor-products}
(A\otimes_\Bbbk B)^!\cong A^!\otimes_\Bbbk^{\mathrm{gr}}B^!,\qquad
(A\otimes_\Bbbk^{\mathrm{gr}}B)^!\cong A^!\otimes_\Bbbk B^!.
\end{equation}
\end{proposition}

The mixed commuting relations become anticommuting relations under the reversed tensor pairing.

The assumption on $S\otimes_\Bbbk T$ holds when $S,T$ are products of matrix algebras over $\Bbbk$, since $\Mat_a(\Bbbk)\otimes_\Bbbk\Mat_b(\Bbbk)\cong\Mat_{ab}(\Bbbk)$.
It does not follow from semisimplicity alone: take $S=T$ to be the field in Example~\ref{ex:inseparable-envelope}.

\begin{proof}
The degree $1$ part of $A\otimes_\Bbbk B$ is $(V\otimes_\Bbbk T)\oplus(S\otimes_\Bbbk W)$.
By Lemma~\ref{lem:tensor-properties}, its tensor square over $S\otimes_\Bbbk T$ decomposes into the four summands
\[
(V\otimes_S V)\otimes_\Bbbk T
\ \oplus\ \underbrace{V\otimes_\Bbbk W}_{\text{order }VW}
\ \oplus\ \underbrace{V\otimes_\Bbbk W}_{\text{order }WV}
\ \oplus\ S\otimes_\Bbbk(W\otimes_T W).
\]
Its quadratic relations consist of $R\otimes_\Bbbk T$, $S\otimes_\Bbbk Q$, and the mixed relations
\begin{equation}\label{eq:mixed-tensor-relations}
(v\otimes1)(1\otimes w)-(1\otimes w)(v\otimes1)
\quad(v\in V,\ w\in W).
\end{equation}
Here the products are taken in the tensor algebra on the degree $1$ part, and their span is an $(S\otimes_\Bbbk T)$-sub-bimodule.
Let $C$ be the quotient of that tensor algebra by the three specified relation spaces.
The generator maps induce algebra homomorphisms $A\to C$, $a\mapsto\overline a$, and $B\to C$, $b\mapsto\overline b$.
Their images commute: the degree $0$ actions already commute across the two factors, and \eqref{eq:mixed-tensor-relations} makes the degree $1$ generators commute.
Consequently $A\otimes_\Bbbk B\to C$, $a\otimes b\mapsto\overline a\,\overline b$, is an algebra homomorphism.
It is inverse to the canonical homomorphism $C\to A\otimes_\Bbbk B$, since both composites fix the generating copies of $S,T,V,W$. Thus the stated relations give exactly $A\otimes_\Bbbk B$.
For the graded tensor product, the same construction uses $\overline b\,\overline a=(-1)^{\deg a\deg b}\overline a\,\overline b$ for homogeneous $a,b$.

The right $(S\otimes_\Bbbk T)$-duals of the two generator summands are $V^*\otimes_\Bbbk T$ and $S\otimes_\Bbbk W^*$, with evaluations
\[
(f\otimes t)(v\otimes t')=f(v)\otimes tt',\qquad
(s\otimes g)(s'\otimes w)=ss'\otimes g(w).
\]
These identifications hold for finite free right modules and hence for their direct summands by Corollary~\ref{cor:semisimple-summands}.
The reversed tensor evaluation pairs each pure summand with its corresponding dual summand and exchanges the two mixed summands; all other pairings are zero.
The pure relation spaces have orthogonal complements $R^\perp\otimes_\Bbbk T$ and $S\otimes_\Bbbk Q^\perp$ in the corresponding pure summands.
The mixed relation space is the set of pairs $(z,-z)$ in the two copies of $V\otimes_\Bbbk W$.
For mixed tensors, the two reversed evaluations of $f\in V^*$ and $g\in W^*$ both give $f(v)\otimes g(w)$ on the corresponding copy.
Thus an annihilating pair of functionals must have equal components. As these simple tensors span the full dual, the mixed annihilator is spanned by
\[
(f\otimes1)(1\otimes g)+(1\otimes g)(f\otimes1)
\quad(f\in V^*,\ g\in W^*).
\]
These are the mixed relations of $A^!\otimes_\Bbbk^{\mathrm{gr}}B^!$, proving the first isomorphism in \eqref{eq:dual-tensor-products}.
For $A\otimes_\Bbbk^{\mathrm{gr}}B$, the minus sign in \eqref{eq:mixed-tensor-relations} is a plus sign. The mixed annihilator then has a minus sign, giving the ordinary tensor product on the dual side.
\end{proof}

\begin{example}\label{ex:commuting-square-dual}
We compute how Proposition~\ref{prop:dual-tensor-products} changes a commutation relation in a matrix example.
Take $A=\Bbbk A_2\otimes_\Bbbk\Bbbk A_2$.
Its degree $0$ algebra has basis $e_{ii}\otimes e_{jj}$ for $i,j\in\{1,2\}$, and its degree $1$ part has basis
\[
a=e_{12}\otimes e_{11},\quad
b=e_{22}\otimes e_{12},\quad
c=e_{11}\otimes e_{12},\quad
d=e_{12}\otimes e_{22}.
\]
By Lemma~\ref{lem:idempotent-tensors}, the tensor square of the degree $1$ part has basis $ab,cd$, and both map to $e_{12}\otimes e_{12}$. Thus the relation space is $\Bbbk(ab-cd)$.
Using the reversed generators from \eqref{eq:dual-basis-idempotents}, its orthogonal complement is $\Bbbk(b^*a^*+d^*c^*)$.
Hence the dual presentation has an anticommuting relation:
\[
ab=cd\text{ in }A,\qquad
b^*a^*=-d^*c^*\text{ in }A^!.
\]
Order the degree $0$ idempotents as
$e_{11}\otimes e_{11},e_{22}\otimes e_{11},e_{11}\otimes e_{22},e_{22}\otimes e_{22}$.
The homogeneous bases give
\[
h_A(t)=\begin{pmatrix}
1&t&t&t^2\\
0&1&0&t\\
0&0&1&t\\
0&0&0&1
\end{pmatrix},\qquad
h_{A^!}(t)=h_A(t)^{\mathsf T}.
\]
The signs agree with Proposition~\ref{prop:dual-tensor-products}; in characteristic $2$, the same presentations remain valid with $-1=1$.
Here replacing $d^*$ by $-d^*$ removes the minus sign. In this example the sign distinguishes the specified presentations, not the isomorphism classes of algebras.
\end{example}

\begin{example}\label{ex:essential-tensor-sign}
The symmetric/exterior pair shows that the sign in a graded tensor product cannot always be omitted.
Let $V=\Bbbk x_1\oplus\Bbbk x_2$ and let $y_1,y_2$ be the dual basis of $W=D(V)$. Proposition~\ref{prop:dual-tensor-products} gives
\[
\SymAlg(V)^!
\cong\bigl(\Bbbk[x_1]\otimes_\Bbbk\Bbbk[x_2]\bigr)^!
\cong\frac{\Bbbk[y_1]}{(y_1^2)}\otimes_\Bbbk^{\mathrm{gr}}
       \frac{\Bbbk[y_2]}{(y_2^2)}
\cong\bigwedge W.
\]
If $\Char\Bbbk\neq2$, the ordinary tensor product of the two square-zero algebras is commutative, whereas the graded tensor product is not: $y_1y_2=-y_2y_1\neq0$.
They are therefore not isomorphic. This sign cannot be removed by a change of generators.
\end{example}
\end{advanced}

\subsection{Exercises}

\begin{exercise}
For $A=\Bbbk A_4/(e_{13})$, write a presentation of $A^!$ using reversed generators and find homogeneous bases of $A$ and $A^!$.
Check the double dual using Proposition~\ref{prop:monomial-dual}.
\emph{Hint.} Among $e_{12}e_{23}$ and $e_{23}e_{34}$, only the first is a relation. Complement the forbidden pairs and reverse their order.
\end{exercise}

\begin{exercise}
Let $S=\Bbbk e_1\oplus\Bbbk e_2$ and let $V$ have basis $a,b$ with $e_1ae_2=a$ and $e_2be_1=b$.
Compute the quadratic dual of $T_S(V)/(ab,ba)$ and the Hilbert series of both algebras.
\emph{Check.} In the order $e_1,e_2$, the two series are
\[
\begin{pmatrix}1&t\\t&1\end{pmatrix}
\qquad\text{and}\qquad
\frac{1}{1-t^2}\begin{pmatrix}1&t\\t&1\end{pmatrix}.
\]
\end{exercise}

\begin{advanced}
\begin{exercise}
For the four generators in Example~\ref{ex:commuting-square-dual}, replace the relation $ab-cd$ by $ab-\lambda cd$, where $\lambda\in\Bbbk$.
Compute the quadratic dual relation for every $\lambda$, including $\lambda=0$, and recover the original relation by dualising again.

\emph{Hint.} Evaluate $\mu b^*a^*+\nu d^*c^*$ on $ab-\lambda cd$. For $\lambda\neq0$, rescale $d$ to compare with Example~\ref{ex:commuting-square-dual}.
\end{exercise}
\end{advanced}


% !TEX root = Lecture_Koszul.tex
\clearpage
\section{Minimal graded resolutions}\label{sec:minimal-resolutions}

% TODO (ChangeProposal9, Lecture 4 follow-up): After approval, refer to
% def:positive-ideal and eq:degree-zero-projection for A_+ and A/A_+ = S.
% When Ext and Tor are introduced in Lectures 6--7, explain
% Hom_A(M,S) = Hom_S(M/MA_+,S) naturally, and
% Ext_A^i(M,S) = Hom_S(Tor_i^A(M,S),S) naturally (ungraded right modules).
% For a minimal graded resolution, identify the graded Ext groups with the
% S-duals of the generator modules P_i/P_iA_+; check the shift sign carefully.
% This is not higher Hom_S(-,S), which is exact. Specify the Yoneda algebra
% action and its opposite convention when formulating Koszul duality.

\subsection{Projective covers and the graded Nakayama lemma}

Let $A=\bigoplus_{d\geq0}A_d$ be a locally finite-dimensional graded algebra with $S=A_0$ finite-dimensional semisimple.
Write
\[
A_+:=\bigoplus_{d>0}A_d.
\]
The grading makes $A_+$ a two-sided ideal, and projection onto degree $0$ identifies $A/A_+$ with $S$.
We regard $S$ as a graded right $A$-module concentrated in degree $0$, with action $s\cdot a=sa_0$ for $s\in S$ and $a=\sum_{d\geq0}a_d\in A$.

\phantomsection\label{not:chain-differentials}
All complexes below are chain complexes of graded right $A$-modules: their terms $P_i$ are indexed by homological degrees $i\geq0$, and their differentials are degree preserving maps $d_i:P_i\to P_{i-1}$ satisfying $d_{i-1}d_i=0$ for $i\geq2$.
An augmented complex to a module $M$ also has a degree preserving map $p_0:P_0\to M$ with $p_0d_1=0$.
The augmented complex is exact when $\Im d_{i+1}=\Ker d_i$ for $i\geq1$, $\Im d_1=\Ker p_0$, and $p_0$ is surjective.

\phantomsection\label{not:graded-top}
For a graded right $A$-module $M$, define its \defn{graded top} by
\[
\overline M:=M/MA_+.
\]
The quotient $\overline M$ is a graded right $S$-module.
We call $M$ \defn{bounded below} if $M_d=0$ for all sufficiently small $d$.
Every finitely generated graded module is bounded below: a finite generating set can be replaced by its finitely many homogeneous components, and $A$ has no negative degrees.

\begin{definition}\label{def:graded-projective}
A graded right $A$-module $P$ is \defn{projective} if, for every surjective degree preserving map $q:X\to Y$ and every degree preserving map $f:P\to Y$, there is a degree preserving map $\widetilde f:P\to X$ such that $q\widetilde f=f$.
\end{definition}

The shifted free module $A(-d)$ has its generator $1$ in degree $d$.
If $Q$ is a finite-dimensional projective right $S$-module, then $Q\otimes_S A(-d)$ is a projective graded right $A$-module: $Q$ is a direct summand of a finite free right $S$-module, so $Q\otimes_S A(-d)$ is a direct summand of a finite direct sum of copies of $A(-d)$.

\begin{lemma}[Graded Nakayama lemma]\label{lem:graded-nakayama}
If a bounded below graded right $A$-module $M$ satisfies $M=MA_+$, then $M=0$.
\end{lemma}

\begin{proof}
If $M\neq0$, let $d$ be the least degree with $M_d\neq0$.
Every degree $d$ element of $MA_+$ is a sum of products $m_j a_{d-j}$ with $j<d$ and $d-j>0$.
All such $m_j$ vanish by the choice of $d$, so $(MA_+)_d=0$, contradicting $M_d=(MA_+)_d$.
\end{proof}

\begin{definition}\label{def:graded-projective-cover}
A \defn{graded projective cover} of a bounded below locally finite-dimensional graded right $A$-module $M$ is a surjective degree preserving map $p:P\to M$ from a bounded below locally finite-dimensional graded projective module such that
\[
\overline p:P/PA_+\xrightarrow{\sim}M/MA_+.
\]
\end{definition}

For a surjection $p:P\to M$, the condition in Definition~\ref{def:graded-projective-cover} is equivalent to $\Ker p\subseteq PA_+$, since $p(PA_+)=MA_+$.

\begin{proposition}\label{prop:graded-projective-covers}
Every bounded below locally finite-dimensional graded right $A$-module has a graded projective cover, unique up to isomorphism over $M$.
\end{proposition}

\begin{proof}
Put $\overline M_d=(M/MA_+)_d$. Since $S$ is semisimple, the quotient map $M_d\to\overline M_d$ has a right $S$-linear section $s_d:\overline M_d\to M_d$.
Set
\[
P:=\bigoplus_{d\in\mathbb Z}(\overline M_d\otimes_S A)(-d).
\]
Each $\overline M_d$ is projective over $S$, so $P$ is a graded projective right $A$-module.
Since $M$ is bounded below and locally finite-dimensional, each degree of $P$ is a finite sum of finite-dimensional spaces; hence $P$ is locally finite-dimensional.
The maps
\[
\overline M_d\otimes_S A\longrightarrow M,\qquad
\overline m\otimes a\longmapsto s_d(\overline m)a
\]
define a degree preserving map $p:P\to M$.
On graded tops, $\overline p$ identifies each copy of $\overline M_d$ with $(M/MA_+)_d$.
Thus $\overline p$ is an isomorphism.
The quotient $C:=M/p(P)$ satisfies $C/CA_+=0$, hence $C=CA_+$; Lemma~\ref{lem:graded-nakayama} gives $C=0$.
Therefore $p$ is a projective cover.

For uniqueness, let $p:P\to M$ and $q:Q\to M$ be projective covers.
Projectivity gives a map $f:P\to Q$ with $qf=p$.
The maps $\overline p$ and $\overline q$ are isomorphisms, so $\overline f=\overline q^{-1}\overline p$ is an isomorphism.
The quotient $Q/f(P)$ has zero graded top and is bounded below, hence is zero by Lemma~\ref{lem:graded-nakayama}.
Thus $f$ is surjective; since $Q$ is projective, $f$ splits.
Its kernel is a direct summand of $P$, and its graded top is zero because $\overline f$ is injective.
Lemma~\ref{lem:graded-nakayama} gives $\Ker f=0$, so $f$ is an isomorphism over $M$.
\end{proof}

\subsection{Minimal resolutions and graded Betti numbers}

\begin{definition}\label{def:minimal-resolution}
A \defn{graded projective resolution} of $M$ is an exact sequence of degree preserving maps
\[
\cdots\longrightarrow P_2\xrightarrow{d_2}P_1\xrightarrow{d_1}P_0\xrightarrow{p_0}M\longrightarrow0
\]
with every $P_i$ graded projective.
It is \defn{minimal} if $p_0:P_0\to M$ is a graded projective cover and, for each $i\geq1$, the map $d_i$ with codomain restricted to $\Ker d_{i-1}$ is a graded projective cover, where $d_0:=p_0$.
\end{definition}

\begin{theorem}\label{thm:minimal-resolution}
Every bounded below locally finite-dimensional graded right $A$-module has a minimal graded projective resolution.
For any two such resolutions $P_\bullet\to M$ and $P'_\bullet\to M$, there are degree preserving isomorphisms $P_i\to P'_i$ commuting with the differentials and the maps to $M$.
\end{theorem}

\begin{proof}
Take a projective cover $p_0:P_0\to M$ using Proposition~\ref{prop:graded-projective-covers}, and set $K_1=\Ker p_0$.
The module $K_1$ is bounded below and locally finite-dimensional, so it has a projective cover $q_1:P_1\to K_1$.
Repeat: after defining $q_i:P_i\to K_i$, set $K_{i+1}=\Ker q_i$ and cover $K_{i+1}$ by $q_{i+1}:P_{i+1}\to K_{i+1}$.
Composing each $q_i$ with the inclusion $K_i\hookrightarrow P_{i-1}$ gives $d_i$ and an exact minimal resolution.

For uniqueness, Proposition~\ref{prop:graded-projective-covers} identifies the covers of $M$ in two minimal resolutions.
The resulting isomorphism identifies their kernels; applying the same proposition to those kernels identifies the next projective covers.
Induction gives an isomorphism in every homological degree, compatible with the differentials and the maps to $M$.
\end{proof}

Choose representatives $L_1,\ldots,L_r$ of the simple right $S$-modules, as in Definition~\ref{def:dimension-vector}.
\phantomsection\label{not:graded-betti}
For a minimal resolution $P_\bullet\to M$, define the \defn{graded Betti numbers} $\beta_{i,d,j}(M)$ by
\[
\left(P_i/P_iA_+\right)_d
\cong\bigoplus_{j=1}^r L_j^{\oplus\beta_{i,d,j}(M)},
\qquad
\beta_{i,d}(M):=(\beta_{i,d,1}(M),\ldots,\beta_{i,d,r}(M)).
\]
Here $i$ is the homological degree, $d$ is the grading degree, and $j$ records the simple right $S$-module.
Theorem~\ref{thm:minimal-resolution} shows these multiplicities do not depend on the chosen minimal resolution.
The projective cover construction also gives
\[
P_i\cong\bigoplus_{d\in\mathbb Z}\bigoplus_{j=1}^r
(L_j\otimes_S A)(-d)^{\oplus\beta_{i,d,j}(M)}.
\]
For $S=\bigoplus_{j=1}^r\Bbbk e_j$, the simple modules are $L_j=e_jS$ and $L_j\otimes_S A\cong e_jA$.

\begin{example}\label{ex:minimal-A2}
This example shows how a non-field degree zero algebra gives distinct projective covers for its different simple modules.
Let $A=\Bbbk A_2$, with $S=\Bbbk e_{11}\oplus\Bbbk e_{22}$ and $A_1=\Bbbk e_{12}$.
The simple right $S$-modules are $L_1=e_{11}S$ and $L_2=e_{22}S$, with projective lifts $Q_1=e_{11}A$ and $Q_2=e_{22}A$.
The quotient map $Q_1\to L_1$ has kernel $\Bbbk e_{12}$, while $Q_2=L_2$.
Sending the degree $1$ generator $e_{22}$ to $e_{12}$ gives the minimal resolution
\[
0\longrightarrow Q_2(-1)\xrightarrow{d_1}Q_1\longrightarrow L_1\longrightarrow0,
\qquad d_1(e_{22})=e_{12}.
\]
\end{example}

\begin{example}\label{ex:minimal-polynomial}
For $A=\Bbbk[x]$, with $\deg x=1$, the degree zero quotient module is $\Bbbk=A/(x)$.
Multiplication by $x$ gives the minimal resolution
\[
0\longrightarrow A(-1)\xrightarrow{d_1}A\xrightarrow{\pi}\Bbbk\longrightarrow0,
\qquad d_1(a)=xa,\quad \pi(f)=f(0).
\]
The map $d_1$ has degree $0$, is injective, and has image $xA=\Ker\pi\subset AA_+$.
Thus $\beta_{0,0}=1$ and $\beta_{1,1}=1$.
\end{example}

\begin{example}\label{ex:minimal-exterior}
The exterior algebra on one generator gives an infinite minimal resolution, although the algebra itself is finite-dimensional.
Let $E=\Bbbk[\epsilon]/(\epsilon^2)$ with $\deg\epsilon=1$ and $E_+=\Bbbk\epsilon$.
The degree zero quotient is $\Bbbk=E/E_+$, and
\[
\cdots\xrightarrow{d_3}E(-2)\xrightarrow{d_2}E(-1)
\xrightarrow{d_1}E\xrightarrow{\pi}\Bbbk\longrightarrow0,
\qquad d_i(a)=\epsilon a,\quad \pi(a)=a_0.
\]
Each $d_i$ is degree preserving and right $E$-linear.
Since $\Ker(\epsilon\cdot -)=\epsilon E=\Im(\epsilon\cdot -)$, the sequence is exact at every $E(-i)$ and at $E$.
Every differential has image in the preceding term times $E_+$, so the resolution is minimal and $\beta_{i,i}=1$ for all $i\geq0$.
\end{example}

\begin{advanced}
If $M$ is finitely generated, its first projective cover is finitely generated.
Without a noetherian hypothesis on $A$, later kernels need not be finitely generated, so later terms of the minimal resolution need not be finitely generated.
The bounded below locally finite-dimensional setting above still gives projective covers and a minimal resolution.
\end{advanced}


% !TEX root = Lecture_Koszul.tex
\clearpage
\section{Koszul complex and Koszul algebras}\label{sec:koszul-complex}

\subsection{The Koszul complex of a polynomial algebra}

Let $V$ be an $n$-dimensional $\Bbbk$-vector space, choose a basis $x_1,\ldots,x_n$, and let $A=\SymAlg(V)=\Bbbk[x_1,\ldots,x_n]$, with each $x_j$ in degree $1$.
Let $\pi:A\to\Bbbk$ be projection onto degree $0$, so $\Bbbk=A/A_+$.
\phantomsection\label{not:koszul-terms}
For $0\leq i\leq n$, put
\[
K_i:=\bigwedge^i V\otimes_\Bbbk A(-i),
\qquad K_0=A.
\]
The right $A$-action is on the second factor.
The factor $\bigwedge^iV$ is placed in grading degree $0$, so the generators of $K_i$ have grading degree $i$.

\begin{definition}\label{def:koszul-complex}
For $1\leq i\leq n$, define the \defn{Koszul differential}
\begin{equation}\label{eq:koszul-differential}
\begin{aligned}
d_i:K_i&\longrightarrow K_{i-1},\\
(v_1\wedge\cdots\wedge v_i)\otimes a
&\longmapsto
\sum_{r=1}^i(-1)^{r-1}
(v_1\wedge\cdots\wedge\widehat{v_r}\wedge\cdots\wedge v_i)
\otimes v_ra .
\end{aligned}
\end{equation}
The formula is multilinear and alternating in $v_1,\ldots,v_i$, so it defines a map on $\bigwedge^iV$.
The sequence $K_\bullet\to\Bbbk$ consists of these maps followed by $\pi:K_0=A\to\Bbbk$.
\end{definition}

\begin{proposition}\label{prop:koszul-complex-resolution}
The sequence
\[
0\longrightarrow K_n\xrightarrow{d_n}\cdots\xrightarrow{d_2}K_1
\xrightarrow{d_1}K_0\xrightarrow{\pi}\Bbbk\longrightarrow0
\]
is an exact minimal graded free resolution of $\Bbbk$ over $A$.
\end{proposition}

\begin{proof}
Each $d_i$ preserves degree and is right $A$-linear.
For $i\geq2$, compute $d_{i-1}d_i$ on a basis wedge.
Every term is indexed by two distinct positions $r<s$ deleted from that wedge.
Deleting position $r$ and then $s$ gives coefficient $(-1)^{r+s-3}x_{j_s}x_{j_r}$; deleting $s$ and then $r$ gives coefficient $(-1)^{r+s-2}x_{j_r}x_{j_s}$.
Since $x_{j_s}x_{j_r}=x_{j_r}x_{j_s}$, these two terms sum to zero. Hence $d_{i-1}d_i=0$; also $\pi d_1=0$ because every term of $d_1$ is divisible by a variable.

To prove exactness, give $A$ and $\bigwedge V$ the multigrading with $\deg x_j$ equal to the $j$th standard basis vector.
For $\alpha\in\mathbb Z_{\geq0}^n$, let $\operatorname{supp}(\alpha)=\{j\mid\alpha_j>0\}$.
The degree $\alpha$ part of $K_i$ has basis
\[
b_J:=(x_{j_1}\wedge\cdots\wedge x_{j_i})\otimes x^{\alpha-\mathbf1_J},
\qquad
J=\{j_1<\cdots<j_i\}\subseteq\operatorname{supp}(\alpha),
\]
where $\mathbf1_J$ has entry $1$ at indices in $J$ and $0$ elsewhere.
The differential is $d_i(b_J)=\sum_{r=1}^i(-1)^{r-1}b_{J\setminus\{j_r\}}$.

If $\alpha\neq0$, let $q$ be the least index in $\operatorname{supp}(\alpha)$.
For each $i\geq0$, define $h_{\alpha,i}:K_{i,\alpha}\to K_{i+1,\alpha}$ on basis vectors by
\[
h_{\alpha,i}(b_J)=
\begin{cases}
b_{\{q\}\cup J},&q\notin J,\\
0,&q\in J,
\end{cases}
\]
inserting $q$ first in the ordered wedge.
For $q\notin J$ and $i\geq1$, direct substitution gives
\[
\begin{aligned}
d_{i+1}h_{\alpha,i}(b_J)
&=b_J-\sum_{r=1}^i(-1)^{r-1}b_{\{q\}\cup(J\setminus\{j_r\})},\\
h_{\alpha,i-1}d_i(b_J)
&=\sum_{r=1}^i(-1)^{r-1}b_{\{q\}\cup(J\setminus\{j_r\})}.
\end{aligned}
\]
For $i=0$, $J=\varnothing$ and $d_1h_{\alpha,0}(b_\varnothing)=b_\varnothing$.
For $q\in J$, the term of $d_i(b_J)$ deleting $q$ is the only term on which $h_{\alpha,i-1}$ is nonzero, and it gives $h_{\alpha,i-1}d_i(b_J)=b_J$ while $d_{i+1}h_{\alpha,i}(b_J)=0$.
Thus $d_{i+1}h_{\alpha,i}+h_{\alpha,i-1}d_i=\id_{K_{i,\alpha}}$ for every $i\geq1$ in positive multidegree; the case $i=0$ was checked separately.
All positive multidegrees are therefore exact.
In multidegree $0$, only $K_0\cong\Bbbk$ is nonzero and $\pi$ is the identity onto $\Bbbk$.
This proves exactness of the displayed sequence.

Finally, each $K_i$ is a finite direct sum of copies of $A(-i)$, and \eqref{eq:koszul-differential} has image in $K_{i-1}A_+$.
Exactness gives $\Ker d_i=\Im d_{i+1}\subseteq K_iA_+$ for $i<n$, and $\Ker d_n=0$.
Also, $\Ker\pi=AA_+$.
Thus each map from $K_i$ onto the preceding kernel has kernel contained in $K_iA_+$, so it is a projective cover by Definition~\ref{def:graded-projective-cover}.
The resolution is therefore minimal.
\end{proof}

\subsection{Koszul algebras}

\begin{definition}\label{def:linear-resolution}
A minimal graded projective resolution $P_\bullet\to M$ is \defn{linear} if $P_i/P_iA_+$ is concentrated in degree $i$ for every $i\geq0$.
This is equivalent to $P_i$ being generated in degree $i$: generation in degree $i$ implies the stated concentration because $A_+$ raises degree.
Conversely, if the top is concentrated in degree $i$, boundedness below and Lemma~\ref{lem:graded-nakayama} show that $P_i$ has no components below degree $i$; the quotient $P_i/(P_i)_iA$ then has zero graded top, so Lemma~\ref{lem:graded-nakayama} makes that quotient zero.
\end{definition}

\begin{definition}\label{def:koszul-algebra}
Let $A=T_S(V)/(R)$ be a quadratic algebra with $S=A_0$ finite-dimensional semisimple.
The algebra $A$ is \defn{Koszul} if the degree zero quotient module $S=A/A_+$ has a linear minimal graded projective resolution as a right $A$-module.
\end{definition}

The generators in homological degree $i$ must then all have grading degree $i$.
Since $S$ can have several simple modules, the definition requires a linear resolution of their direct sum, not just of one chosen simple module.

\begin{proposition}\label{prop:polynomial-koszul}
For every finite-dimensional vector space $V$, the symmetric algebra $\SymAlg(V)$ is Koszul.
\end{proposition}

\begin{proof}
Proposition~\ref{prop:koszul-complex-resolution} gives a minimal resolution with $K_i=\bigwedge^iV\otimes_\Bbbk A(-i)$.
Thus $K_i$ is generated in degree $i$ for every $i$.
\end{proof}

\begin{example}\label{ex:koszul-nonconnected}
The definition also applies when $S$ has several simple modules.
For $A=\Bbbk A_2$, the resolution in Example~\ref{ex:minimal-A2} is linear for $L_1$, and $L_2=e_{22}A$ is projective in degree $0$.
Taking their direct sum gives a linear resolution of $S=L_1\oplus L_2$; hence $\Bbbk A_2$ is Koszul.
\end{example}

\begin{example}
The exterior algebra $E=\Bbbk[\epsilon]/(\epsilon^2)$ is quadratic, and Example~\ref{ex:minimal-exterior} gives a minimal resolution of its degree zero quotient with $P_i=E(-i)$.
Thus $E$ is Koszul, although its minimal resolution is infinite.
\end{example}


% TODO (Deferred global request): Insert a linked table of notation, with a
% page break before it, before this bibliography. This lies outside the present
% Lecture 1-only proposal; use first-definition labels when it is authorised.
\clearpage
\begin{thebibliography}{BGS96}
\bibitem[Bez23]{Bez23}
R.~Bezrukavnikov,
\emph{Noncommutative Algebra}, MIT OpenCourseWare 18.706 (2023),
\href{https://ocw.mit.edu/courses/18-706-noncommutative-algebra-spring-2023/resources/mit18_706_s23_lec02_pdf/}{Lecture~2: Semisimple Modules, Socles, Artinian Rings, Wedderburn's Theorem}.

\bibitem[BGS96]{BGS96}
A.~Beilinson, V.~Ginzburg, and W.~Soergel,
\emph{Koszul duality patterns in representation theory},
J.~Amer.~Math.~Soc.~9 (1996), 473--527.

\bibitem[DK94]{DK94}
Yu.~A.~Drozd and V.~V.~Kirichenko,
\emph{Finite Dimensional Algebras},
Springer-Verlag, Berlin, 1994,
\href{https://doi.org/10.1007/978-3-642-76244-4}{doi:10.1007/978-3-642-76244-4}.

\bibitem[DNN17]{DNN17}
S.~D\u{a}sc\u{a}lescu, C.~N\u{a}st\u{a}sescu, and L.~N\u{a}st\u{a}sescu,
\emph{Graded semisimple algebras are symmetric},
J.~Algebra~491 (2017), 207--218,
\href{https://arxiv.org/abs/1504.04868}{arXiv:1504.04868}.

\bibitem[GI20]{GI20}
J.~Grant and O.~Iyama,
\emph{Higher preprojective algebras, Koszul algebras, and superpotentials},
Compos.~Math.~156 (2020), 2588--2627.
\end{thebibliography}

\end{document}
